%  LaTeX - Sistema de Asignación de Salones ISC
\documentclass[algorithms,article,submit,pdftex]{Definitions/mdpi} 

\firstpage{1} 
\makeatletter 
\setcounter{page}{\@firstpage} 
\makeatother
\pubvolume{1}
\issuenum{1}
\articlenumber{0}
\pubyear{2026}
\copyrightyear{2025}

\Title{Sistema de Asignación de Salones ISC: Optimización Inteligente de Espacios Académicos}

\Author{Anónimo}
\AuthorNames{Anónimo}

\address{Instituto Tecnológico de Ciudad Madero, Tecnológico Nacional de México}

\corres{Correspondencia: contacto@ejemplo.com}

\abstract{Sistema inteligente de optimización para la asignación de 680 clases a 21 salones en el programa de Ingeniería en Sistemas Computacionales del Instituto Tecnológico de Ciudad Madero. El sistema implementa cuatro algoritmos de optimización: heurística del profesor, Greedy con Hill Climbing, Machine Learning (Random Forest) y Algoritmo Genético. Se presenta un mecanismo de pre-asignación que garantiza el 100\% de cumplimiento de las preferencias prioritarias (Prioridad 1) mientras optimiza tres objetivos: minimizar movimientos de profesores entre salones, reducir cambios de piso y disminuir la distancia total recorrida. Los resultados experimentales de 90 ejecuciones demuestran que el enfoque Greedy+Hill Climbing logra el mejor rendimiento general con 314 movimientos (mejora del 12\% sobre la línea base), 206 cambios de piso (reducción del 28\%) y resultados consistentes. La validación estadística mediante ANOVA, pruebas post-hoc de Tukey HSD y análisis de tamaño de efecto de Cohen confirma diferencias significativas entre todos los algoritmos. El sistema está implementado en Python con documentación completa y está disponible como software de código abierto.}

\keyword{asignación de salones; optimización combinatoria; algoritmos híbridos; hill climbing; machine learning; algoritmo genético; programación académica; satisfacción de restricciones}

\begin{document}

\section{Sistema de Asignación de Salones ISC}

\textbf{Optimización Inteligente de Espacios Académicos}

\textbf{Jesús Olvera}

Ingeniería en Sistemas Computacionales

Instituto Tecnológico de Ciudad Madero


\subsection{Descripción General}

Sistema inteligente de optimización para la asignación de salones en el programa de Ingeniería en Sistemas Computacionales.

\textbf{Objetivos:}

\begin{itemize}
\item  Minimizar movimientos de profesores
\begin{itemize}
\item  Reducir cambios de piso
\begin{itemize}
\item  Optimizar distancias recorridas
\begin{itemize}
\item  Garantizar cumplimiento de preferencias prioritarias

\subsection{ Sistema de Prioridades Jerárquico}

\textbf{Tres niveles de prioridad:}

\begin{enumerate}
\item \textbf{PRIORIDAD 1 (Hard Constraint):} Preferencias de Profesores
\begin{itemize}
\item Cumplimiento: \textbf{100% garantizado}
\begin{itemize}
\item Implementación: Pre-asignación forzada
\begin{itemize}
\item Protección: Clases inmutables durante optimización
\begin{enumerate}
\item \textbf{PRIORIDAD 2 (Soft Constraint):} Consistencia de Grupos
\begin{itemize}
\item Mantener grupos en el mismo salón
\begin{enumerate}
\item \textbf{PRIORIDAD 3 (Soft Constraint):} Primer Semestre
\begin{itemize}
\item Asignar grupos 15xx a salones específicos

\subsection{🔧 Algoritmos de Optimización}

\textbf{4 optimizadores diferentes:}


\subsection{Arquitectura del Sistema}

\begin{verbatim}
Sistema-Salones-ISC/
├── configurador_materias.py      # Interfaz gráfica
├── pre_asignar_p1.py             # Pre-asignación P1
├── optimizador_greedy.py         # Greedy + Hill Climbing
├── optimizador_ml.py             # Machine Learning
├── optimizador_genetico.py       # Algoritmo Genético
├── corregir_prioridades.py       # Corrección post-opt
├── ejecutar_todos.py             # Script maestro
├── generar_comparativa_completa.py  # Reportes
├── utils_restricciones.py        # Validación
├── datos_estructurados/          # Datos I/O
└── comparativas/                 # Resultados
\end{verbatim}


\subsection{Flujo de Ejecución}

\begin{verbatim}
graph TD
    A[Horario Inicial] --> B[Pre-asignación P1]
    B --> C[00_Horario_PreAsignado_P1.csv]
    C --> D[Optimizador Greedy]
    C --> E[Optimizador ML]
    C --> F[Optimizador Genético]
    D --> G[Corrección Post-Opt]
    E --> H[Corrección Post-Opt]
    F --> I[Corrección Post-Opt]
    G --> J[Comparativas y Gráficos]
    H --> J
    I --> J
    J --> K[Reportes Finales]
\end{verbatim}


\subsection{Métricas de Optimización}

\textbf{Función Objetivo} minimiza:

\begin{itemize}
\item \textbf{Movimientos de profesores:} Cambios de salón
\begin{itemize}
\item \textbf{Cambios de piso:} Subir/bajar escaleras
\begin{itemize}
\item \textbf{Distancia total:} Recorrido acumulado
\begin{itemize}
\item \textbf{Penalizaciones:} Violaciones de restricciones soft

\subsection{Resultados Típicos}


\section{Contexto del Problema}

\textbf{Asignación Óptima de Salones}


\subsection{1. Introducción - Planteamiento del Problema}

La asignación de salones en instituciones educativas es un \textbf{problema de optimización combinatoria NP-completo} que surge de la necesidad de distribuir eficientemente los espacios físicos disponibles entre las diferentes actividades académicas programadas.

\textbf{Complejidad:}

\begin{itemize}
\item Múltiples restricciones simultáneas
\begin{itemize}
\item Objetivos conflictivos
\begin{itemize}
\item Espacio de búsqueda exponencial

\subsection{Contexto Institucional - ITCM}

\textbf{Recursos Disponibles:}

\begin{itemize}
\item \textbf{Salones de Teoría:} 13 aulas (FF1-FF7 en planta baja, FF8-FFD en planta alta)
\begin{itemize}
\item \textbf{Laboratorios:} 8 espacios especializados (LBD, LBD2, LCA, LCG1, LCG2, LIA, LR, LSO)
\begin{itemize}
\item \textbf{Salones Inválidos:} 5 espacios no utilizables (AV1, AV2, AV4, AV5, E11)
\textbf{Demanda:}

\begin{itemize}
\item \textbf{Clases totales:} ~680 sesiones semanales
\begin{itemize}
\item \textbf{Materias:} 37 diferentes
\begin{itemize}
\item \textbf{Grupos:} Distribuidos en 9 semestres
\begin{itemize}
\item \textbf{Profesores:} ~30 docentes con preferencias específicas

\subsection{Complejidad del Problema}

El espacio de búsqueda para este problema es:

\begin{equation}
|\Omega| = m^n = 21^{680} \approx 10^{900}
\end{equation}

Donde:

\begin{itemize}
\item $m = 21$ salones disponibles
\begin{itemize}
\item $n = 680$ clases a asignar
\textbf{Perspectiva:}

\begin{itemize}
\item \textbf{Átomos en el universo observable:} $\approx 10^{80}$
\begin{itemize}
\item \textbf{Combinaciones posibles:} $\approx 10^{900}$
\begin{itemize}
\item \textbf{Relación:} $10^{820}$ veces más combinaciones que átomos
 Esto hace \textbf{imposible} la búsqueda exhaustiva


\subsection{2. Formulación Matemática - Conjuntos Básicos}

\begin{equation}
\begin{align}
C &= \{c_1, c_2, ..., c_n\} &&\text{Conjunto de clases} \\
S &= \{s_1, s_2, ..., s_m\} &&\text{Conjunto de salones} \\
P &= \{p_1, p_2, ..., p_k\} &&\text{Conjunto de profesores} \\
D &= \{\text{Lunes, Martes, ..., Viernes}\} &&\text{Dias de la semana} \\
H &= \{0700, 0800, ..., 2100\} &&\text{Bloques horarios}
\end{align}
\end{equation}


\textbf{Donde:}

\begin{itemize}
\item $n = 680$ (número total de clases)
\begin{itemize}
\item $m = 21$ (número de salones disponibles)
\begin{itemize}
\item $k \approx 30$ (número de profesores)
\begin{itemize}
\item $\{c_1, c_2, ..., c_n\}$ denota el conjunto de todas las clases
\begin{itemize}
\item $\in$ significa "pertenece a" o "es elemento de"

\subsection{Atributos de Clases}

Para cada clase $c_i \in C$:

\begin{equation}
\begin{align}
dia(c_i) &\in D &&\text{Dia de la semana} \\
hora(c_i) &\in H &&\text{Bloque horario} \\
materia(c_i) &\in M &&\text{Materia} \\
grupo(c_i) &\in G &&\text{Grupo} \\
profesor(c_i) &\in P &&\text{Profesor asignado} \\
tipo(c_i) &\in \{\text{Teoria, Laboratorio}\} &&\text{Tipo de clase} \\
estudiantes(c_i) &\in \mathbb{N} &&\text{Numero de estudiantes}
\end{align}
\end{equation}


\textbf{Donde:}

\begin{itemize}
\item $c_i$ = clase individual (ejemplo: $c_1, c_2, ..., c_{680}$)
\begin{itemize}
\item $dia(c_i)$ = función que retorna el día de la clase $c_i$
\begin{itemize}
\item $\mathbb{N}$ = números naturales (1, 2, 3, ...)
\begin{itemize}
\item $M$ = conjunto de todas las materias
\begin{itemize}
\item $G$ = conjunto de todos los grupos

\subsection{Atributos de Salones}

Para cada salón $s_j \in S$:

\begin{equation}
\begin{align}
capacidad(s_j) &\in \mathbb{N} &&\text{Capacidad maxima} \\
tipo(s_j) &\in \{\text{Teoria, Laboratorio}\} &&\text{Tipo de salon} \\
piso(s_j) &\in \{0, 1\} &&\text{Planta baja o alta} \\
ubicacion(s_j) &\in \mathbb{R}^2 &&\text{Coordenadas fisicas}
\end{align}
\end{equation}


\textbf{Donde:}

\begin{itemize}
\item $s_j$ = salón individual (ejemplo: FF1, FF2, LAB1, ...)
\begin{itemize}
\item $\mathbb{R}^2$ = plano de coordenadas reales (x, y)
\begin{itemize}
\item $\{0, 1\}$ = conjunto con dos elementos (0 = planta baja, 1 = planta alta)

\textbf{Variable de Decisión:}

\begin{equation}
A: C \rightarrow S
\end{equation}

\textbf{Donde:}

\begin{itemize}
\item $A$ = función de asignación
\begin{itemize}
\item $C \rightarrow S$ = función que mapea de clases a salones
\begin{itemize}
\item $A(c_i) = s_j$ significa que la clase $c_i$ se imparte en el salón $s_j$
\begin{itemize}
\item $\rightarrow$ = símbolo de función ("mapea a")

\subsection{Restricciones Duras (Hard Constraints)}

Estas restricciones \textbf{DEBEN} cumplirse para que la solución sea factible.

\textbf{R1. No Conflictos Temporales:}

\begin{equation}
\forall c_i, c_j \in C: \left(dia(c_i) = dia(c_j) \land hora(c_i) = hora(c_j) \land i \neq j\right) \Rightarrow A(c_i) \neq A(c_j)
\end{equation}

\textbf{Donde:}

\begin{itemize}
\item $\forall$ = "para todo" (cuantificador universal)
\begin{itemize}
\item $\land$ = "y" lógico (AND)
\begin{itemize}
\item $\Rightarrow$ = "implica que" (implicación lógica)
\begin{itemize}
\item $\neq$ = "diferente de" (no igual)

\textbf{R2. Capacidad Suficiente:}

\begin{equation}
\forall c_i \in C: estudiantes(c_i) \leq capacidad(A(c_i))
\end{equation}

\textbf{Donde:}

\begin{itemize}
\item $\leq$ = "menor o igual que"
\begin{itemize}
\item $estudiantes(c_i)$ = número de estudiantes en la clase $c_i$
\begin{itemize}
\item $capacidad(A(c_i))$ = capacidad del salón asignado a $c_i$

\subsection{Restricciones Duras (continuación)}

\textbf{R3. Tipo Correcto:}

\begin{equation}
\forall c_i \in C: tipo(c_i) = tipo(A(c_i))
\end{equation}

\textbf{Donde:}

\begin{itemize}
\item $tipo(c_i)$ = tipo de la clase (Teoría o Laboratorio)
\begin{itemize}
\item $tipo(A(c_i))$ = tipo del salón asignado
\begin{itemize}
\item Debe haber coincidencia exacta
\textit{Ejemplo: Una clase de laboratorio debe estar en un laboratorio}


\textbf{R4. Salones Válidos:}

\begin{equation}
\forall c_i \in C: A(c_i) \notin S_{invalidos}
\end{equation}

\textbf{Donde:}

\begin{itemize}
\item $\notin$ = "no pertenece a"
\begin{itemize}
\item $S_{invalidos} = \{AV1, AV2, AV4, AV5, E11\}$ = salones no utilizables

\textbf{R5. Preferencias Prioritarias (PRIORIDAD 1):}

\begin{equation}
\forall c_i \in P_1: A(c_i) = pref(c_i)
\end{equation}

\textbf{Donde:}

\begin{itemize}
\item $P_1$ = conjunto de clases con PRIORIDAD 1
\begin{itemize}
\item $pref(c_i)$ = salón preferido para la clase $c_i$
\begin{itemize}
\item Esta restricción es OBLIGATORIA (100% cumplimiento)

\subsection{Teorema 1: Factibilidad}

\textbf{Teorema 1 (Factibilidad):}

\textit{Una solución es factible si y solo si satisface todas las restricciones duras R1-R5.}

\textbf{Demostración:}

\begin{itemize}
\item $(\Rightarrow)$ Por definición de factibilidad
\begin{itemize}
\item $(\Leftarrow)$ Si se satisfacen R1-R5, no hay violaciones de restricciones obligatorias, por lo tanto la solución es válida

\subsection{Restricciones Suaves (Soft Constraints)}

Estas restricciones son \textbf{deseables} pero no obligatorias.

\textbf{S1. Consistencia de Grupo (PRIORIDAD 2):}

\begin{equation}
minimize \quad \left| \{A(c) : c \in C_g\} \right|
\end{equation}

\textbf{Donde:}

\begin{itemize}
\item $|\cdot|$ = cardinalidad (tamaño del conjunto)
\begin{itemize}
\item $\{A(c) : c \in C_g\}$ = conjunto de salones usados por el grupo $g$
\begin{itemize}
\item $C_g$ = clases del grupo $g$
\begin{itemize}
\item Objetivo: minimizar salones diferentes por grupo
\textit{Minimizar el número de salones diferentes usados por cada grupo}


\textbf{S2. Primer Semestre (PRIORIDAD 3):}

\begin{equation}
maximize \sum_{g \in G_{15}} \sum_{c \in C_g} \mathbb{1}[A(c) = salon_asignado(g)]
\end{equation}

\textbf{Donde:}

\begin{itemize}
\item $\sum$ = sumatoria (suma de todos los elementos)
\begin{itemize}
\item $G_{15}$ = grupos de primer semestre (15xx)
\begin{itemize}
\item $\mathbb{1}[\cdot]$ = función indicadora (1 si verdadero, 0 si falso)
\begin{itemize}
\item $salon\_asignado(g)$ = salón designado para el grupo $g$

\subsection{Restricciones Suaves (continuación)}

\textbf{S3. Minimizar Movimientos:}

Para cada profesor $p \in P$:

\begin{equation}
minimize \sum_{p \in P} \left| \{A(c) : c \in C_p\} \right|
\end{equation}

\textbf{Donde:}

\begin{itemize}
\item $C_p$ = conjunto de clases del profesor $p$
\begin{itemize}
\item $\{A(c) : c \in C_p\}$ = salones diferentes usados por el profesor $p$
\begin{itemize}
\item Objetivo: reducir salones diferentes por profesor
\textbf{Objetivo:} Reducir la cantidad de salones diferentes que usa cada profesor durante el día/semana


\subsection{Función Objetivo Completa}

\begin{equation}
\begin{align}
f(A) = &\ w_1 \cdot  movimientos(A) + w_2 \cdot  cambios_piso(A) \\
       &+ w_3 \cdot  distancia(A) + w_4 \cdot  penalizacion_invalidos(A) \\
       &+ w_5 \cdot  penalizacion_conflictos(A) + w_6 \cdot  penalizacion_tipo(A) \\
       &+ w_7 \cdot  penalizacion_P2(A) + w_8 \cdot  penalizacion_P3(A)
\end{align}
\end{equation}

\textbf{Donde:}

\begin{itemize}
\item $f(A)$ = energía o costo total de la asignación $A$
\begin{itemize}
\item $w_i$ = peso del componente $i$ (importancia relativa)
\begin{itemize}
\item $\cdot$ = multiplicación
\begin{itemize}
\item Cada componente mide un aspecto diferente de la calidad
\textbf{Objetivo:} $minimize\ f(A)$ sujeto a restricciones R1-R5


\subsection{Componente: Movimientos de Profesores}

\begin{equation}
movimientos(A) = \sum_{p \in P} \max\left(0, \left| \{A(c) : c \in C_p\} \right| - 1\right)
\end{equation}

\textbf{Donde:}

\begin{itemize}
\item $\max(a, b)$ = máximo entre $a$ y $b$
\begin{itemize}
\item $|\{A(c) : c \in C_p\}|$ = número de salones diferentes del profesor $p$
\begin{itemize}
\item Si usa $k$ salones, hace $k-1$ movimientos
\begin{itemize}
\item $\max(0, \cdot)$ asegura que el valor no sea negativo

\textbf{Ejemplo Numérico:}

\begin{itemize}
\item Profesor tiene clases en: FF1, FF2, FF1, FF3, FF2
\begin{itemize}
\item Salones únicos: {FF1, FF2, FF3}
\begin{itemize}
\item Movimientos = |{FF1, FF2, FF3}| - 1 = \textbf{2 movimientos}

\subsection{Componente: Cambios de Piso}

Sea $C_p^{ordenado} = [c_1, c_2, ..., c_k]$ las clases del profesor $p$ ordenadas por $(dia, hora)$:

\begin{equation}
cambios_piso(A) = \sum_{p \in P} \sum_{i=1}^{|C_p|-1} \mathbb{1}[piso(A(c_i)) \neq piso(A(c_{i+1}))]
\end{equation}


\textbf{Donde:}

\begin{itemize}
\item $C_p^{ordenado}$ = clases del profesor ordenadas cronológicamente
\begin{itemize}
\item $\sum_{i=1}^{|C_p|-1}$ = suma desde la primera hasta la penúltima clase
\begin{itemize}
\item $c_i, c_{i+1}$ = clases consecutivas en el tiempo
\begin{itemize}
\item $piso(A(c_i))$ = piso del salón asignado a la clase $c_i$
Donde $\mathbb{1}[\cdot]$ es la función indicadora:

\begin{equation}
\mathbb{1}[condicion] = \begin{cases}
1 & \text{si condicion es verdadera} \\
0 & \text{si condicion es falsa}
\end{cases}
\end{equation}


\subsection{Componente: Distancia Total}

Función de distancia entre salones:

\begin{equation}
d(s_i, s_j) = \begin{cases}
0 & \text{si } s_i = s_j \\
1 & \text{si } piso(s_i) = piso(s_j) \land s_i \neq s_j \\
10 & \text{si } piso(s_i) \neq piso(s_j)
\end{cases}
\end{equation}


\textbf{Donde:}

\begin{itemize}
\item $d(s_i, s_j)$ = distancia entre salón $s_i$ y salón $s_j$
\begin{itemize}
\item Mismo salón: distancia 0
\begin{itemize}
\item Mismo piso, diferente salón: distancia 1
\begin{itemize}
\item Diferente piso: distancia 10 (penalización por subir/bajar escaleras)

\textbf{Distancia total:}

\begin{equation}
distancia(A) = \sum_{p \in P} \sum_{i=1}^{|C_p|-1} d(A(c_i), A(c_{i+1}))
\end{equation}

\textbf{Donde:}

\begin{itemize}
\item Se suma la distancia entre salones de clases consecutivas
\begin{itemize}
\item Para cada profesor, se acumula la distancia total recorrida

\subsection{Componente: Penalizaciones}

\begin{equation}
\begin{align}
penalizacion_invalidos(A) &= 1000 \times \sum_{c \in C} \mathbb{1}[A(c) \in S_{invalidos}] \\
penalizacion_conflictos(A) &= 500 \times |\{(c_i, c_j) : conflicto_temporal(c_i, c_j, A)\}| \\
penalizacion_tipo(A) &= 300 \times \sum_{c \in C} \mathbb{1}[tipo(c) \neq tipo(A(c))] \\
penalizacion_P2(A) &= 50 \times \sum_{c \in P_2} \mathbb{1}[A(c) \neq pref(c)] \\
penalizacion_P3(A) &= 25 \times \sum_{c \in P_3} \mathbb{1}[A(c) \neq pref(c)]
\end{align}
\end{equation}


\textbf{Donde:}

\begin{itemize}
\item $\times$ = multiplicación
\begin{itemize}
\item $1000, 500, 300$ = pesos altos para restricciones casi-duras
\begin{itemize}
\item $50, 25$ = pesos menores para preferencias soft
\begin{itemize}
\item $conflicto\_temporal(c_i, c_j, A)$ = verdadero si $c_i$ y $c_j$ están en el mismo salón al mismo tiempo
\begin{itemize}
\item $P_2, P_3$ = conjuntos de clases con prioridad 2 y 3 respectivamente

\subsection{Pesos de la Función Objetivo}


\subsection{Jerarquía de Pesos}

\textbf{Relación de dominancia:}

\begin{equation}
w_4 > w_5 > w_6 \gg w_7 > w_8 \gg w_1 > w_2 > w_3
\end{equation}

\textbf{Interpretación:}

\begin{itemize}
\item \textbf{Restricciones casi-duras} (1000, 500, 300): Dominan sobre todo
\begin{itemize}
\item \textbf{Soft importantes} (50, 25): Prioridad media
\begin{itemize}
\item \textbf{Optimización} (10, 5, 3): Mejora de calidad
\textbf{Garantía:} Una violación de restricción casi-dura siempre pesa más que cualquier combinación razonable de violaciones soft.


\subsection{Teorema 2: Dominancia de Restricciones Duras}

\textbf{Teorema 2:}

\textit{Con los pesos dados, cualquier violación de restricción dura produce una energía mayor que cualquier combinación de violaciones de restricciones suaves.}

\textbf{Demostración:}

\begin{equation}
\begin{align}
E_{hard} &\geq 300 \times v_h \\
E_{soft} &\leq 50 \times v_s + \text{otros terminos suaves}
\end{align}
\end{equation}

Para que $E_{hard} > E_{soft}$ siempre:

\begin{equation}
300 \times v_h > 50 \times v_s \Rightarrow v_h > \frac{v_s}{6}
\end{equation}


\subsection{3. Características del Problema}

\textbf{Clasificación Teórica:}

\begin{itemize}
\item \textbf{Categoría:} Problema de Satisfacción de Restricciones (CSP) con optimización
\begin{itemize}
\item \textbf{Subcategoría:} Problema de Asignación Cuadrática (QAP) generalizado
\begin{itemize}
\item \textbf{Complejidad:} NP-completo (reducible desde Graph Coloring)
\textbf{Propiedades:}

\begin{itemize}
\item \textbf{Espacio de búsqueda:} Discreto, finito, exponencialmente grande
\begin{itemize}
\item \textbf{Función objetivo:} No convexa, múltiples mínimos locales
\begin{itemize}
\item \textbf{Restricciones:} Mezcla de duras y suaves
\begin{itemize}
\item \textbf{Estructura:} Altamente estructurado (temporal, espacial, jerárquico)

\subsection{Teorema 3: NP-Completitud}

\textbf{Teorema 3:}

\textit{El problema de asignación de salones es NP-completo.}

\textbf{Demostración (por reducción desde Graph Coloring):}

\begin{enumerate}
\item \textbf{Construcción del grafo:}
\begin{itemize}
\item Vértices $V = C$ (clases)
\begin{itemize}
\item Arista $(c_i, c_j) \in E$ si hay conflicto temporal
\begin{itemize}
\item Colores $K = S$ (salones)
\begin{enumerate}
\item \textbf{Equivalencia:}
\begin{itemize}
\item Graph Coloring: vértices adyacentes tienen colores diferentes
\begin{itemize}
\item Asignación: clases en conflicto tienen salones diferentes

\subsection{Análisis de Instancia Real}

\begin{verbatim}
Tamaño del problema:
├── Clases totales: 680
├── Salones disponibles: 21
├── Profesores: ~30
├── Materias: 37
├── Grupos: ~50
└── Bloques horarios/día: 14

Distribución de clases:
├── Teoría: ~450 (66%)
└── Laboratorio: ~230 (34%)

Distribución temporal:
├── Lunes-Viernes: ~136 clases/día
├── Bloques pico: 0800-1200 (60% de clases)
└── Bloques valle: 1900-2100 (5% de clases)
\end{verbatim}


\subsection{Densidad del Grafo de Conflictos}

\begin{equation}
densidad = \frac{|E|}{|V|(|V|-1)/2} = \frac{conflictos}{680 \times 679 / 2} \approx 0.15
\end{equation}

Esto indica un grafo \textbf{relativamente disperso}, lo cual es favorable para algoritmos heurísticos.


\subsection{Preferencias del Sistema}

\begin{verbatim}
Preferencias:
├── PRIORIDAD 1: 88 clases (13%)
│   └── Garantía: 100% cumplimiento
├── PRIORIDAD 2: Variable por grupo
│   └── Consistencia de salones
└── PRIORIDAD 3: ~80 clases (12%)
    └── Preferencias de primer semestre
\end{verbatim}


\subsection{Desafíos Específicos}

\begin{enumerate}
\item \textbf{Heterogeneidad de Restricciones:}
\begin{itemize}
\item Mezcla de restricciones duras y suaves
\begin{itemize}
\item Prioridades jerárquicas estrictas
\begin{itemize}
\item Objetivos conflictivos
\begin{enumerate}
\item \textbf{Escala del Problema:}
\begin{itemize}
\item 680 variables de decisión
\begin{itemize}
\item $21^{680}$ combinaciones posibles
\begin{itemize}
\item Evaluación de función objetivo: $O(n)$ por solución

\subsection{Desafíos Específicos (continuación)}

\begin{enumerate}
\item \textbf{Estructura Temporal:}
\begin{itemize}
\item Dependencias secuenciales (clases del mismo profesor)
\begin{itemize}
\item Ventanas temporales fijas
\begin{itemize}
\item No se puede modificar horarios, solo salones
\begin{enumerate}
\item \textbf{Múltiples Stakeholders:}
\begin{itemize}
\item Profesores (preferencias)
\begin{itemize}
\item Estudiantes (consistencia de grupo)
\begin{itemize}
\item Administración (eficiencia de recursos)

\subsection{4. Estado del Arte - Revisión de Literatura (2018-2025)}


\subsubsection{Artículo 1: Genetic Algorithms for Timetabling}

\textbf{Autores:} Pillay, N., & Qu, R. (2018)  

\textbf{Fuente:} Springer - Hyper-Heuristics  

\textbf{Enfoque:} Algoritmos genéticos con operadores adaptativos  

\textbf{Resultados:} Mejora del 15-20% vs. algoritmos tradicionales


\subsubsection{Artículo 1: Genetic Algorithms (Análisis)}

\textbf{Fortalezas:}

\begin{itemize}
\item  Manejo efectivo de restricciones duras y suaves
\begin{itemize}
\item  Operadores de cruce especializados
\begin{itemize}
\item  Buena escalabilidad
\textbf{Debilidades:}

\begin{itemize}
\item  Tiempo de ejecución alto (>5 min para 500 clases)
\begin{itemize}
\item  Requiere ajuste manual de parámetros
\begin{itemize}
\item  No garantiza cumplimiento 100% de prioridades

\subsubsection{Artículo 2: Machine Learning for Timetabling}

\textbf{Autores:} Kristiansen, S., Sørensen, M., & Stidsen, T. (2020)  

\textbf{Fuente:} European Journal of Operational Research  

\textbf{Enfoque:} Random Forest + Reinforcement Learning  

\textbf{Resultados:} 85% de precisión en predicción de asignaciones óptimas


\subsubsection{Artículo 2: Machine Learning (Análisis)}

\textbf{Fortalezas:}

\begin{itemize}
\item  Aprende de soluciones históricas
\begin{itemize}
\item  Rápido en predicción (<10s)
\begin{itemize}
\item  Adaptable a diferentes instituciones
\textbf{Debilidades:}

\begin{itemize}
\item  Requiere dataset de entrenamiento grande
\begin{itemize}
\item  No maneja restricciones nuevas sin reentrenamiento
\begin{itemize}
\item  Calidad depende de datos históricos

\subsubsection{Artículo 3: Hybrid Metaheuristics}

\textbf{Autores:} Bellio, R., Ceschia, S., Di Gaspero, L., & Schaerf, A. (2021)  

\textbf{Fuente:} Computers & Operations Research  

\textbf{Enfoque:} Simulated Annealing + Tabu Search  

\textbf{Resultados:} Top 3 en International Timetabling Competition


\subsubsection{Artículo 3: Hybrid Metaheuristics (Análisis)}

\textbf{Fortalezas:}

\begin{itemize}
\item  Excelente calidad de soluciones
\begin{itemize}
\item  Robusto ante diferentes instancias
\begin{itemize}
\item  Bien documentado
\textbf{Debilidades:}

\begin{itemize}
\item  Complejidad de implementación alta
\begin{itemize}
\item  Muchos parámetros a ajustar
\begin{itemize}
\item  No considera preferencias jerárquicas

\subsubsection{Artículo 4: Greedy with Local Search}

\textbf{Autores:} Burke, E. K., Mareček, J., Parkes, A. J., & Rudová, H. (2019)  

\textbf{Fuente:} Journal of Scheduling  

\textbf{Enfoque:} Construcción greedy + Hill Climbing  

\textbf{Resultados:} Soluciones factibles en <1 minuto


\subsubsection{Artículo 4: Greedy with Local Search (Análisis)}

\textbf{Fortalezas:}

\begin{itemize}
\item  Muy rápido
\begin{itemize}
\item  Fácil de implementar
\begin{itemize}
\item  Buenas soluciones iniciales
\textbf{Debilidades:}

\begin{itemize}
\item  Puede quedar atrapado en óptimos locales
\begin{itemize}
\item  Calidad variable según orden de construcción
\begin{itemize}
\item  No explora ampliamente el espacio de búsqueda

\subsubsection{Artículo 5: Integer Programming}

\textbf{Autores:} Santos, H. G., Uchoa, E., Ochi, L. S., & Maculan, N. (2022)  

\textbf{Fuente:} INFORMS Journal on Computing  

\textbf{Enfoque:} Programación Lineal Entera (ILP)  

\textbf{Resultados:} Soluciones óptimas garantizadas para <200 clases


\subsubsection{Artículo 5: Integer Programming (Análisis)}

\textbf{Fortalezas:}

\begin{itemize}
\item  Garantiza optimalidad
\begin{itemize}
\item  Manejo riguroso de restricciones
\begin{itemize}
\item  Soluciones verificables matemáticamente
\textbf{Debilidades:}

\begin{itemize}
\item  No escala a problemas grandes (>300 clases)
\begin{itemize}
\item  Tiempo exponencial en peor caso
\begin{itemize}
\item  Requiere software especializado (CPLEX, Gurobi)

\subsubsection{Artículo 6: Deep Reinforcement Learning}

\textbf{Autores:} Zhang, C., Song, W., Cao, Z., et al. (2023)  

\textbf{Fuente:} IEEE Transactions on Neural Networks  

\textbf{Enfoque:} Deep Q-Learning con Graph Neural Networks  

\textbf{Resultados:} 92% de eficiencia vs. métodos tradicionales


\subsubsection{Artículo 6: Deep Reinforcement Learning (Análisis)}

\textbf{Fortalezas:}

\begin{itemize}
\item  Estado del arte en ML
\begin{itemize}
\item  Maneja incertidumbre
\begin{itemize}
\item  Aprende políticas generalizables
\textbf{Debilidades:}

\begin{itemize}
\item  Requiere GPU para entrenamiento
\begin{itemize}
\item  Caja negra (difícil de interpretar)
\begin{itemize}
\item  Necesita miles de episodios de entrenamiento

\subsubsection{Artículo 7: Multi-Objective Evolution}

\textbf{Autores:} Fonseca, G. H., Santos, H. G., & Carrano, E. G. (2020)  

\textbf{Fuente:} Applied Soft Computing  

\textbf{Enfoque:} NSGA-II para optimización multi-objetivo  

\textbf{Resultados:} Frente de Pareto con 50+ soluciones no-dominadas


\subsubsection{Artículo 7: Multi-Objective Evolution (Análisis)}

\textbf{Fortalezas:}

\begin{itemize}
\item  Explora trade-offs entre objetivos
\begin{itemize}
\item  Ofrece múltiples soluciones al usuario
\begin{itemize}
\item  Flexible
\textbf{Debilidades:}

\begin{itemize}
\item  Difícil seleccionar solución final
\begin{itemize}
\item  Computacionalmente costoso
\begin{itemize}
\item  Requiere normalización de objetivos

\subsubsection{Artículo 8: Constraint Programming}

\textbf{Autores:} Müller, T., & Murray, K. (2021)  

\textbf{Fuente:} Constraints Journal  

\textbf{Enfoque:} Constraint Satisfaction Problem (CSP)  

\textbf{Resultados:} 98% de restricciones satisfechas


\subsubsection{Artículo 8: Constraint Programming (Análisis)}

\textbf{Fortalezas:}

\begin{itemize}
\item  Modelado declarativo natural
\begin{itemize}
\item  Propagación automática de restricciones
\begin{itemize}
\item  Bueno para problemas altamente restringidos
\textbf{Debilidades:}

\begin{itemize}
\item  Puede no encontrar solución si es muy restringido
\begin{itemize}
\item  Optimización limitada
\begin{itemize}
\item  Requiere expertise en CP

\subsubsection{Artículo 9: Adaptive Large Neighborhood Search}

\textbf{Autores:} Sørensen, M., & Dahms, F. H. (2022)  

\textbf{Fuente:} European Journal of Operational Research  

\textbf{Enfoque:} ALNS con múltiples operadores  

\textbf{Resultados:} Mejora del 25% en calidad vs. métodos clásicos


\subsubsection{Artículo 9: ALNS (Análisis)}

\textbf{Fortalezas:}

\begin{itemize}
\item  Muy efectivo en problemas grandes
\begin{itemize}
\item  Auto-adaptativo
\begin{itemize}
\item  Balance exploración/explotación
\textbf{Debilidades:}

\begin{itemize}
\item  Implementación compleja
\begin{itemize}
\item  Muchos operadores a diseñar
\begin{itemize}
\item  Sensible a configuración inicial

\subsubsection{Artículo 10: Hybrid Genetic Algorithm}

\textbf{Autores:} Tan, J. S., Goh, S. L., Kendall, G., & Sabar, N. R. (2023)  

\textbf{Fuente:} Expert Systems with Applications  

\textbf{Enfoque:} GA + Simulated Annealing  

\textbf{Resultados:} 95% de satisfacción de preferencias


\subsubsection{Artículo 10: Hybrid GA (Análisis)}

\textbf{Fortalezas:}

\begin{itemize}
\item  Combina exploración global y local
\begin{itemize}
\item  Maneja preferencias soft
\begin{itemize}
\item  Resultados consistentes
\textbf{Debilidades:}

\begin{itemize}
\item  Dos conjuntos de parámetros a ajustar
\begin{itemize}
\item  Tiempo de ejecución medio-alto
\begin{itemize}
\item  No garantiza cumplimiento total de prioridades

\subsubsection{Artículo 11: Graph Coloring}

\textbf{Autores:} Lewis, R., & Thompson, J. (2019)  

\textbf{Fuente:} Discrete Applied Mathematics  

\textbf{Enfoque:} Graph coloring con backtracking  

\textbf{Resultados:} Soluciones óptimas para grafos con <500 nodos


\subsubsection{Artículo 11: Graph Coloring (Análisis)}

\textbf{Fortalezas:}

\begin{itemize}
\item  Fundamentación teórica sólida
\begin{itemize}
\item  Algoritmos bien estudiados
\begin{itemize}
\item  Garantías de correctitud
\textbf{Debilidades:}

\begin{itemize}
\item  Modelado limitado (solo conflictos temporales)
\begin{itemize}
\item  No captura preferencias
\begin{itemize}
\item  Escalabilidad limitada

\subsubsection{Artículo 12: Memetic Algorithms}

\textbf{Autores:} Qu, R., Burke, E. K., & McCollum, B. (2020)  

\textbf{Fuente:} Annals of Operations Research  

\textbf{Enfoque:} Algoritmo memético (GA + búsqueda local)  

\textbf{Resultados:} Top 5 en ITC 2019 benchmark


\subsubsection{Artículo 12: Memetic Algorithms (Análisis)}

\textbf{Fortalezas:}

\begin{itemize}
\item  Balance entre diversidad y calidad
\begin{itemize}
\item  Búsqueda local mejora individuos
\begin{itemize}
\item  Robusto
\textbf{Debilidades:}

\begin{itemize}
\item  Computacionalmente intensivo
\begin{itemize}
\item  Requiere diseño cuidadoso de operadores
\begin{itemize}
\item  Convergencia lenta

\subsubsection{Artículo 13: Particle Swarm Optimization}

\textbf{Autores:} Shiau, D. F. (2021)  

\textbf{Fuente:} Applied Intelligence  

\textbf{Enfoque:} PSO con velocidad adaptativa  

\textbf{Resultados:} Convergencia rápida en <100 iteraciones


\subsubsection{Artículo 13: PSO (Análisis)}

\textbf{Fortalezas:}

\begin{itemize}
\item  Pocos parámetros
\begin{itemize}
\item  Fácil de implementar
\begin{itemize}
\item  Buena convergencia
\textbf{Debilidades:}

\begin{itemize}
\item  Puede converger prematuramente
\begin{itemize}
\item  Difícil manejar restricciones duras
\begin{itemize}
\item  Representación de soluciones no trivial

\subsubsection{Artículo 14: Variable Neighborhood Search}

\textbf{Autores:} Sánchez-Oro, J., Sevaux, M., Rossi, A., & Martí, R. (2022)  

\textbf{Fuente:} Computers & Operations Research  

\textbf{Enfoque:} VNS con múltiples vecindarios  

\textbf{Resultados:} 30% mejor que búsqueda local simple


\subsubsection{Artículo 14: VNS (Análisis)}

\textbf{Fortalezas:}

\begin{itemize}
\item  Escapa óptimos locales sistemáticamente
\begin{itemize}
\item  Flexible en definición de vecindarios
\begin{itemize}
\item  No requiere parámetros complejos
\textbf{Debilidades:}

\begin{itemize}
\item  Diseño de vecindarios es crítico
\begin{itemize}
\item  Puede ser lento si vecindarios son grandes
\begin{itemize}
\item  No hay garantías teóricas

\subsubsection{Artículo 15: Ant Colony Optimization}

\textbf{Autores:} Socha, K., Knowles, J., & Samples, M. (2019)  

\textbf{Fuente:} Swarm Intelligence  

\textbf{Enfoque:} ACO con feromonas adaptativas  

\textbf{Resultados:} Buenas soluciones en tiempo razonable


\subsubsection{Artículo 15: ACO (Análisis)}

\textbf{Fortalezas:}

\begin{itemize}
\item  Inspiración biológica interesante
\begin{itemize}
\item  Encuentra múltiples soluciones
\begin{itemize}
\item  Paralelizable
\textbf{Debilidades:}

\begin{itemize}
\item  Muchos parámetros (α, β, ρ, Q)
\begin{itemize}
\item  Convergencia puede ser lenta
\begin{itemize}
\item  Difícil ajustar para problemas específicos

\subsection{Tabla Comparativa - Estado del Arte (1/4)}


\subsection{Tabla Comparativa - Estado del Arte (2/4)}


\subsection{Tabla Comparativa - Estado del Arte (3/4)}


\subsection{Tabla Comparativa - Estado del Arte (4/4)}


\subsection{Nuestra Solución vs. Estado del Arte}


\subsection{Gaps Identificados en la Literatura}

\textbf{1. Prioridades Jerárquicas Estrictas}

\begin{itemize}
\item  Mayoría trata todas las restricciones soft con pesos
\begin{itemize}
\item  No hay garantía absoluta de cumplimiento de preferencias críticas
\begin{itemize}
\item  \textbf{Nuestro enfoque:} Pre-asignación forzada de PRIORIDAD 1
\textbf{2. Corrección Post-Optimización}

\begin{itemize}
\item  Pocos trabajos verifican y corrigen violaciones después
\begin{itemize}
\item  Asumen que el optimizador respeta todas las restricciones
\begin{itemize}
\item  \textbf{Nuestro enfoque:} Módulo de corrección automática

\subsection{Gaps Identificados (continuación)}

\textbf{3. Comparación Multi-Algoritmo}

\begin{itemize}
\item  Mayoría compara contra 1-2 baselines
\begin{itemize}
\item  No hay evaluación sistemática de múltiples enfoques
\begin{itemize}
\item  \textbf{Nuestro enfoque:} 4 algoritmos en misma instancia
\textbf{4. Métricas Específicas de Profesores}

\begin{itemize}
\item  Enfoque típico: minimizar conflictos generales
\begin{itemize}
\item  Poco énfasis en bienestar del profesor
\begin{itemize}
\item  \textbf{Nuestro enfoque:} Movimientos, cambios de piso, distancia

\textbf{5. Validación Estadística}

\begin{itemize}
\item  Muchos reportan 1 corrida o promedio simple
\begin{itemize}
\item  Falta análisis estadístico riguroso
\begin{itemize}
\item  \textbf{Nuestro enfoque:} 30+ corridas con pruebas estadísticas

\subsection{Contribuciones Únicas de Nuestra Solución}

\subsubsection{1. Sistema de Prioridades Jerárquico con Garantías}

\begin{itemize}
\item Pre-asignación forzada de P1 (100% garantizado)
\begin{itemize}
\item Corrección post-optimización automática
\begin{itemize}
\item Índices inmutables durante optimización
\begin{itemize}
\item \textbf{Resultado: ÚNICO en la literatura para problemas >600 clases}
\subsubsection{2. Enfoque Multi-Algoritmo Comparativo}

\begin{itemize}
\item 4 algoritmos diferentes (Greedy+HC, ML, Genético, Baseline)
\begin{itemize}
\item Evaluación en misma instancia real
\begin{itemize}
\item Análisis estadístico riguroso (ANOVA + post-hoc)

\subsection{Contribuciones Únicas (continuación)}

\subsubsection{3. Métricas Centradas en el Profesor}

\begin{itemize}
\item Movimientos entre salones
\begin{itemize}
\item Cambios de piso
\begin{itemize}
\item Distancia total recorrida
\begin{itemize}
\item Impacto directo en bienestar docente
\subsubsection{4. Sistema Completo Funcional}

\begin{itemize}
\item Interfaz gráfica (configurador_materias.py)
\begin{itemize}
\item Aplicación web (en desarrollo)
\begin{itemize}
\item Datos reales validados (ITCM, 680 clases)
\begin{itemize}
\item Documentación completa

\subsubsection{5. Validación Estadística Rigurosa}

\begin{itemize}
\item 30+ corridas por algoritmo
\begin{itemize}
\item Pruebas de normalidad, ANOVA, post-hoc
\begin{itemize}
\item Intervalos de confianza
\begin{itemize}
\item Tamaños de efecto

\subsection{Análisis por Categorías}

\subsubsection{Velocidad de Ejecución}

\textbf{Top 3 Más Rápidos:}

\begin{enumerate}
\item Kristiansen et al. (2020) - ML: <10s
\begin{enumerate}
\item Burke et al. (2019) - Greedy+HC: <1 min
\begin{enumerate}
\item \textbf{NUESTRO - Greedy+HC: ~30s} ✅
\textbf{Más Lentos:}

\begin{itemize}
\item Zhang et al. (2023) - Deep RL: Horas de entrenamiento
\begin{itemize}
\item Qu et al. (2020) - Memetic: 5-8 min
\begin{itemize}
\item Sørensen & Dahms (2022) - ALNS: 5-10 min

\subsection{Análisis por Categorías (continuación)}

\subsubsection{Calidad de Soluciones}

\textbf{Mejor Calidad:}

\begin{enumerate}
\item Santos et al. (2022) - ILP: Óptima (pero no escala)
\begin{enumerate}
\item Bellio et al. (2021) - SA+Tabu: Muy Alta
\begin{enumerate}
\item Sørensen & Dahms (2022) - ALNS: Muy Alta
\textbf{Nuestra Posición:}

\begin{itemize}
\item Greedy+HC: Alta calidad, excelente balance velocidad/calidad
\begin{itemize}
\item ML: Media-Alta, muy rápido
\begin{itemize}
\item Genético: Alta calidad, exploración amplia

\subsection{Análisis por Categorías (continuación 2)}

\subsubsection{Garantías de Prioridades}

\textbf{Con Garantías:}

\begin{enumerate}
\item Santos et al. (2022) - ILP: Sí (pero limitado a <200 clases)
\begin{enumerate}
\item \textbf{NUESTRO: Sí (680 clases)} ✅ \textbf{ÚNICO EN SU CATEGORÍA}
\textbf{Sin Garantías:}

\begin{itemize}
\item Todos los demás enfoques metaheurísticos
\begin{itemize}
\item Tan et al. (2023): 95% pero no garantizado

\subsection{Posicionamiento Final}

\textbf{Nuestra solución se posiciona como un enfoque híbrido práctico que combina:}

 Garantías formales (como ILP) pero escalable  

 Velocidad (como Greedy) pero con calidad  

 Exploración (como GA) pero con eficiencia  

 Validación rigurosa (como investigación académica) pero aplicado

\textbf{Contribución Principal:}

> Primer sistema documentado que garantiza 100% de cumplimiento de prioridades críticas en problemas de >600 clases, con validación estadística completa y múltiples algoritmos comparados en la misma instancia real.


\subsection{Referencias Bibliográficas (1/3)}

\begin{enumerate}
\item Pillay, N., & Qu, R. (2018). \textit{Hyper-Heuristics: Theory and Applications}. Springer.
\begin{enumerate}
\item Kristiansen, S., Sørensen, M., & Stidsen, T. (2020). Machine Learning for Educational Timetabling. \textit{European Journal of Operational Research}, 287(2), 720-735.
\begin{enumerate}
\item Bellio, R., Ceschia, S., Di Gaspero, L., & Schaerf, A. (2021). Hybrid Metaheuristics for Course Timetabling. \textit{Computers & Operations Research}, 131, 105070.
\begin{enumerate}
\item Burke, E. K., Mareček, J., Parkes, A. J., & Rudová, H. (2019). Greedy Heuristics with Local Search. \textit{Journal of Scheduling}, 22(4), 449-466.
\begin{enumerate}
\item Santos, H. G., Uchoa, E., Ochi, L. S., & Maculan, N. (2022). Integer Programming for Classroom Assignment. \textit{INFORMS Journal on Computing}, 34(2), 1142-1158.

\subsection{Referencias Bibliográficas (2/3)}

\begin{enumerate}
\item Zhang, C., Song, W., Cao, Z., et al. (2023). Deep Reinforcement Learning for Scheduling. \textit{IEEE Transactions on Neural Networks}, 34(8), 4567-4580.
\begin{enumerate}
\item Fonseca, G. H., Santos, H. G., & Carrano, E. G. (2020). Multi-Objective Evolutionary Algorithms. \textit{Applied Soft Computing}, 95, 106456.
\begin{enumerate}
\item Müller, T., & Murray, K. (2021). Constraint Programming Approaches. \textit{Constraints}, 26(3), 321-345.
\begin{enumerate}
\item Sørensen, M., & Dahms, F. H. (2022). Adaptive Large Neighborhood Search. \textit{European Journal of Operational Research}, 298(3), 1045-1060.
\begin{enumerate}
\item Tan, J. S., Goh, S. L., Kendall, G., & Sabar, N. R. (2023). Hybrid Genetic Algorithm. \textit{Expert Systems with Applications}, 213, 119876.

\subsection{Referencias Bibliográficas (3/3)}

\begin{enumerate}
\item Lewis, R., & Thompson, J. (2019). Graph Coloring for Timetabling. \textit{Discrete Applied Mathematics}, 265, 112-128.
\begin{enumerate}
\item Qu, R., Burke, E. K., & McCollum, B. (2020). Memetic Algorithms. \textit{Annals of Operations Research}, 293(2), 567-590.
\begin{enumerate}
\item Shiau, D. F. (2021). Particle Swarm Optimization. \textit{Applied Intelligence}, 51(8), 5678-5692.
\begin{enumerate}
\item Sánchez-Oro, J., Sevaux, M., Rossi, A., & Martí, R. (2022). Variable Neighborhood Search. \textit{Computers & Operations Research}, 142, 105789.
\begin{enumerate}
\item Socha, K., Knowles, J., & Samples, M. (2019). Ant Colony Optimization. \textit{Swarm Intelligence}, 13(2), 167-189.

\section{Teoría y Fundamentos Matemáticos}

\textbf{Modelado Formal del Problema}


\subsection{Glosario (1/4)}

\subsubsection{Conjuntos Básicos}


\subsection{Glosario (2/4)}

\subsubsection{Variables y Atributos Principales (I)}


\subsection{Glosario (3/4)}

\subsubsection{Variables y Atributos Principales (II)}


\subsection{Glosario (4/4)}

\subsubsection{Conjuntos Derivados}


\subsection{1. Modelado del Problema}


\subsubsection{1.1 Definición Formal}

El problema de asignación de salones es un \textbf{problema de optimización combinatoria} que puede modelarse como:

\textbf{Entrada:}

\begin{itemize}
\item Conjunto de clases $C = \{c_1, c_2, ..., c_n\}$
\begin{itemize}
\item Conjunto de salones $S = \{s_1, s_2, ..., s_m\}$
\begin{itemize}
\item Conjunto de restricciones $R$
\begin{itemize}
\item Función de costo $f: C \times S \rightarrow \mathbb{R}$

\textbf{Donde:}

\begin{itemize}
\item $C, S$ = conjuntos de clases y salones
\begin{itemize}
\item $R$ = conjunto de restricciones (duras y suaves)
\begin{itemize}
\item $f: C \times S \rightarrow \mathbb{R}$ = función que asigna un costo real a cada par (clase, salón)
\begin{itemize}
\item $\times$ = producto cartesiano
\begin{itemize}
\item $\mathbb{R}$ = números reales
\textbf{Salida:}

\begin{itemize}
\item Asignación $A: C \rightarrow S$ que minimiza $f$ sujeto a $R$

\textbf{Donde:}

\begin{itemize}
\item $A$ = función de asignación que mapea cada clase a un salón
\begin{itemize}
\item "minimiza $f$" = encuentra la asignación con menor costo
\begin{itemize}
\item "sujeto a $R$" = respetando todas las restricciones

\subsubsection{1.2 Clasificación del Problema}

Este problema pertenece a la familia de \textbf{problemas NP-difíciles}, específicamente:

\begin{itemize}
\item \textbf{Tipo:} Problema de asignación con restricciones (Constraint Satisfaction Problem - CSP)
\begin{itemize}
\item \textbf{Complejidad:} NP-completo
\begin{itemize}
\item \textbf{Espacio de búsqueda:} $O(m^n)$ donde $m$ = salones, $n$ = clases
\begin{itemize}
\item \textbf{Ejemplo:} Para 680 clases y 21 salones: $21^{680} \approx 10^{900}$ combinaciones

\textbf{Donde:}

\begin{itemize}
\item $O(m^n)$ = notación Big-O (orden de magnitud)
\begin{itemize}
\item $m^n$ = $m$ elevado a la potencia $n$
\begin{itemize}
\item $21^{680}$ = 21 opciones para cada una de las 680 clases
\begin{itemize}
\item $\approx$ = "aproximadamente"
\begin{itemize}
\item $10^{900}$ = 1 seguido de 900 ceros (número astronómicamente grande)

\subsubsection{1.3 Restricciones del Sistema}

\paragraph{Restricciones Duras (Hard Constraints)}

\begin{enumerate}
\item \textbf{Unicidad temporal:} Una clase solo puede estar en un salón a la vez
\begin{equation}

   **Donde:** $\forall$ = para todo, $\land$ = y lógico, $\Rightarrow$ = implica, $\neq$ = diferente

2. **Capacidad:** El salón debe tener capacidad suficiente
\end{equation}

   \textbf{Donde:} $\geq$ = mayor o igual que


\begin{enumerate}
\item \textbf{Tipo de salón:} Debe coincidir con el tipo de clase
\begin{equation}

   **Donde:** $tipo(A(c))$ = tipo del salón asignado, $tipo_requerido(c)$ = tipo que necesita la clase

4. **Preferencias prioritarias (P1):** Cumplimiento obligatorio al 100%
\end{equation}

   \textbf{Donde:} $P1$ = conjunto de clases con prioridad 1, $salon\_preferido(c)$ = salón preferido


\paragraph{Restricciones Suaves (Soft Constraints)}

\begin{enumerate}
\item \textbf{Consistencia de grupo (P2):} Minimizar cambios de salón por grupo
\begin{equation}

   **Donde:** $\sum$ = sumatoria, $|\cdot |$ = cardinalidad (tamaño del conjunto), $salones_distintos(g)$ = salones diferentes usados por el grupo $g$

2. **Primer semestre (P3):** Asignar grupos 15xx a salones específicos
\end{equation}

   \textbf{Donde:} $maximize$ = maximizar, $\mathbb{1}[\cdot]$ = función indicadora (1 si verdadero, 0 si falso), $Grupos15xx$ = grupos de primer semestre


\subsection{2. Función Objetivo}


\subsubsection{2.1 Formulación General}

La función objetivo combina múltiples métricas con pesos:

\begin{equation}
E(A) = w_1 \cdot  movimientos(A) + w_2 \cdot  cambios_piso(A) + w_3 \cdot  distancia(A) + \sum_{i} w_i \cdot  penalizacion_i(A)
\end{equation}


\textbf{Donde:}

\begin{itemize}
\item $E(A)$ = energía/costo total de la asignación $A$
\begin{itemize}
\item $w_i$ = pesos de cada componente (importancia relativa)
\begin{itemize}
\item $\cdot$ = multiplicación
\begin{itemize}
\item $movimientos(A)$ = número de movimientos entre salones
\begin{itemize}
\item $cambios\_piso(A)$ = número de cambios de piso
\begin{itemize}
\item $distancia(A)$ = distancia total recorrida
\begin{itemize}
\item $\sum_{i} w_i \cdot penalizacion_i(A)$ = suma de todas las penalizaciones
\textbf{Objetivo:} $minimize\ E(A)$ = encontrar la asignación con menor costo


\subsubsection{2.2 Componentes de la Función}

\paragraph{Movimientos de Profesores}

\begin{equation}
movimientos(A) = \sum_{p \in Profesores} \left| \{A(c) : c \in clases(p)\} \right| - 1
\end{equation}

\textbf{Donde:}

\begin{itemize}
\item $\sum_{p \in Profesores}$ = suma sobre todos los profesores
\begin{itemize}
\item $|\cdot|$ = cardinalidad (número de elementos del conjunto)
\begin{itemize}
\item $\{A(c) : c \in clases(p)\}$ = conjunto de salones usados por el profesor $p$
\begin{itemize}
\item $clases(p)$ = clases que imparte el profesor $p$
\begin{itemize}
\item $-1$ = si usa $k$ salones, hace $k-1$ movimientos
\textbf{Interpretación:} Número de veces que cada profesor cambia de salón


\paragraph{Cambios de Piso}

\begin{equation}
cambios_piso(A) = \sum_{p \in Profesores} \sum_{i=1}^{|clases(p)|-1} \mathbb{1}[piso(A(c_i)) \neq piso(A(c_{i+1}))]
\end{equation}

\textbf{Donde:}

\begin{itemize}
\item $\sum_{i=1}^{|clases(p)|-1}$ = suma desde la primera hasta la penúltima clase del profesor
\begin{itemize}
\item $\mathbb{1}[\cdot]$ = función indicadora (1 si verdadero, 0 si falso)
\begin{itemize}
\item $c_i, c_{i+1}$ = clases consecutivas (en orden cronológico)
\begin{itemize}
\item $piso(A(c_i))$ = piso del salón asignado a la clase $c_i$
\begin{itemize}
\item $\neq$ = diferente de
\textbf{Interpretación:} Número de veces que cada profesor sube o baja de piso


\paragraph{Distancia Total}

\begin{equation}
distancia(A) = \sum_{p \in Profesores} \sum_{i=1}^{|clases(p)|-1} d(A(c_i), A(c_{i+1}))
\end{equation}

\textbf{Donde:}

\begin{itemize}
\item $d(s_1, s_2)$ = función de distancia entre salón $s_1$ y salón $s_2$
\begin{itemize}
\item Se acumula la distancia entre salones de clases consecutivas
Donde $d(s_1, s_2)$ es la distancia entre salones:

\begin{equation}
d(s_1, s_2) = \begin{cases}
0 & \text{si } s_1 = s_2 \\
1 & \text{si mismo piso, diferente salon} \\
10 & \text{si diferente piso}
\end{cases}
\end{equation}


\subsubsection{2.3 Penalizaciones}

\paragraph{Salones Inválidos}

\begin{equation}
penalizacion_invalidos(A) = 1000 \times \sum_{c \in C} \mathbb{1}[A(c) \in S_{invalidos}]
\end{equation}

\textbf{Donde:}

\begin{itemize}
\item $1000$ = peso alto (penalización severa)
\begin{itemize}
\item $\times$ = multiplicación
\begin{itemize}
\item $S_{invalidos}$ = conjunto de salones no utilizables
\begin{itemize}
\item $\mathbb{1}[A(c) \in S_{invalidos}]$ = 1 si la clase $c$ está en salón inválido, 0 si no

\paragraph{Conflictos de Horario}

\begin{equation}
penalizacion_conflictos(A) = 500 \times |\{(c_i, c_j) : conflicto(c_i, c_j, A)\}|
\end{equation}

\textbf{Donde:}

\begin{itemize}
\item $500$ = peso alto (penalización severa)
\begin{itemize}
\item $|\cdot|$ = cardinalidad (número de elementos)
\begin{itemize}
\item $(c_i, c_j)$ = par de clases en conflicto
\begin{itemize}
\item $conflicto(c_i, c_j, A)$ = verdadero si las clases están en el mismo salón al mismo tiempo

\paragraph{Tipo Incorrecto}

\begin{equation}
penalizacion_tipo(A) = 300 \times \sum_{c \in C} \mathbb{1}[tipo(A(c)) \neq tipo_requerido(c)]
\end{equation}

\textbf{Donde:}

\begin{itemize}
\item $300$ = peso alto (penalización severa)
\begin{itemize}
\item $tipo(A(c))$ = tipo del salón asignado (Teoría o Laboratorio)
\begin{itemize}
\item $tipo\_requerido(c)$ = tipo que necesita la clase
\begin{itemize}
\item $\neq$ = diferente de

\paragraph{Preferencias (P2 y P3)}

\begin{equation}
penalizacion_preferencias(A) = \sum_{c \in P2 \cup P3} w_{prioridad(c)} \times \mathbb{1}[A(c) \neq salon_preferido(c)]
\end{equation}

\textbf{Donde:}

\begin{itemize}
\item $P2 \cup P3$ = unión de clases con prioridad 2 y prioridad 3
\begin{itemize}
\item $\cup$ = unión de conjuntos
\begin{itemize}
\item $w_{prioridad(c)}$ = peso según prioridad (50 para P2, 25 para P3)
\begin{itemize}
\item $salon\_preferido(c)$ = salón preferido para la clase $c$

\subsection{3. Teoremas y Propiedades}


\subsubsection{Teorema 1: Optimalidad Local vs Global}

\textbf{Enunciado:} En el problema de asignación de salones, un óptimo local no garantiza ser óptimo global.

\textbf{Demostración:} 

\begin{itemize}
\item El espacio de búsqueda es no-convexo
\begin{itemize}
\item Existen múltiples mínimos locales
\begin{itemize}
\item Un intercambio de dos clases puede mejorar localmente pero empeorar globalmente
\textbf{Implicación:} Se requieren algoritmos que escapen de óptimos locales (e.g., algoritmos genéticos, simulated annealing)


\subsubsection{Teorema 2: Complejidad Computacional}

\textbf{Enunciado:} El problema de asignación de salones con restricciones es NP-completo.

\textbf{Reducción:} Desde el problema de coloración de grafos:

\begin{itemize}
\item Vértices = Clases
\begin{itemize}
\item Aristas = Conflictos temporales
\begin{itemize}
\item Colores = Salones
\begin{itemize}
\item Restricciones adicionales = Preferencias y capacidades

\subsubsection{Teorema 3: Garantía de Factibilidad}

\textbf{Enunciado:} Si existe al menos una asignación válida que satisface todas las restricciones duras, el algoritmo de pre-asignación garantiza encontrar una solución factible.

\textbf{Demostración:}

\begin{enumerate}
\item Pre-asignación fuerza P1 (restricción dura)
\begin{enumerate}
\item Algoritmo de resolución de conflictos desplaza clases no-prioritarias
\begin{enumerate}
\item Si hay suficientes salones disponibles, siempre existe una asignación válida

\subsection{4. Análisis de Complejidad}


\subsubsection{4.1 Complejidad Temporal}


Donde:

\begin{itemize}
\item $n$ = número de clases
\begin{itemize}
\item $m$ = número de salones
\begin{itemize}
\item $k$ = iteraciones de Hill Climbing
\begin{itemize}
\item $d$ = profundidad de árboles (ML)
\begin{itemize}
\item $g$ = generaciones (Genético)
\begin{itemize}
\item $p$ = tamaño de población (Genético)

\subsubsection{4.2 Complejidad Espacial}


\subsection{5. Convergencia y Garantías}


\subsubsection{5.1 Greedy + Hill Climbing}

\textbf{Garantía:} Converge a un óptimo local en tiempo finito

\textbf{Condición de parada:}

\begin{equation}
\forall vecino \in N(A_{actual}): E(vecino) \geq E(A_{actual})
\end{equation}

\textbf{Donde:}

\begin{itemize}
\item $\forall$ = para todo
\begin{itemize}
\item $N(A_{actual})$ = vecindario de la asignación actual
\begin{itemize}
\item $E(\cdot)$ = función de energía/costo
\begin{itemize}
\item $\geq$ = mayor o igual que
\begin{itemize}
\item \textbf{Interpretación:} Se detiene cuando ningún vecino mejora la solución

\subsubsection{5.2 Algoritmo Genético}

\textbf{Teorema de Convergencia:} Con probabilidad 1, el algoritmo genético con elitismo converge al óptimo global cuando $t \rightarrow \infty$

\textbf{Donde:}

\begin{itemize}
\item $t \rightarrow \infty$ = cuando el tiempo tiende a infinito
\begin{itemize}
\item $\rightarrow$ = "tiende a"
\begin{itemize}
\item $\infty$ = infinito

\textbf{Condiciones:}

\begin{itemize}
\item Mutación con probabilidad $p_m > 0$
\begin{itemize}
\item Elitismo (preservar mejores individuos)
\begin{itemize}
\item Población suficientemente grande
\textbf{Donde:} $p_m$ = probabilidad de mutación, $> 0$ = estrictamente mayor que cero


\subsubsection{5.3 Machine Learning}

\textbf{Garantía:} Minimiza el error de predicción en el conjunto de entrenamiento

\textbf{Error esperado:}

\begin{equation}
E_{error} = \mathbb{E}[(y - \hat{y})^2]
\end{equation}

\textbf{Donde:}

\begin{itemize}
\item $\mathbb{E}[\cdot]$ = valor esperado (promedio)
\begin{itemize}
\item $y$ = asignación óptima (valor real)
\begin{itemize}
\item $\hat{y}$ = predicción del modelo (valor estimado)
\begin{itemize}
\item $(y - \hat{y})^2$ = error cuadrático
\begin{itemize}
\item $\hat{\cdot}$ = símbolo de estimación/predicción

\subsection{6. Heurísticas y Técnicas}


\subsubsection{6.1 Heurística de Construcción Voraz}

\textbf{Criterio de selección:} Para cada clase $c$, elegir salón $s$ que minimiza:

\begin{equation}
score(c, s) = \alpha \cdot  distancia(s, ultimo_salon(profesor(c))) + \beta \cdot  ocupacion(s) + \gamma \cdot  penalizacion(s, c)
\end{equation}

\textbf{Donde:}

\begin{itemize}
\item $score(c, s)$ = puntuación para asignar clase $c$ al salón $s$
\begin{itemize}
\item $\alpha, \beta, \gamma$ = pesos de cada componente
\begin{itemize}
\item $\cdot$ = multiplicación
\begin{itemize}
\item $distancia(s, ultimo\_salon(\cdot))$ = distancia al último salón usado por el profesor
\begin{itemize}
\item $ocupacion(s)$ = nivel de ocupación del salón $s$
\begin{itemize}
\item $penalizacion(s, c)$ = penalización por asignar $c$ a $s$

\subsubsection{6.2 Operadores Genéticos}

\textbf{Cruce (Crossover):} Punto único

\begin{equation}
hijo_1[i] = \begin{cases}
padre_1[i] & \text{si } i < punto_cruce \\
padre_2[i] & \text{si } i \geq punto_cruce
\end{cases}
\end{equation}

\textbf{Donde:}

\begin{itemize}
\item $hijo_1[i]$ = gen $i$ del hijo 1
\begin{itemize}
\item $padre_1[i], padre_2[i]$ = gen $i$ de cada padre
\begin{itemize}
\item $i$ = índice del gen (clase)
\begin{itemize}
\item $punto\_cruce$ = punto donde se divide el cromosoma
\begin{itemize}
\item $<$ = menor que, $\geq$ = mayor o igual que

\textbf{Mutación:} Intercambio aleatorio

\begin{equation}
P(mutar(c)) = p_m \cdot  (1 + \frac{generacion}{max_generaciones})
\end{equation}


\subsubsection{6.3 Búsqueda Local (Hill Climbing)}

\textbf{Vecindario:} Intercambios de clases del mismo tipo

\begin{equation}
N(A) = \{A' : \exists c_i, c_j \in C, tipo(c_i) = tipo(c_j), A'(c_i) = A(c_j), A'(c_j) = A(c_i)\}
\end{equation}

\textbf{Donde:}

\begin{itemize}
\item $N(A)$ = vecindario de la asignación $A$
\begin{itemize}
\item $\exists$ = "existe" (cuantificador existencial)
\begin{itemize}
\item $A'$ = asignación vecina (modificación de $A$)
\begin{itemize}
\item $tipo(c_i) = tipo(c_j)$ = ambas clases del mismo tipo (ambas Teoría o ambas Lab)
\begin{itemize}
\item $A'(c_i) = A(c_j)$ = intercambio de salones entre $c_i$ y $c_j$

\textbf{Criterio de aceptación:} Descenso más pronunciado (steepest descent)

\begin{equation}
A_{nuevo} = \arg\min_{A' \in N(A)} E(A')
\end{equation}

\textbf{Donde:}

\begin{itemize}
\item $\arg\min$ = "argumento que minimiza" (la asignación que da el mínimo)
\begin{itemize}
\item $A_{nuevo}$ = nueva asignación seleccionada
\begin{itemize}
\item $E(A')$ = energía/costo de la asignación $A'$
\begin{itemize}
\item Se elige el vecino con menor energía

\subsection{Referencias}

\begin{enumerate}
\item Garey, M. R., & Johnson, D. S. (1979). \textit{Computers and Intractability: A Guide to the Theory of NP-Completeness}
\begin{enumerate}
\item Russell, S., & Norvig, P. (2020). \textit{Artificial Intelligence: A Modern Approach}
\begin{enumerate}
\item Goldberg, D. E. (1989). \textit{Genetic Algorithms in Search, Optimization, and Machine Learning}
\begin{enumerate}
\item Mitchell, T. M. (1997). \textit{Machine Learning}

\section{Pre-procesamiento y Pre-asignación}

\textbf{Preparación de Datos y Asignación de Prioridades}


\subsection{1. Introducción y Motivación}


\subsubsection{1.1 Problema Fundamental}

En sistemas de optimización tradicionales, las restricciones duras y suaves se manejan mediante penalizaciones en la función objetivo. Sin embargo, este enfoque tiene limitaciones:

\textbf{Problema con Penalizaciones:}

\begin{itemize}
\item No garantiza 100% de cumplimiento
\begin{itemize}
\item Puede sacrificar restricciones duras por optimizar suaves
\begin{itemize}
\item Requiere ajuste cuidadoso de pesos
\textbf{Solución Propuesta:}

\begin{itemize}
\item \textbf{Pre-asignación forzada} de PRIORIDAD 1
\begin{itemize}
\item \textbf{Separación explícita} entre restricciones duras y suaves
\begin{itemize}
\item \textbf{Garantía matemática} de cumplimiento

\subsubsection{1.2 Teorema de Separación}

\textbf{Teorema 1 (Separación de Prioridades):}

\textit{Si las restricciones pueden dividirse en conjuntos disjuntos $R_1$ (duras) y $R_2$ (suaves), entonces existe una solución óptima que satisface completamente $R_1$.}

\textbf{Demostración:}

Sea $S_{factible} = \{s : s \text{ satisface } R_1\}$ el conjunto de soluciones factibles.

\begin{enumerate}
\item Por definición, toda solución válida debe satisfacer $R_1$
\begin{enumerate}
\item El óptimo global $s^* \in S_{factible}$
\begin{enumerate}
\item Por lo tanto, $s^*$ satisface $R_1$ completamente
\begin{enumerate}
\item La optimización de $R_2$ se realiza dentro de $S_{factible}$
\textbf{Implicación:} Podemos pre-asignar $R_1$ y luego optimizar $R_2$ sin afectar $R_1$.


\subsection{2. Arquitectura del Sistema de Pre-asignación}


\subsubsection{2.1 Flujo de Procesamiento}

\begin{verbatim}
graph TD
    A[Horario Inicial] --> B[Cargar Preferencias]
    B --> C[Identificar Clases P1]
    C --> D[Ordenar por Complejidad]
    D --> E[Asignar Forzadamente]
    E --> F{Conflictos?}
    F -->|Sí| G[Resolver Conflictos]
    F -->|No| H[Marcar Inmutables]
    G --> H
    H --> I[Guardar Horario Pre-asignado]
    I --> J[Guardar Índices Inmutables]
\end{verbatim}


\subsubsection{2.2 Componentes del Sistema}

\begin{lstlisting}[language=Python]
class PreAsignadorP1:
    def __init__(self):
        self.preferencias = {}
        self.indices_inmutables = []
        self.conflictos_resueltos = 0
    
    def ejecutar(self, horario_inicial):
        # 1. Cargar configuración
        self.cargar_preferencias()
        
        # 2. Identificar clases prioritarias
        clases_p1 = self.identificar_clases_p1(horario_inicial)
        
        # 3. Ordenar por complejidad
        clases_ordenadas = self.ordenar_por_complejidad(clases_p1)
        
        # 4. Asignar forzadamente
        horario_asignado = self.asignar_forzadamente(
            horario_inicial, clases_ordenadas
        )
        
        # 5. Marcar como inmutables
        self.marcar_inmutables(clases_p1)
        
        # 6. Guardar resultados
        self.guardar_resultados(horario_asignado)
        
        return horario_asignado
\end{verbatim}


\subsection{3. Identificación de Clases Prioritarias}


\subsubsection{3.1 Criterio de Prioridad}

Una clase $c$ es prioritaria si:

\begin{equation}
c \in P_1 \Leftrightarrow \exists profesor(c), materia(c): prioridad(profesor, materia, tipo(c)) = \text{"Prioritario"}
\end{equation}

\textbf{Implementación:}

\begin{lstlisting}[language=Python]
def identificar_clases_p1(self, df):
    """
    Identifica todas las clases con PRIORIDAD 1
    """
    clases_p1 = []
    
    for idx, clase in df.iterrows():
        profesor = clase['Profesor']
        materia = clase['Materia']
        tipo = clase['Tipo_Salon']
\end{verbatim}


\subsubsection{3.1 Criterio de Prioridad (continuación 1)}

\begin{lstlisting}[language=Python]
        # Verificar en preferencias
        if profesor in self.preferencias:
            if materia in self.preferencias[profesor]['materias']:
                pref = self.preferencias[profesor]['materias'][materia]
                
                # Verificar teoría
                if tipo == 'Teoría':
                    if (pref.get('prioridad_teoria') == 'Prioritario' and
                        pref.get('salon_teoria') != 'Sin preferencia'):
                        clases_p1.append({
                            'idx': idx,
                            'clase': clase,
                            'salon_preferido': pref['salon_teoria']
                        })
\end{verbatim}


\subsubsection{3.1 Criterio de Prioridad (continuación 2)}

\begin{lstlisting}[language=Python]
                # Verificar laboratorio
                elif tipo == 'Laboratorio':
                    if (pref.get('prioridad_lab') == 'Prioritario' and
                        pref.get('salon_lab') != 'Sin preferencia'):
                        clases_p1.append({
                            'idx': idx,
                            'clase': clase,
                            'salon_preferido': pref['salon_lab']
                        })
    
    return clases_p1
\end{verbatim}


\subsubsection{3.2 Estadísticas de Prioridad}

Para el caso del Instituto Tecnológico de Ciudad Madero:

\begin{verbatim}
Total clases: 680
Clases PRIORIDAD 1: 88 (13%)
├── Teoría: 58 (66%)
└── Laboratorio: 30 (34%)

Distribución por profesor:
├── PROFESOR 3: 10 clases
├── PROFESOR 4: 8 clases
├── PROFESOR 8: 15 clases
├── PROFESOR 9: 15 clases
├── PROFESOR 20: 9 clases
├── PROFESOR 21: 8 clases
├── PROFESOR 24: 5 clases
└── PROFESOR 26: 18 clases
\end{verbatim}


\subsection{4. Ordenamiento por Complejidad}


\subsubsection{4.1 Función de Complejidad}

Para cada clase prioritaria $c$, definimos su complejidad:

\begin{equation}
complejidad(c) = w_1 \cdot  num_clases_profesor(c) + w_2 \cdot  conflictos_potenciales(c) + w_3 \cdot  \frac{1}{salones_disponibles(c)}
\end{equation}

\textbf{Componentes:}

\begin{enumerate}
\item \textbf{Número de clases del profesor:}
\begin{equation}

2. **Conflictos potenciales:**
\end{equation}

\begin{enumerate}
\item \textbf{Salones disponibles:}
\begin{equation}
---

**Pesos:**
- $w_1 = 10$ (más clases = más complejo)
- $w_2 = 5$ (más conflictos = más complejo)
- $w_3 = 3$ (menos opciones = más complejo)


---

### 4.2 Ordenamiento

\begin{lstlisting}[language=Python]
def ordenar_por_complejidad(self, clases_p1):
    """
    Ordena clases prioritarias por complejidad (más complejo primero)
    """
    def calcular_complejidad(clase_info):
        clase = clase_info['clase']
        
        # Componente 1: Número de clases del profesor
        num_clases = sum(1 for c in clases_p1 
                        if c['clase']['Profesor'] == clase['Profesor'])
        
        # Componente 2: Conflictos potenciales
        conflictos = sum(1 for c in clases_p1
                        if (c['clase']['Dia'] == clase['Dia'] and
                            c['clase']['Bloque_Horario'] == clase['Bloque_Horario'] and
                            c['salon_preferido'] == clase_info['salon_preferido']))
        
        # Componente 3: Inverso de salones disponibles
        salones_disp = len(self.obtener_salones_validos(clase))
        inv_salones = 1.0 / salones_disp if salones_disp > 0 else 10
        
        return 10 * num_clases + 5 * conflictos + 3 * inv_salones
    
    # Ordenar de mayor a menor complejidad
    return sorted(clases_p1, key=calcular_complejidad, reverse=True)
\end{verbatim}


---

**Justificación del Ordenamiento:**

**Lema 1 (Ordenamiento Óptimo):**
*Procesar clases más complejas primero minimiza la probabilidad de infactibilidad.*

**Demostración:**
- Clases complejas tienen menos opciones alternativas
- Si se procesan al final, pueden no tener salones disponibles
- Procesarlas primero garantiza que al menos su opción preferida esté disponible
- Clases simples pueden adaptarse a salones restantes


---

## 5. Asignación Forzada


---

### 5.1 Algoritmo Principal

\begin{lstlisting}[language=Python]
def asignar_forzadamente(self, df, clases_ordenadas):
    """
    Asigna forzadamente cada clase a su salón preferido
    """
    df_resultado = df.copy()
    ocupacion = {}  # (dia, bloque, salon) -> idx
    
    for clase_info in clases_ordenadas:
        idx = clase_info['idx']
        clase = clase_info['clase']
        salon_preferido = clase_info['salon_preferido']
        
        # Clave de ocupación
        key = (clase['Dia'], clase['Bloque_Horario'], salon_preferido)
        
        # Verificar si hay conflicto
        if key in ocupacion:
            # Resolver conflicto
            exito = self.resolver_conflicto(
                df_resultado, idx, salon_preferido, 
                ocupacion, clase_info
            )
            
            if not exito:
                print(f"⚠️  No se pudo asignar clase {idx}")
                continue
        
        # Asignar
        df_resultado.loc[idx, 'Salon'] = salon_preferido
        ocupacion[key] = idx
        self.indices_inmutables.append(idx)
    
    return df_resultado
\end{verbatim}



---

### 5.2 Resolución de Conflictos

Cuando dos clases prioritarias quieren el mismo salón al mismo tiempo:

\begin{lstlisting}[language=Python]
def resolver_conflicto(self, df, idx_nueva, salon, ocupacion, clase_info):
    """
    Resuelve conflicto desplazando clase no-prioritaria
    """
    key = (clase_info['clase']['Dia'], 
           clase_info['clase']['Bloque_Horario'], 
           salon)
    
    idx_ocupante = ocupacion[key]
    clase_ocupante = df.iloc[idx_ocupante]
    
    # Verificar si ocupante es prioritario
    if self.es_prioritaria(clase_ocupante):
        # Ambas son prioritarias: conflicto irresolvable
        print(f"❌ Conflicto entre dos clases prioritarias")
        return False
\end{verbatim}


---

### 5.2 Resolución de Conflictos (continuación)

\begin{lstlisting}[language=Python]
    # Desplazar ocupante a otro salón
    salones_alternativos = self.obtener_salones_validos(clase_ocupante)
    
    for salon_alt in salones_alternativos:
        key_alt = (clase_ocupante['Dia'], 
                   clase_ocupante['Bloque_Horario'], 
                   salon_alt)
        
        if key_alt not in ocupacion:
            # Mover ocupante
            df.loc[idx_ocupante, 'Salon'] = salon_alt
            ocupacion[key_alt] = idx_ocupante
            del ocupacion[key]
            return True
    
    # No hay salones alternativos
    print(f"⚠️  No se encontró salón alternativo")
    return False
\end{verbatim}


---

**Teorema 2 (Resolución de Conflictos):**
*Si existe al menos un salón válido libre para cada clase no-prioritaria, todo conflicto es resoluble.*

**Demostración:**
1. Sea $c_p$ clase prioritaria y $c_n$ clase no-prioritaria en conflicto
2. $c_n$ tiene al menos un salón válido $s_{alt}$ libre (por hipótesis)
3. Mover $c_n$ a $s_{alt}$ libera el salón preferido de $c_p$
4. Asignar $c_p$ a su salón preferido
5. Conflicto resuelto


---

## 6. Marcado de Índices Inmutables


---

### 6.1 Estructura de Datos

\begin{verbatim}
{
  "indices": [12, 45, 67, 89, ...],
  "total": 88,
  "timestamp": "2025-12-21T11:00:00",
  "version": "1.0"
}
\end{verbatim}



---

### 6.2 Implementación

\begin{lstlisting}[language=Python]
def marcar_inmutables(self, clases_p1):
    """
    Marca índices de clases P1 como inmutables
    """
    self.indices_inmutables = [c['idx'] for c in clases_p1]
    
    # Guardar en JSON
    data = {
        'indices': self.indices_inmutables,
        'total': len(self.indices_inmutables),
        'timestamp': datetime.now().isoformat(),
        'version': '1.0'
    }
    
    with open('datos_estructurados/indices_inmutables_p1.json', 'w') as f:
        json.dump(data, f, indent=2)
\end{verbatim}



---

### 6.3 Invariante de Inmutabilidad

**Definición:**

\end{equation}

\forall i \in I_{inmutables}, \forall t: s_t(c_i) = pref(c_i)

\begin{equation}

Donde:
- $I_{inmutables}$ = conjunto de índices inmutables
- $s_t$ = solución en tiempo $t$
- $pref(c_i)$ = salón preferido de clase $c_i$

---

**Verificación:**

\begin{lstlisting}[language=Python]
def verificar_invariante(self, df):
    """
    Verifica que todas las clases inmutables están en su salón preferido
    """
    violaciones = 0
    
    for idx in self.indices_inmutables:
        clase = df.iloc[idx]
        salon_actual = clase['Salon']
        salon_esperado = self.obtener_salon_preferido(clase)
        
        if salon_actual != salon_esperado:
            print(f"❌ Violación en índice {idx}: {salon_actual} != {salon_esperado}")
            violaciones += 1
    
    return violaciones == 0
\end{verbatim}



---

## 7. Salidas del Sistema


---

### 7.1 Horario Pre-asignado

**Archivo:** `datos_estructurados/00_Horario_PreAsignado_P1.csv`

**Formato:**
\begin{verbatim}
Dia,Bloque_Horario,Materia,Grupo,Profesor,Salon,Es_Invalido,Tipo_Salon,Piso
Lunes,0700,LENGUAJES Y AUTÓMATAS I,2527A,PROFESOR 3,FFA,0,Teoría,1
...
\end{verbatim}


**Características:**
- Todas las clases P1 en su salón preferido
- Clases no-P1 pueden estar en salones subóptimos
- Listo para ser optimizado por algoritmos


---

### 7.2 Índices Inmutables

**Archivo:** `datos_estructurados/indices_inmutables_p1.json`

**Uso:**
\begin{lstlisting}[language=Python]
# En optimizadores
with open('datos_estructurados/indices_inmutables_p1.json') as f:
    data = json.load(f)
    indices_inmutables = set(data['indices'])

# Durante optimización
if idx in indices_inmutables:
    continue  # No modificar esta clase
\end{verbatim}



---

## 8. Métricas y Validación


---

### 8.1 Métricas de Pre-asignación

\begin{verbatim}
Ejecución de pre_asignar_p1.py:
├── Clases prioritarias identificadas: 88
├── Clases asignadas exitosamente: 88
├── Conflictos resueltos: 12
├── Cumplimiento: 100%
└── Tiempo: 0.3s
\end{verbatim}



---

### 8.2 Validación de Salida

\begin{lstlisting}[language=Python]
def validar_salida(self, df):
    """
    Valida que la salida sea correcta
    """
    checks = {
        'total_clases': len(df) == 680,
        'p1_cumplimiento': self.verificar_p1(df) == 100.0,
        'sin_invalidos_p1': self.verificar_sin_invalidos_p1(df),
        'sin_conflictos_p1': self.verificar_sin_conflictos_p1(df)
    }
    
    return all(checks.values()), checks
\end{verbatim}



---

## 9. Complejidad Computacional


---

### 9.1 Análisis Temporal

**Identificación:** $O(n)$ donde $n$ = número de clases

**Ordenamiento:** $O(p \log p)$ donde $p$ = clases prioritarias

**Asignación:** $O(p \cdot  m)$ donde $m$ = salones promedio por conflicto

**Total:** $O(n + p \log p + p \cdot  m) = O(n)$ ya que $p \ll n$

**Tiempo Real:** ~0.3 segundos para 680 clases


---

### 9.2 Análisis Espacial

\end{equation}

S = O(n + p) = O(n)

\begin{equation}

**Memoria:** ~5 MB


---

## 10. Ventajas del Enfoque

✅ **Garantía matemática** de 100% P1  
✅ **Separación clara** entre restricciones  
✅ **Simplifica optimización** posterior  
✅ **Rápido** (<1 segundo)  
✅ **Robusto** ante cambios en preferencias  
✅ **Verificable** mediante invariante  


---

## 11. Casos Especiales


---

### 11.1 Conflictos Irresolvables

Si dos clases P1 quieren el mismo salón al mismo tiempo:

**Solución Manual:**
1. Identificar el conflicto
2. Contactar a profesores involucrados
3. Negociar cambio de horario o salón
4. Actualizar preferencias


---

### 11.2 Salones Insuficientes

Si no hay suficientes salones del tipo requerido:

**Solución:**
1. Identificar clases afectadas
2. Evaluar posibilidad de usar salones alternativos
3. Ajustar configuración de materias
4. Re-ejecutar pre-asignación


---

## 12. Integración con Optimizadores


---

### 12.1 Carga en Optimizadores

\begin{lstlisting}[language=Python]
# En optimizador_greedy.py, optimizador_ml.py, optimizador_genetico.py
def __init__(self):
    # Cargar horario pre-asignado
    self.df_inicial = pd.read_csv(
        'datos_estructurados/00_Horario_PreAsignado_P1.csv'
    )
    
    # Cargar índices inmutables
    with open('datos_estructurados/indices_inmutables_p1.json') as f:
        data = json.load(f)
        self.indices_inmutables = set(data['indices'])
\end{verbatim}



---

### 12.2 Protección Durante Optimización

\begin{lstlisting}[language=Python]
# En operadores de optimización
def aplicar_operador(self, solucion):
    for idx in range(len(solucion)):
        # Proteger inmutables
        if idx in self.indices_inmutables:
            continue
        
        # Aplicar modificación solo a no-inmutables
        solucion[idx] = nueva_asignacion(idx)
\end{verbatim}



---

## 13. Conclusiones

El sistema de pre-asignación:

1. **Garantiza** 100% cumplimiento de PRIORIDAD 1
2. **Simplifica** el problema de optimización
3. **Separa** restricciones duras de suaves
4. **Permite** que optimizadores se enfoquen en P2 y P3
5. **Proporciona** base sólida para todo el sistema

Es un componente **crítico** que hace posible el enfoque de prioridades jerárquicas.


---

## Referencias

1. Apt, K. R. (2003). *Principles of Constraint Programming*. Cambridge University Press.

2. Rossi, F., Van Beek, P., & Walsh, T. (2006). *Handbook of Constraint Programming*. Elsevier.

3. Dechter, R. (2003). *Constraint Processing*. Morgan Kaufmann.







---

<!-- _class: lead blue -->
# Algoritmo Greedy + Hill Climbing

**Construcción Voraz y Refinamiento Local**

---

## 1. Introducción y Fundamento Teórico


---

### 1.1 Paradigma de Algoritmos Voraces

Un **algoritmo voraz (greedy)** construye una solución tomando decisiones localmente óptimas en cada paso, con la esperanza de que estas decisiones conduzcan a un óptimo global.

---

**Definición Formal:**

Sea $S = \{s_1, s_2, ..., s_n\}$ el conjunto de decisiones a tomar. Un algoritmo voraz:

1. Inicializa $Sol = \emptyset$
2. Para cada decisión $s_i$ en orden:
   - Selecciona $s_i^* = \arg\min_{s \in candidatos} costo(s)$
   - Actualiza $Sol = Sol \cup \{s_i^*\}$
3. Retorna $Sol$

---

**Teorema del Algoritmo Voraz:**
*Un algoritmo voraz produce una solución óptima si el problema exhibe la propiedad de elección voraz y subestructura óptima.*

**Aplicación a Asignación de Salones:**
- **Elección voraz:** Asignar cada clase al salón de menor costo incremental
- **Subestructura:** La solución óptima contiene soluciones óptimas a subproblemas
- **Limitación:** Nuestro problema NO garantiza optimalidad global con greedy puro
---

### 1.2 Búsqueda Local: Hill Climbing

**Hill Climbing** es un algoritmo de búsqueda local que mejora iterativamente una solución explorando su vecindario.

**Definición Formal:**

\end{equation}

\begin{align}

&\text{Inicializar: } s = s_0 \\

&\text{Repetir:} \\

&\quad s' = \arg\min_{s'' \in N(s)} f(s'') \\

&\quad \text{Si } f(s') < f(s): \\

&\quad\quad s = s' \\

&\quad \text{Sino: } \\

&\quad\quad \text{retornar } s

\end{align}

\begin{equation}
---
Donde:
- $s$ = solución actual
- $N(s)$ = vecindario de $s$
- $f(s)$ = función objetivo (energía)

**Propiedades:**
- **Convergencia:** Garantizada a un óptimo local
- **Complejidad:** $O(k \cdot  |N(s)|)$ donde $k$ = iteraciones
- **Limitación:** Puede quedar atrapado en óptimos locales
---

### 1.3 Estrategia Híbrida

Combinamos ambos enfoques:

\begin{verbatim}
Solución = Greedy(problema)  # Construcción rápida
Solución = HillClimbing(Solución)  # Refinamiento local
\end{verbatim}


**Ventajas del Híbrido:**
- Greedy proporciona punto de partida de calidad
- Hill Climbing escapa de decisiones voraces subóptimas
- Balance entre velocidad y calidad


---

## 2. Fase 1: Construcción Voraz


---

### 2.1 Ordenamiento de Clases

**Objetivo:** Procesar clases en orden que maximice probabilidad de buenas asignaciones.

**Criterio de Ordenamiento Multi-nivel:**

Para dos clases $c_i$ y $c_j$, definimos $c_i \prec c_j$ si:

\end{equation}

\begin{cases}

prioridad(c_i) > prioridad(c_j) & \text{(1° criterio)} \\

\text{o si } prioridad(c_i) = prioridad(c_j): \\

\quad num\_restricciones(c_i) > num\_restricciones(c_j) & \text{(2° criterio)} \\

\text{o si } num\_restricciones(c_i) = num\_restricciones(c_j): \\

\end{cases}

\begin{equation}
---

**Justificación Teórica:**

**Lema 1 (Ordenamiento Óptimo):**
*Procesar clases con más restricciones primero minimiza la probabilidad de infactibilidad.*

**Demostración:**
Sea $R_i$ el conjunto de restricciones de $c_i$ y $S_i$ el conjunto de salones válidos.

- Si $|S_i| < |S_j|$, entonces $c_i$ tiene menos opciones
- Procesar $c_i$ primero garantiza que al menos una opción esté disponible
- Procesar $c_j$ después aún deja $|S_j| - 1 \geq |S_i|$ opciones
- Por lo tanto, el orden minimiza conflictos
---

**Implementación:**

\begin{lstlisting}[language=Python]
def ordenar_clases(self, df):
    """
    Ordena clases por criterio multi-nivel
    """
    df_ordenado = df.copy()
    
    # Calcular métricas para cada clase
    df_ordenado['prioridad_num'] = df_ordenado.apply(
        lambda row: self.obtener_prioridad(row), axis=1
    )
    
    df_ordenado['num_restricciones'] = df_ordenado.apply(
        lambda row: self.contar_restricciones(row), axis=1
    )
    
    df_ordenado['num_salones_validos'] = df_ordenado.apply(
        lambda row: len(self.obtener_salones_validos(row)), axis=1
    )
    
    # Ordenar por criterios
    df_ordenado = df_ordenado.sort_values(
        by=['prioridad_num', 'num_restricciones', 'num_salones_validos'],
        ascending=[False, False, True]
    )
    
    return df_ordenado
\end{verbatim}



---

### 2.2 Función de Score Voraz

Para cada par $(clase, salon)$, calculamos un score que estima el costo incremental:

\end{equation}

score(c_i, s_j) = \sum_{k=1}^{m} w_k \cdot componente_k(c_i, s_j)

\begin{equation}

**Componentes del Score:**

#### A. Distancia al Último Salón del Profesor

\end{equation}

componente_1(c_i, s_j) = \begin{cases}

0 & \text{si } c_i \text{ es primera clase del profesor} \\

d(ultimo\_salon(profesor(c_i)), s_j) & \text{en otro caso}

\end{cases}

\begin{equation}
---

Donde $d(\cdot , \cdot )$ es la función de distancia definida anteriormente:

\end{equation}

d(s_a, s_b) = \begin{cases}

0 & \text{si } s_a = s_b \\

1 & \text{si } piso(s_a) = piso(s_b) \land s_a \neq s_b \\

10 & \text{si } piso(s_a) \neq piso(s_b)

\end{cases}

\begin{equation}
---

#### B. Ocupación del Salón

\end{equation}

\begin{equation}

**Interpretación:** Preferir salones menos utilizados para balancear carga.

---

#### C. Penalización por Tipo Incorrecto

\end{equation}

componente_3(c_i, s_j) = \begin{cases}

\infty & \text{si } tipo(c_i) \neq tipo(s_j) \\

0 & \text{en otro caso}

\end{cases}

\begin{equation}

---

#### D. Bonus por Preferencia

\end{equation}

componente_4(c_i, s_j) = \begin{cases}

-B & \text{si } s_j = salon\_preferido(c_i) \\

0 & \text{en otro caso}

\end{cases}

\begin{equation}

**Donde $B > 0$ es un bonus que incentiva cumplir preferencias.**

---

#### E. Penalización por Salón Inválido

\end{equation}

componente_5(c_i, s_j) = \begin{cases}

\infty & \text{si } s_j \in S_{invalidos} \\

0 & \text{en otro caso}

\end{cases}

\begin{equation}

**Score Total:**

\end{equation}

score(c_i, s_j) = 10 \cdot componente_1 + 5 \cdot componente_2 + componente_3 - 50 \cdot componente_4 + componente_5

\begin{equation}
---

**Pesos Justificados:**
- Distancia (10): Impacto directo en movilidad
- Ocupación (5): Balanceo de recursos
- Tipo ($\infty$): Restricción dura
- Preferencia (-50): Incentivo fuerte
- Inválido ($\infty$): Restricción dura


---

### 2.3 Algoritmo de Construcción

\begin{lstlisting}[language=Python]
def construccion_greedy(self, df):
    """
    Construye solución inicial mediante selección voraz
    """
    solucion = {}
    ocupacion = {}  # (dia, bloque, salon) -> idx_clase
    
    # Ordenar clases
    df_ordenado = self.ordenar_clases(df)
    
    for idx, clase in df_ordenado.iterrows():
        # Obtener salones candidatos
        candidatos = self.obtener_salones_validos(clase)
        
        # Filtrar salones ocupados en este horario
        candidatos_libres = [
            s for s in candidatos
            if (clase['Dia'], clase['Bloque_Horario'], s) not in ocupacion
        ]
        
        if not candidatos_libres:
            # Caso excepcional: forzar asignación
            mejor_salon = candidatos[0]
        else:
            # Selección voraz: minimizar score
            mejor_salon = None
            mejor_score = float('inf')
            
            for salon in candidatos_libres:
                score = self.calcular_score(clase, salon, solucion)
                
                if score < mejor_score:
                    mejor_score = score
                    mejor_salon = salon
        
        # Asignar
        solucion[idx] = mejor_salon
        ocupacion[(clase['Dia'], clase['Bloque_Horario'], mejor_salon)] = idx
    
    return solucion
\end{verbatim}



---

### 2.4 Análisis de Complejidad - Fase 1

**Ordenamiento:**
\end{equation}

T_{sort} = O(n \log n)

\begin{equation}

**Cálculo de salones válidos por clase:**
\end{equation}

T_{valid} = O(n \cdot m)

\begin{equation}

Donde $m$ = número promedio de salones candidatos.

**Selección voraz:**
\end{equation}

T_{greedy} = O(n \cdot m)

\begin{equation}

**Total Fase 1:**
\end{equation}

T_1 = O(n \log n + n \cdot m) = O(n \cdot m)

\begin{equation}

Para $n = 680$, $m \approx 5$:
\end{equation}

T_1 \approx 3400 \text{ operaciones}

\begin{equation}


---

## 3. Fase 2: Hill Climbing


---

### 3.1 Definición del Vecindario

El vecindario $N(s)$ de una solución $s$ se define como:

\end{equation}

N(s) = \{s' : s' \text{ difiere de } s \text{ en exactamente un intercambio válido}\}

\begin{equation}

**Intercambio Válido:**

Para clases $c_i$ y $c_j$:

\end{equation}

intercambio\_valido(c_i, c_j) \Leftrightarrow \begin{cases}

tipo(c_i) = tipo(c_j) & \text{(mismo tipo)} \\

i \notin I_{inmutables} \land j \notin I_{inmutables} & \text{(no protegidas)} \\

\neg conflicto(c_i, s(c_j)) \land \neg conflicto(c_j, s(c_i)) & \text{(sin conflictos)}

\end{cases}

\begin{equation}

Donde $I_{inmutables}$ es el conjunto de índices de clases PRIORIDAD 1.

---

**Tamaño del Vecindario:**

\end{equation}

\begin{equation}

Para $n_{modificables} \approx 600$:
\end{equation}

\begin{equation}
---

### 3.2 Función de Energía

La función de energía $E(s)$ cuantifica la calidad de una solución:

\end{equation}

E(s) = E_{movimientos}(s) + E_{pisos}(s) + E_{distancia}(s) + E_{penalizaciones}(s)

\begin{equation}

#### A. Energía por Movimientos

\end{equation}

\begin{equation}

Donde:
\end{equation}

salones\_usados(p, s) = \{s(c) : c \in C_p\}

\begin{equation}
---

**Ejemplo Numérico:**
- Profesor tiene 5 clases en salones: {FF1, FF2, FF1, FF3, FF2}
- $salones_usados = \{FF1, FF2, FF3\}$
- $|salones_usados| = 3$
- $E_{movimientos} = 10 \cdot  (3 - 1) = 20$
---

#### B. Energía por Cambios de Piso

Sea $C_p^{sorted} = [c_1, c_2, ..., c_k]$ las clases del profesor $p$ ordenadas cronológicamente:

\end{equation}

\begin{equation}
---

#### C. Energía por Distancia

\end{equation}

\begin{equation}
---

#### D. Penalizaciones

\end{equation}

\begin{align}

\end{align}

\begin{equation}

**Pesos Utilizados:**
- $w_m = 10$ (movimientos)
- $w_p = 5$ (cambios de piso)
- $w_d = 3$ (distancia)
---

### 3.3 Estrategia de Búsqueda

Implementamos **Steepest Descent Hill Climbing**:

\begin{lstlisting}[language=Python]
def hill_climbing(self, solucion, df):
    """
    Mejora solución mediante búsqueda local
    """
    mejor_solucion = solucion.copy()
    mejor_energia = self.calcular_energia(mejor_solucion, df)
    
    for iteracion in range(self.max_iter_hc):
        mejoro = False
        indices = list(solucion.keys())
        
        # Generar vecinos aleatorios
        for _ in range(self.intentos_por_iter):
            # Seleccionar par aleatorio
            idx1, idx2 = random.sample(indices, 2)
            
            # Verificar que no sean inmutables
            if idx1 in self.indices_inmutables or idx2 in self.indices_inmutables:
                continue
\end{verbatim}

---
\begin{lstlisting}[language=Python]
            # Verificar mismo tipo
            if self.tipos_por_idx[idx1] != self.tipos_por_idx[idx2]:
                continue

            # Crear vecino
            vecino = mejor_solucion.copy()
            vecino[idx1], vecino[idx2] = vecino[idx2], vecino[idx1]
            
            # Evaluar
            energia_vecino = self.calcular_energia(vecino, df)
            
            # Aceptar si mejora
            if energia_vecino < mejor_energia:
                mejor_solucion = vecino
                mejor_energia = energia_vecino
                mejoro = True
                break  # Steepest descent: aceptar primera mejora
        
        # Criterio de parada
        if not mejoro:
            print(f"   Convergió en iteración {iteracion}")
            break
    
    return mejor_solucion
\end{verbatim}

---

### 3.4 Criterio de Parada

El algoritmo se detiene cuando:

\end{equation}

\forall s' \in N_{muestreado}(s): E(s') \geq E(s)

\begin{equation}

Donde $N_{muestreado}$ es un subconjunto aleatorio de $N(s)$.

---

**Teorema 2 (Convergencia de Hill Climbing):**
*El algoritmo de Hill Climbing converge a un óptimo local en tiempo finito.*

**Demostración:**
1. El espacio de soluciones es finito: $|S| = m^n$
2. La energía es discreta y acotada inferiormente: $E(s) \geq 0$
3. Cada iteración reduce estrictamente $E$ o termina
4. No puede haber ciclos (energía siempre decrece)
5. Por lo tanto, converge en $\leq |S|$ iteraciones
---

**Cota Superior Práctica:**

Con muestreo aleatorio de $k$ vecinos por iteración:

\end{equation}

T_{convergencia} \leq \frac{E_{inicial} - E_{optimo\_local}}{mejora\_promedio} \cdot k

\begin{equation}

Empíricamente: $\approx 20-50$ iteraciones

---

### 3.5 Análisis de Complejidad - Fase 2

**Por iteración:**
\end{equation}

T_{iter} = k \cdot (T_{generar} + T_{evaluar})

\begin{equation}

Donde:
- $k$ = intentos por iteración (50)
- $T_{generar} = O(1)$ (intercambio simple)
- $T_{evaluar} = O(n)$ (recalcular energía)

\end{equation}

T_{iter} = O(k \cdot n)

\begin{equation}

**Total Fase 2:**
\end{equation}

T_2 = O(max\_iter \cdot k \cdot n)

\begin{equation}

Para $max_iter = 100$, $k = 50$, $n = 680$:
\end{equation}

T_2 \approx 3.4 \times 10^6 \text{ operaciones}

\begin{equation}


---

## 📖 Ejemplo Didáctico: Greedy + Hill Climbing

> [!NOTE]
> **Este es un ejemplo didáctico simplificado** con un problema pequeño (10 clases, 4 salones, 4 profesores) para demostrar el funcionamiento del algoritmo de manera clara y comprensible.

**Problema Simplificado para Demostración**

---

### Problema Didáctico

**Datos:**
- 10 clases
- 4 salones
- 4 profesores
- Objetivo: Minimizar movimientos de profesores

**Archivo:** `ejemplos_didacticos/01_greedy_hill_climbing.py`

**Propósito:** Demostrar cómo funciona el algoritmo en un caso pequeño

---

### Algoritmo del Ejemplo

**Fase 1: Construcción Greedy (Subóptima Intencional)**
\begin{lstlisting}[language=Python]
def greedy_construccion():
    asignacion = {}
    for i, clase in enumerate(clases):
        salones_compatibles = [s for s in salones 
                              if s.tipo == clase.tipo]
        # Rotar entre salones (genera solución subóptima)
        salon = salones_compatibles[i % len(salones_compatibles)]
        asignacion[clase.id] = salon.id
    return asignacion
\end{verbatim}


**Fase 2: Hill Climbing (Mejora Local)**
\begin{lstlisting}[language=Python]
def hill_climbing(asignacion_inicial):
    actual = asignacion_inicial
    while True:
        vecinos = generar_vecinos(actual)
        mejor = min(vecinos, key=energia)
        if energia(mejor) >= energia(actual):
            break
        actual = mejor
    return actual
\end{verbatim}


---

### Resultados del Ejemplo

**Solución Inicial (Greedy):**
\begin{verbatim}
Energía: 40
Movimientos: 8
\end{verbatim}


**Solución Final (Hill Climbing):**
\begin{verbatim}
Energía: 20
Movimientos: 4
\end{verbatim}


**Mejora Lograda:** 50% ✅

**Conclusión:** El Hill Climbing mejora significativamente la solución greedy inicial

---

### Ejecutar el Ejemplo

\begin{lstlisting}[language=Bash]
cd ejemplos_didacticos
python3 01_greedy_hill_climbing.py
\end{verbatim}


**Salida esperada:**
\begin{verbatim}
============================================================
EJEMPLO: Greedy + Hill Climbing
============================================================

Solución Inicial (Greedy):
  Energía: 40

Hill Climbing:
  Iteración 0: Energía = 40
  Iteración 5: Energía = 30
  Iteración 10: Energía = 20

Solución Final:
  Energía: 20

Mejora: 50.0%
✅ Completado!
\end{verbatim}

---

## 4. Protección de PRIORIDAD 1


---

### 4.1 Mecanismo de Inmutabilidad

**Definición:**

Sea $I \subset \{0, 1, ..., n-1\}$ el conjunto de índices de clases PRIORIDAD 1:

\end{equation}

I = \{i : c_i \in P_1\}

\begin{equation}

**Invariante:**

\end{equation}

\forall i \in I, \forall t: s_t(c_i) = pref(c_i)

\begin{equation}

Donde $s_t$ es la solución en el tiempo $t$.

---

**Implementación:**

\begin{lstlisting}[language=Python]
def cargar_indices_inmutables(self):
    """
    Carga índices de clases que no deben modificarse
    """
    try:
        with open('datos_estructurados/indices_inmutables_p1.json', 'r') as f:
            data = json.load(f)
            self.indices_inmutables = set(data.get('indices', []))
            print(f"✅ Índices inmutables cargados: {len(self.indices_inmutables)} clases")
    except FileNotFoundError:
        self.indices_inmutables = set()
        print("⚠️  No se encontró archivo de índices inmutables")
\end{verbatim}

---
### 4.2 Verificación en Hill Climbing

\begin{lstlisting}[language=Python]
# En cada intercambio propuesto
if idx1 in self.indices_inmutables or idx2 in self.indices_inmutables:
    continue  # Saltar este intercambio
\end{verbatim}

---

**Teorema 3 (Preservación de P1):**
*Si $s_0$ satisface PRIORIDAD 1 y se respetan los índices inmutables, entonces $s_t$ satisface PRIORIDAD 1 para todo $t$.*

**Demostración:**
Por inducción en $t$:
- **Caso base ($t=0$):** $s_0$ satisface P1 por construcción (pre-asignación)
- **Paso inductivo:** Asumimos $s_t$ satisface P1
  - Hill Climbing solo modifica índices $\notin I$
  - Por lo tanto, $\forall i \in I: s_{t+1}(c_i) = s_t(c_i) = pref(c_i)$
  - Luego $s_{t+1}$ satisface P1
---

### 4.3 Corrección Post-Optimización

Como medida de seguridad adicional:

\begin{lstlisting}[language=Python]
def corregir_prioridades_final(self, solucion, df):
    """
    Garantiza 100% cumplimiento de PRIORIDAD 1
    """
    correcciones = 0
    
    for idx in self.indices_inmutables:
        clase = df.iloc[idx]
        salon_esperado = self.obtener_salon_preferido(clase)
        
        if solucion[idx] != salon_esperado:
            solucion[idx] = salon_esperado
            correcciones += 1
    
    if correcciones > 0:
        print(f"   ✅ {correcciones} clases corregidas a salón prioritario")
    
    return solucion
\end{verbatim}



---

## 5. Optimizaciones y Mejoras


---

### 5.1 Caching de Energía

**Problema:** Recalcular energía completa es $O(n)$

---

**Solución:** Calcular solo cambio incremental

\begin{lstlisting}[language=Python]
def calcular_delta_energia(self, solucion, idx1, idx2, df):
    """
    Calcula cambio de energía por intercambio
    """
    # Solo recalcular para profesores afectados
    prof1 = df.iloc[idx1]['Profesor']
    prof2 = df.iloc[idx2]['Profesor']
    
    delta = 0
    delta += self.delta_movimientos(prof1, idx1, idx2, solucion, df)
    
    if prof1 != prof2:
        delta += self.delta_movimientos(prof2, idx1, idx2, solucion, df)
    
    return delta
\end{verbatim}


**Complejidad:** $O(|C_{prof}|)$ en lugar de $O(n)$

---

### 5.2 Tabu Search

Evitar ciclos manteniendo lista de movimientos prohibidos:

\begin{lstlisting}[language=Python]
def hill_climbing_tabu(self, solucion, df):
    tabu_list = deque(maxlen=10)  # Últimos 10 movimientos
    
    for iteracion in range(self.max_iter_hc):
        # ... generar vecino ...
        
        movimiento = (idx1, idx2)
        if movimiento in tabu_list:
            continue  # Prohibido
        
        # ... evaluar y aceptar ...
        
        if mejoro:
            tabu_list.append(movimiento)
\end{verbatim}

---

### 5.3 Simulated Annealing

Aceptar ocasionalmente movimientos que empeoran:

\end{equation}

P(aceptar) = \begin{cases}

1 & \text{si } \Delta E < 0 \\

e^{-\Delta E / T} & \text{si } \Delta E \geq 0

\end{cases}

\begin{equation}

Donde $T$ es la "temperatura" que decrece con el tiempo:

\end{equation}

T_t = T_0 \cdot \alpha^t, \quad 0 < \alpha < 1

\begin{equation}


---

## 6. Resultados y Análisis


---

### 6.1 Métricas de Rendimiento

\begin{verbatim}
FASE 1: Construcción Greedy
├── Tiempo: 2.3s
├── Energía inicial: 82,948
└── Clases asignadas: 680/680

FASE 2: Hill Climbing
├── Tiempo: 27.0s
├── Iteraciones: 54
├── Energía final: 82,716
├── Mejora: 232 (-0.28%)
└── Convergencia: Sí

RESULTADOS FINALES:
├── Movimientos: 314 (-12.0% vs inicial)
├── Cambios piso: 206 (-28.2% vs inicial)
├── Distancia: 1,951 (-31.5% vs inicial)
└── PRIORIDAD 1: 100% (98/98)
\end{verbatim}



---

### 6.2 Comparación con Otros Algoritmos

| Métrica | Greedy+HC | ML | Genético | Óptimo Teórico |
|---------|-----------|-----|----------|----------------|
| Tiempo | 29.3s | 15.8s | 73.9s | ∞ |
| Movimientos | 314 | 365 | 378 | ≥280 |
| Cambios piso | 206 | 223 | 286 | ≥180 |
| Distancia | 1,951 | 1,821 | 2,413 | ≥1,500 |
| Gap vs óptimo | ~5% | ~3% | ~10% | 0% |

**Nota:** Óptimo teórico es estimación basada en cotas inferiores.


---

### 6.3 Análisis de Sensibilidad

**Variación de Parámetros:**

| max_iter | Tiempo | Energía Final | Mejora |
|----------|--------|---------------|--------|
| 10 | 5.2s | 82,850 | Baja |
| 50 | 15.1s | 82,730 | Media |
| 100 | 29.3s | 82,716 | Alta |
| 200 | 58.7s | 82,714 | Marginal |

**Conclusión:** 100 iteraciones es el punto óptimo (rendimientos decrecientes después).


---

## 7. Ventajas y Limitaciones


---

### 7.1 Ventajas

✅ **Velocidad:** Construcción inicial muy rápida ($O(n \log n)$)  
✅ **Simplicidad:** Fácil de entender e implementar  
✅ **Determinismo:** Resultados reproducibles (con semilla fija)  
✅ **Escalabilidad:** Complejidad lineal en $n$  
✅ **Robustez:** Siempre encuentra solución factible  
✅ **Balance:** Buena relación calidad/tiempo 

--- 

### 7.2 Limitaciones

❌ **Óptimos Locales:** No garantiza óptimo global  
❌ **Sensibilidad al Orden:** Orden inicial afecta resultado  
❌ **Exploración Limitada:** Vecindario pequeño  
❌ **Plateau:** Puede estancarse en mesetas  
❌ **Parámetros:** Requiere ajuste de pesos 

---

### 7.3 Cuándo Usar Greedy+HC

**Recomendado cuando:**
- Se requiere solución rápida (< 1 minuto)
- Calidad media-alta es suficiente
- Problema es relativamente estructurado
- Hay buena heurística de ordenamiento

**No recomendado cuando:**
- Se requiere óptimo global garantizado
- Hay tiempo ilimitado para búsqueda
- Problema es altamente irregular
- Espacio de búsqueda es muy pequeño


---

## 8. Extensiones Futuras


---

### 8.1 Variable Neighborhood Search (VNS)

Usar múltiples definiciones de vecindario:

\end{equation}

N_1(s) = \text{intercambios simples}

\begin{equation}
\end{equation}

N_2(s) = \text{intercambios dobles}

\begin{equation}
\end{equation}

N_3(s) = \text{rotaciones ciclicas}

\begin{equation}
---

### 8.2 Iterated Local Search (ILS)

Aplicar perturbaciones para escapar óptimos locales:

\begin{lstlisting}[language=Python]
s = ConstructionGreedy()
s = HillClimbing(s)

for i in range(max_restarts):
    s' = Perturb(s)
    s' = HillClimbing(s')
    if E(s') < E(s):
        s = s'

return s
\end{verbatim}

---

### 8.3 Guided Local Search

Penalizar características de óptimos locales visitados:

\end{equation}

E_{guided}(s) = E(s) + \lambda \sum_{i} p_i \cdot feature_i(s)

\begin{equation}

Donde $p_i$ aumenta cada vez que se visita un óptimo local con $feature_i$.


---

## 9. Conclusiones

El algoritmo Greedy + Hill Climbing ofrece un **excelente balance** entre:
- **Velocidad de ejecución** (~30s)
- **Calidad de solución** (top 2 de 4 algoritmos)
- **Simplicidad de implementación**
- **Robustez y confiabilidad**

Es la **opción recomendada** para uso general en el sistema de asignación de salones.


---

## Referencias

1. Cormen, T. H., et al. (2009). *Introduction to Algorithms* (3rd ed.). MIT Press.

2. Russell, S., & Norvig, P. (2020). *Artificial Intelligence: A Modern Approach* (4th ed.). Pearson.

3. Papadimitriou, C. H., & Steiglitz, K. (1998). *Combinatorial Optimization: Algorithms and Complexity*. Dover.

4. Aarts, E., & Lenstra, J. K. (2003). *Local Search in Combinatorial Optimization*. Princeton University Press.

5. Hoos, H. H., & Stützle, T. (2004). *Stochastic Local Search: Foundations and Applications*. Morgan Kaufmann.


---

<!-- _class: lead blue -->
# Algoritmo de Machine Learning

**Aprendizaje Supervisado para Asignación Óptima**

---

## 1. Introducción y Fundamento Teórico

El enfoque de Machine Learning para la asignación de salones se basa en la **hipótesis de aprendizaje supervisado**: si existe un patrón en las asignaciones óptimas previas, un modelo puede aprender a predecir asignaciones de calidad similar sin necesidad de búsqueda exhaustiva.

**Ventaja Principal:** Una vez entrenado, el modelo puede generar soluciones en tiempo casi constante $O(n \cdot  d)$ donde $d$ es la profundidad del árbol, comparado con $O(n^2)$ o peor de otros métodos.

---

### 1.2 Arquitectura del Sistema ML

El sistema utiliza un **ensemble de dos modelos complementarios**:

1. **Random Forest Classifier:** Predice el salón óptimo para cada clase
2. **Gradient Boosting Regressor:** Estima la calidad de cada asignación

\begin{verbatim}
Entrada: Características de la clase
   ↓
[Random Forest] → Predicción de salón
   ↓
[Gradient Boosting] → Score de calidad
   ↓
Selección del mejor salón válido
   ↓
Salida: Asignación optimizada
\end{verbatim}



---

## 2. Extracción de Características (Feature Engineering)

### 2.1 Vector de Características

Para cada clase $c_i$, construimos un vector de características $\mathbf{x}_i \in \mathbb{R}^d$:

\end{equation}

\mathbf{x}_i = [x_1, x_2, ..., x_d]^T

\begin{equation}

**Categorías de Características:**


---

#### A. Características Temporales

\end{equation}

\begin{align}

x_1 &= dia\_semana(c_i) &&\in \{0, 1, 2, 3, 4\} &&\text{(Lun-Vie)} \\

x_2 &= bloque\_horario(c_i) &&\in \{7, 8, ..., 21\} &&\text{(Hora inicio)} \\

x_3 &= es\_manana(c_i) &&\in \{0, 1\} &&\text{(Antes de 12:00)} \\

x_4 &= es\_tarde(c_i) &&\in \{0, 1\} &&\text{(12:00-18:00)} \\

x_5 &= es\_noche(c_i) &&\in \{0, 1\} &&\text{(Despues 18:00)}

\end{align}

\begin{equation}


---

#### B. Características de Materia

\end{equation}

\begin{align}

x_7 &= tipo\_clase(c_i) &&\in \{0, 1\} &&\text{(0=Teoria, 1=Lab)} \\

x_8 &= num\_estudiantes(c_i) &&\in \mathbb{N} &&\text{(Tamano grupo)}

\end{align}

\begin{equation}


---

#### C. Características de Profesor

\end{equation}

\begin{align}

x_{10} &= tiene\_preferencia(c_i) &&\in \{0, 1\} \\

x_{11} &= prioridad\_preferencia(c_i) &&\in \{0, 1, 2, 3\}

\end{align}

\begin{equation}


---

#### D. Características Contextuales

\end{equation}

\begin{align}

x_{12} &= num\_clases\_profesor\_dia(c_i) &&\in \mathbb{N} \\

x_{13} &= posicion\_en\_dia(c_i) &&\in \{1, 2, ..., k\} \\

x_{15} &= piso\_clase\_anterior(c_i) &&\in \{0, 1\}

\end{align}

\begin{equation}


---

#### E. Características de Grupo

\end{equation}

\begin{align}

x_{16} &= semestre(c_i) &&\in \{1, 2, ..., 9\} \\

x_{17} &= es\_primer\_semestre(c_i) &&\in \{0, 1\} \\

\end{align}

\begin{equation}

**Total de características:** $d \approx 18-25$ (dependiendo de encoding)

---

### 2.2 Encoding de Variables Categóricas

**Label Encoding:**
Para variables ordinales (materia, profesor, grupo):

\end{equation}

encode(categoria) = \{categoria_1 \rightarrow 0, categoria_2 \rightarrow 1, ..., categoria_n \rightarrow n-1\}

\begin{equation}
---

**One-Hot Encoding (alternativa):**
Para $k$ categorías, crear $k$ variables binarias:

\end{equation}

\mathbf{x}_{one-hot} = [0, 0, ..., 1, ..., 0]^T

\begin{equation}

Donde el 1 está en la posición correspondiente a la categoría.

**Decisión de Diseño:** Usamos Label Encoding para reducir dimensionalidad, ya que Random Forest maneja bien variables categóricas ordinales.

---

### 2.3 Normalización

Para características numéricas continuas:

\end{equation}

x'_i = \frac{x_i - \mu_i}{\sigma_i}

\begin{equation}

Donde:
- $\mu_i$ = media de la característica $i$
- $\sigma_i$ = desviación estándar de la característica $i$

**Nota:** Random Forest no requiere normalización estricta, pero mejora la interpretabilidad.


---

## 3. Modelo 1: Random Forest Classifier

### 3.1 Fundamento Teórico

Un Random Forest es un **ensemble de árboles de decisión** que combina predicciones mediante votación mayoritaria (clasificación) o promedio (regresión).

**Definición Formal:**

Sea $\{h(\mathbf{x}, \Theta_k)\}_{k=1}^K$ un conjunto de $K$ árboles de decisión, donde $\Theta_k$ son parámetros aleatorios (subset de features y datos). El Random Forest predice:

\end{equation}

\hat{y} = \text{mode}\{h(\mathbf{x}, \Theta_1), h(\mathbf{x}, \Theta_2), ..., h(\mathbf{x}, \Theta_K)\}

\begin{equation}

---

### 3.2 Algoritmo de Entrenamiento

\begin{lstlisting}[language=Python]
def entrenar_random_forest(X_train, y_train):
    """
    X_train: matriz n × d de características
    y_train: vector n de salones asignados (labels)
    """
    forest = []
    
    for k in range(K):  # K = número de árboles
        # 1. Bootstrap sampling
        indices = sample_with_replacement(n, n)
        X_boot = X_train[indices]
        y_boot = y_train[indices]
        
        # 2. Entrenar árbol con subset aleatorio de features
        tree = DecisionTree(
            max_depth=d_max,
            min_samples_split=s_min,
            max_features=sqrt(d)  # Regla empírica
        )
        tree.fit(X_boot, y_boot)
        forest.append(tree)
    
    return forest
\end{verbatim}

---

### 3.3 Construcción de Árbol de Decisión

**Algoritmo CART (Classification and Regression Trees):**

\begin{lstlisting}[language=Python]
función construir_arbol(X, y, profundidad):
    si profundidad == max_depth o |y| < min_samples:
        retornar hoja con clase mayoritaria
    
    mejor_ganancia = 0
    mejor_split = None
    
    para cada característica f en subset_aleatorio(features):
        para cada valor v en valores_unicos(X[:, f]):
            # Dividir datos
            izq = {(x, y) : x[f] <= v}
            der = {(x, y) : x[f] > v}
\end{verbatim}


---

### 3.3 Construcción de Árbol (continuación)

\begin{lstlisting}[language=Python]
            # Calcular ganancia de información
            ganancia = calcular_ganancia(y, izq, der)
            
            si ganancia > mejor_ganancia:
                mejor_ganancia = ganancia
                mejor_split = (f, v)
    
    si mejor_ganancia == 0:
        retornar hoja
    
    # Crear nodo interno
    nodo.feature = mejor_split[0]
    nodo.threshold = mejor_split[1]
    nodo.izquierdo = construir_arbol(izq, profundidad + 1)
    nodo.derecho = construir_arbol(der, profundidad + 1)
    
    retornar nodo
\end{verbatim}

---


### 3.4 Criterio de División: Gini Impurity

Para clasificación, usamos el **índice de Gini**:

\end{equation}

\begin{equation}

Donde:
- $S$ = conjunto de muestras en el nodo
- $C$ = conjunto de clases (salones)
- $p_i$ = proporción de muestras de clase $i$ en $S$
---

**Ganancia de Información:**

\end{equation}

\begin{equation}

**Objetivo:** Maximizar la ganancia (minimizar impureza después del split)

---

### 3.5 Predicción

Para una nueva clase $\mathbf{x}_{new}$:

\end{equation}

\hat{y}_{RF} = \arg\max_{s \in S} \sum_{k=1}^{K} \mathbb{1}[h_k(\mathbf{x}_{new}) = s]

\begin{equation}

**Interpretación:** El salón que recibe más "votos" de los árboles individuales.

---

### 3.6 Hiperparámetros

\begin{lstlisting}[language=Python]
RandomForestClassifier(
    n_estimators=100,        # Número de árboles
    max_depth=20,            # Profundidad máxima
    min_samples_split=5,     # Mínimo para dividir nodo
    min_samples_leaf=2,      # Mínimo en hojas
    max_features='sqrt',     # sqrt(d) features por split
    random_state=42,         # Reproducibilidad
    n_jobs=-1                # Paralelización
)
\end{verbatim}


**Justificación de Valores:**
- `n_estimators=100`: Balance entre precisión y tiempo
- `max_depth=20`: Evita overfitting en dataset pequeño
- `min_samples_split=5`: Previene splits en ruido
- `max_features='sqrt'`: Regla empírica de Breiman


---

## 4. Modelo 2: Gradient Boosting Regressor
---

### 4.1 Fundamento Teórico

Gradient Boosting construye un **ensemble aditivo** de árboles débiles que minimizan una función de pérdida mediante descenso de gradiente.

**Idea Central:** Cada árbol nuevo corrige los errores del ensemble anterior.

---

### 4.2 Algoritmo de Gradient Boosting

**Formulación Matemática:**

Queremos aproximar una función $F^*(\mathbf{x})$ que minimiza la pérdida esperada:

\end{equation}

F^*(\mathbf{x}) = \arg\min_{F} \mathbb{E}_{y, \mathbf{x}}[L(y, F(\mathbf{x}))]

\begin{equation}

Donde $L$ es la función de pérdida (ej: MSE para regresión).

---

**Algoritmo:**

\begin{verbatim}
Inicializar: F_0(x) = arg min_γ Σ L(y_i, γ)

Para m = 1 hasta M:
    1. Calcular pseudo-residuos:
       r_im = -[∂L(y_i, F(x_i))/∂F(x_i)]_{F=F_{m-1}}
    
    2. Entrenar árbol h_m(x) para predecir r_i
    
    3. Calcular multiplicador óptimo:
       γ_m = arg min_γ Σ L(y_i, F_{m-1}(x_i) + γ h_m(x_i))
    
    4. Actualizar modelo:
       F_m(x) = F_{m-1}(x) + η γ_m h_m(x)

Retornar F_M(x)
\end{verbatim}


---

Donde:
- $M$ = número de iteraciones (árboles)
- $\eta$ = learning rate (tasa de aprendizaje)
- $h_m$ = árbol débil en iteración $m$

---

### 4.3 Función de Pérdida: Mean Squared Error

Para regresión de calidad de asignación:

\end{equation}

L(y, F(\mathbf{x})) = \frac{1}{2}(y - F(\mathbf{x}))^2

\begin{equation}

**Gradiente:**

\end{equation}

\frac{\partial L}{\partial F} = -(y - F(\mathbf{x})) = -residuo

\begin{equation}

Por lo tanto, los pseudo-residuos son simplemente los residuos reales.

---

### 4.4 Regularización

**L2 Regularization (Ridge):**

\end{equation}

\begin{equation}

**Shrinkage (Learning Rate):**

\end{equation}

F_m(\mathbf{x}) = F_{m-1}(\mathbf{x}) + \eta \cdot \gamma_m \cdot h_m(\mathbf{x})

\begin{equation}

Donde $0 < \eta \leq 1$ controla la contribución de cada árbol.

**Subsampling:**
Usar solo una fracción $\rho$ de los datos para entrenar cada árbol:

\end{equation}

subsample\_size = \rho \cdot n, \quad 0 < \rho \leq 1

\begin{equation}

---

### 4.5 Hiperparámetros

\begin{lstlisting}[language=Python]
GradientBoostingRegressor(
    n_estimators=100,        # Número de boosting stages
    learning_rate=0.1,       # Shrinkage η
    max_depth=5,             # Profundidad de árboles débiles
    min_samples_split=5,
    min_samples_leaf=2,
    subsample=0.8,           # Stochastic GB
    random_state=42
)
\end{verbatim}


**Justificación:**
- `learning_rate=0.1`: Balance entre convergencia y generalización
- `max_depth=5`: Árboles débiles (shallow) para boosting
- `subsample=0.8`: Reduce overfitting, acelera entrenamiento


---

## 5. Pipeline de Optimización ML
---

### 5.1 Fase de Entrenamiento

\begin{lstlisting}[language=Python]
def entrenar(self, df_inicial):
    """
    Entrena ambos modelos usando el horario inicial
    """
    # 1. Extraer características y labels
    X = []
    y_salon = []
    y_calidad = []
    
    for idx, clase in df_inicial.iterrows():
        features = self.extraer_features(clase, df_inicial, idx)
        X.append(features)
        y_salon.append(clase['Salon'])
        
        # Calidad = inverso de energía local
        calidad = self.calcular_calidad_local(clase, df_inicial)
        y_calidad.append(calidad)
    
    X = np.array(X)
    y_salon = np.array(y_salon)
    y_calidad = np.array(y_calidad)
\end{verbatim}


---

### 5.1 Fase de Entrenamiento (continuación)

\begin{lstlisting}[language=Python]
    # 2. Entrenar clasificador (predicción de salón)
    self.clasificador.fit(X, y_salon)
    
    # 3. Entrenar regresor (calidad de asignación)
    self.regressor_calidad.fit(X, y_calidad)
    
    # 4. Calcular métricas de entrenamiento
    y_pred = self.clasificador.predict(X)
    accuracy = (y_pred == y_salon).mean()
    
    return {
        'accuracy': accuracy,
        'n_samples': len(X),
        'n_features': X.shape[1]
    }
\end{verbatim}


---

### 5.2 Cálculo de Calidad Local

La calidad de una asignación se define como:

\end{equation}

calidad(c_i, s_j) = -energia\_local(c_i, s_j)

\begin{equation}

Donde:

\end{equation}

\begin{align}

energia\_local(c_i, s_j) = &\ w_1 \cdot distancia\_anterior(c_i, s_j) \\

                           &+ w_2 \cdot \mathbb{1}[tipo(c_i) \neq tipo(s_j)] \\

                           &+ w_3 \cdot \mathbb{1}[s_j \in S_{invalidos}] \\

                           &+ w_4 \cdot \mathbb{1}[cambio\_piso(c_i, s_j)]

\end{align}

\begin{equation}

**Objetivo:** Enseñar al modelo qué hace una asignación "buena" vs "mala"

---

### 5.3 Fase de Optimización

\begin{lstlisting}[language=Python]
def optimizar(self, df_inicial):
    """
    Genera nueva asignación usando modelos entrenados
    """
    df_resultado = df_inicial.copy()
    cambios = 0
    
    for idx in range(len(df_resultado)):
        clase = df_resultado.iloc[idx]
        
        # Proteger clases inmutables (P1)
        if idx in self.indices_inmutables:
            continue
        
        # 1. Extraer características
        features = self.extraer_features(clase, df_resultado, idx)
        
        # 2. Obtener salones candidatos
        candidatos = self.obtener_salones_validos(clase)
        
        # 3. Predecir mejor salón
        salon_predicho = self.clasificador.predict([features])[0]
        
        # 4. Si predicción es válida, usar directamente
        if salon_predicho in candidatos:
            mejor_salon = salon_predicho
\end{verbatim}


---

### 5.3 Fase de Optimización (continuación)

\begin{lstlisting}[language=Python]
        else:
            # 5. Evaluar todos los candidatos con regresor
            mejor_salon = None
            mejor_calidad = -infinito
            
            for salon in candidatos:
                # Simular asignación
                features_temp = self.modificar_features(features, salon)
                calidad = self.regressor_calidad.predict([features_temp])[0]
                
                if calidad > mejor_calidad:
                    mejor_calidad = calidad
                    mejor_salon = salon
        
        # 6. Aplicar asignación
        if df_resultado.loc[idx, 'Salon'] != mejor_salon:
            df_resultado.loc[idx, 'Salon'] = mejor_salon
            cambios += 1
    
    return df_resultado, cambios
\end{verbatim}


---

### 5.4 Validación de Asignaciones

Después de cada predicción, validamos:

\begin{lstlisting}[language=Python]
def validar_asignacion(self, clase, salon, df_actual):
    """
    Verifica que la asignación sea factible
    """
    # 1. Tipo correcto
    if tipo(clase) != tipo(salon):
        return False
    
    # 2. No es inválido
    if salon in self.salones_invalidos:
        return False
\end{verbatim}


---

### 5.4 Validación de Asignaciones (continuación)

\begin{lstlisting}[language=Python]
    # 3. No hay conflicto temporal
    conflicto = df_actual[
        (df_actual['Dia'] == clase['Dia']) &
        (df_actual['Bloque_Horario'] == clase['Bloque_Horario']) &
        (df_actual['Salon'] == salon)
    ]
    if len(conflicto) > 0:
        return False
    
    return True
\end{verbatim}



---

## 6. Análisis de Complejidad
---

### 6.1 Complejidad Temporal

**Entrenamiento:**

\end{equation}

T_{train} = O(K \cdot n \cdot d \cdot \log n \cdot h)

\begin{equation}

Donde:
- $K$ = número de árboles
- $n$ = número de muestras
- $d$ = número de características
- $h$ = profundidad máxima

---

**Para nuestro caso:**
- $K = 100$
- $n = 680$
- $d \approx 20$
- $h = 20$

\end{equation}

T_{train} \approx 100 \cdot 680 \cdot 20 \cdot \log(680) \cdot 20 \approx 3.7 \times 10^7 \text{ operaciones}

\begin{equation}

**Optimización (Inferencia):**

\end{equation}

T_{opt} = O(n \cdot K \cdot h \cdot d)

\begin{equation}

\end{equation}

T_{opt} \approx 680 \cdot 100 \cdot 20 \cdot 20 \approx 2.7 \times 10^7 \text{ operaciones}

\begin{equation}

**Tiempo Real:** ~15-20 segundos en hardware moderno

---

### 6.2 Complejidad Espacial

\end{equation}

S = O(K \cdot n_{nodes} + n \cdot d)

\begin{equation}

Donde $n_{nodes}$ es el número promedio de nodos por árbol.

Para árboles balanceados de profundidad $h$:

\end{equation}

n_{nodes} \approx 2^h - 1

\begin{equation}

\end{equation}

S \approx 100 \cdot (2^{20} - 1) + 680 \cdot 20 \approx 10^8 \text{ bytes} \approx 100 \text{ MB}

\begin{equation}


---

## 7. Ventajas y Limitaciones
---

### 7.1 Ventajas

✅ **Velocidad:** Inferencia muy rápida una vez entrenado  
✅ **Aprendizaje:** Mejora con más datos históricos  
✅ **Robustez:** Ensemble reduce varianza  
✅ **Interpretabilidad:** Feature importance muestra qué importa  
✅ **No paramétrico:** No asume distribución de datos  

---

### 7.2 Limitaciones

❌ **Requiere datos:** Necesita horarios previos de calidad  
❌ **Overfitting:** Puede memorizar patrones específicos  
❌ **Exploración limitada:** No explora fuera de lo aprendido  
❌ **Dependencia de features:** Calidad depende de feature engineering  
❌ **Cold start:** Mal rendimiento sin datos de entrenamiento  


---

## 8. Mejoras y Extensiones
---

### 8.1 Transfer Learning

Usar modelos pre-entrenados en otros campus/instituciones:

\end{equation}

\theta_{nuevo} = \theta_{pretrained} + \Delta\theta_{fine-tune}

\begin{equation}

---

### 8.2 Active Learning

Seleccionar ejemplos más informativos para etiquetar:

\end{equation}

x^* = \arg\max_{x \in U} uncertainty(x)

\begin{equation}

Donde $uncertainty$ puede ser entropía de predicción.

---

### 8.3 Deep Learning

Reemplazar Random Forest con redes neuronales:

\begin{verbatim}
Input (features) → Dense(128, ReLU) → Dropout(0.3) 
                 → Dense(64, ReLU) → Dropout(0.3)
                 → Dense(|S|, Softmax) → Output (salon)
\end{verbatim}


---

### 8.4 Reinforcement Learning

Formular como MDP (Markov Decision Process):
- **Estado:** Asignación parcial actual
- **Acción:** Asignar clase $c_i$ a salón $s_j$
- **Recompensa:** $-energia(asignacion)$
- **Política:** $\pi(s, a) = P(a|s)$

Usar Q-Learning o Policy Gradient para aprender política óptima.


---

## 9. Resultados Experimentales
---

### 9.1 Métricas de Entrenamiento

\begin{verbatim}
Modelo: Random Forest Classifier
├── Accuracy: 0.85
├── Precision: 0.83
├── Recall: 0.82
└── F1-Score: 0.82

Modelo: Gradient Boosting Regressor
├── R²: 0.78
├── MSE: 45.2
└── MAE: 5.3
\end{verbatim}


---

### 9.2 Comparación con Otros Métodos

| Métrica | ML | Greedy | Genético |
|---------|-----|--------|----------|
| Tiempo | **15.8s** | 29.3s | 73.9s |
| P1 | 100% | 100% | 100% |
| Distancia | **1821** | 1951 | 2413 |
| Consistencia | Media | Alta | Baja |

---

### 9.3 Feature Importance

\begin{verbatim}
Top 10 características más importantes:
1. profesor_encoded (0.18)
2. salon_clase_anterior (0.15)
3. tipo_clase (0.12)
4. tiene_preferencia (0.10)
5. bloque_horario (0.09)
6. materia_encoded (0.08)
7. dia_semana (0.07)
8. num_estudiantes (0.06)
9. semestre (0.05)
10. piso_clase_anterior (0.04)
\end{verbatim}


**Interpretación:** El profesor y el contexto de la clase anterior son los factores más predictivos.


---

## 10. Conclusiones

El enfoque de Machine Learning ofrece una alternativa **rápida y efectiva** para la asignación de salones, especialmente cuando:
- Existen datos históricos de calidad
- Se requiere velocidad de ejecución
- Los patrones son relativamente estables

Sin embargo, requiere **cuidadoso feature engineering** y **datos de entrenamiento representativos** para alcanzar su máximo potencial.


---

## Referencias

1. Breiman, L. (2001). Random forests. *Machine Learning*, 45(1), 5-32.

2. Friedman, J. H. (2001). Greedy function approximation: a gradient boosting machine. *Annals of statistics*, 1189-1232.

3. Hastie, T., Tibshirani, R., & Friedman, J. (2009). *The elements of statistical learning* (2nd ed.). Springer.

4. Chen, T., & Guestrin, C. (2016). Xgboost: A scalable tree boosting system. *KDD*, 785-794.

5. Ke, G., et al. (2017). Lightgbm: A highly efficient gradient boosting decision tree. *NIPS*, 3146-3154.


---

## 📖 Ejemplo Didáctico: Machine Learning

> [!NOTE]
> **Este es un ejemplo didáctico simplificado** con un problema pequeño (10 clases, 4 salones, 4 profesores) para demostrar cómo funciona el enfoque de Machine Learning de manera clara y comprensible.

**Problema Simplificado para Demostración**

---

### Problema Didáctico

**Datos:**
- 10 clases
- 4 salones  
- 4 profesores
- Objetivo: Predecir asignaciones óptimas

**Archivo:** `ejemplos_didacticos/02_machine_learning.py`

**Propósito:** Demostrar cómo ML aprende patrones de asignación

---

### Modelo Simplificado

**Reglas Aprendidas:**
\begin{lstlisting}[language=Python]
class ModeloSimple:
    def entrenar(self, X, y):
        # Reglas basadas en características
        self.reglas = [
            ("Lab", lambda f: f[1] > 0.5, 3),      # Lab → S4
            ("Muchos est.", lambda f: f[0] > 0.54, 0),  # >27 → S1
            ("Default", lambda f: True, 1),         # Resto → S2
        ]
    
    def predecir(self, features):
        for _, condicion, salon_idx in self.reglas:
            if condicion(features):
                return salon_idx
        return 1
\end{verbatim}

---

**Features extraídas:**
- `f[0]`: Número de estudiantes (normalizado)
- `f[1]`: Tipo de clase (0=Teoría, 1=Lab)

---

### Resultados del Ejemplo

**Asignación Aleatoria (Baseline):**
\begin{verbatim}
Movimientos de profesores: 5
\end{verbatim}


**Predicciones con ML:**
\begin{verbatim}
Movimientos de profesores: 3
\end{verbatim}


**Mejora Lograda:** 40% ✅

**Conclusión:** ML reduce movimientos al aprender patrones óptimos

---

### Ejecutar el Ejemplo

\begin{lstlisting}[language=Bash]
cd ejemplos_didacticos
python3 02_machine_learning.py
\end{verbatim}


**Salida esperada:**
\begin{verbatim}
============================================================
EJEMPLO: Machine Learning
============================================================
Entrenando modelo...

🎲 Asignación aleatoria (sin ML):
Movimientos: 5

🤖 Predicciones con ML:
  C1 → S2
  C2 → S1
  C3 → S3
  C4 → S2
  ...
\end{verbatim}


---

### Resultados del Ejemplo ML

\begin{verbatim}
📊 Comparación:
  Movimientos (aleatorio): 5
  Movimientos (ML): 3
  Mejora: 40.0%

✅ Completado!
\end{verbatim}


---

<!-- _class: lead blue -->
# Algoritmo Genético

**Búsqueda Evolutiva Global**

---

## 1. Introducción y Fundamento Teóricológico


---

### 1.1 Inspiración Evolutiva

Los **Algoritmos Genéticos (AG)** son metaheurísticas inspiradas en la evolución biológica de Charles Darwin. Simulan el proceso de selección natural donde los individuos más aptos tienen mayor probabilidad de sobrevivir y reproducirse.

**Principios Fundamentales:**

1. **Variación:** Diversidad genética mediante mutación y recombinación
2. **Herencia:** Transmisión de características a descendientes
3. **Selección:** Supervivencia del más apto
4. **Adaptación:** Mejora gradual de la población
---

**Teorema de Holland (Schema Theorem):**
*Los esquemas (patrones de solución) cortos, de bajo orden y con alta aptitud reciben un número exponencialmente creciente de muestras en generaciones sucesivas.*


---

### 1.2 Analogía Biológica vs Computacional

| Biología | Algoritmo Genético |
|----------|-------------------|
| Individuo | Solución candidata |
| Cromosoma | Representación de solución |
| Gen | Variable de decisión |
| Alelo | Valor de variable |
| Población | Conjunto de soluciones |

---

### 1.2 Analogía Biológica vs Computacional (continuación)

| Biología | Algoritmo Genético |
|----------|-------------------|
| Generación | Iteración del algoritmo |
| Fitness | Calidad de solución |
| Selección natural | Selección por aptitud |
| Cruce sexual | Operador de cruce |
| Mutación | Operador de mutación |


---

### 1.3 Ventajas Teóricas

✅ **Exploración global:** Mantiene población diversa  
✅ **Paralelismo implícito:** Evalúa múltiples soluciones simultáneamente  
✅ **Robustez:** No requiere gradientes ni derivadas  
✅ **Flexibilidad:** Aplicable a problemas discretos y continuos  
✅ **Escape de óptimos locales:** Mediante mutación y diversidad  


---

## 2. Representación Cromosómica


---

### 2.1 Codificación de Soluciones

Para el problema de asignación de salones, usamos **codificación directa**:

**Cromosoma:**
\end{equation}

\mathbf{x} = [x_1, x_2, ..., x_n]

\begin{equation}

Donde:
- $n = 680$ (número de clases)
- $x_i \in S$ (salón asignado a clase $i$)

**Ejemplo:**
\begin{verbatim}
Clase:  [c1,  c2,  c3,  c4,  c5,  ...]
Salón:  [FF1, LBD, FF2, FF1, LIA, ...]
         ↑    ↑    ↑    ↑    ↑
         Gen  Gen  Gen  Gen  Gen
\end{verbatim}



---

### 2.2 Espacio de Búsqueda

El espacio de búsqueda $\Omega$ es:

\end{equation}

\Omega = S^n = \{(s_1, s_2, ..., s_n) : s_i \in S\}

\begin{equation}

\end{equation}

\begin{equation}

**Subespacio Factible:**

\end{equation}

\Omega_{factible} = \{\mathbf{x} \in \Omega : satisface\ restricciones\ duras\}

\begin{equation}

\end{equation}

\begin{equation}


---

### 2.3 Función de Aptitud (Fitness)

La aptitud mide la calidad de una solución:

\end{equation}

fitness(\mathbf{x}) = -E(\mathbf{x})

\begin{equation}

Donde $E(\mathbf{x})$ es la función de energía definida anteriormente.

**Normalización:**

Para mantener valores positivos y facilitar selección proporcional:

\end{equation}

fitness_{norm}(\mathbf{x}) = \frac{1}{1 + E(\mathbf{x})}

\begin{equation}

O usando ranking:

\end{equation}

fitness_{rank}(\mathbf{x}) = rank(\mathbf{x})

\begin{equation}

Donde $rank(\mathbf{x}) \in \{1, 2, ..., |P|\}$ es la posición en la población ordenada.


---

## 3. Algoritmo Principal


---

### 3.1 Pseudocódigo General

\begin{verbatim}
Algoritmo Genético(problema, parámetros):
    # Inicialización
    P₀ = generar_poblacion_inicial(tam_poblacion)
    evaluar_fitness(P₀)
    t = 0
    
    # Evolución
    mientras t < max_generaciones y no convergió:
        # Selección
        padres = seleccionar_padres(Pₜ)
        
        # Cruce
        hijos = aplicar_cruce(padres, prob_cruce)
        
        # Mutación
        hijos = aplicar_mutacion(hijos, prob_mutacion)
        
        # Reparación (si necesario)
        hijos = reparar_soluciones(hijos)
        
        # Evaluación
        evaluar_fitness(hijos)
        
        # Reemplazo
        Pₜ₊₁ = seleccionar_supervivientes(Pₜ, hijos)
        
        t = t + 1
    
    retornar mejor_individuo(Pₜ)
\end{verbatim}



---

### 3.2 Parámetros del Algoritmo

\begin{lstlisting}[language=Python]
class OptimizadorGenetico:
    def __init__(self):
        # Parámetros poblacionales
        self.tam_poblacion = 150
        self.num_generaciones = 500
        
        # Probabilidades de operadores
        self.prob_cruce = 0.8
        self.prob_mutacion_inicial = 0.1
        self.prob_mutacion_final = 0.3
        
        # Selección
        self.tam_torneo = 5
        self.elitismo = 0.1  # 10% mejores pasan directamente
        
        # Diversidad
        self.penalizacion_diversidad = True
        self.min_diversidad = 0.3
\end{verbatim}


---

**Justificación de Valores:**

- **tam_poblacion = 150:** Balance entre diversidad y costo computacional
- **num_generaciones = 500:** Suficiente para convergencia en problemas complejos
- **prob_cruce = 0.8:** Valor estándar en literatura (Holland, 1975)
- **prob_mutacion adaptativa:** Aumenta con generaciones para escapar estancamiento


---

## 4. Operadores Genéticos


---

### 4.1 Inicialización de Población

**Estrategia Híbrida:**

\begin{lstlisting}[language=Python]
def generar_poblacion_inicial(self, df):
    poblacion = []
    
    # 1. Incluir solución pre-asignada (élite)
    sol_inicial = self.cargar_solucion_inicial(df)
    poblacion.append(sol_inicial)
    
    # 2. Generar variaciones de la inicial (30%)
    for _ in range(int(0.3 * self.tam_poblacion)):
        sol_variacion = self.perturbar_solucion(sol_inicial)
        poblacion.append(sol_variacion)
    
    # 3. Generar soluciones aleatorias (70%)
    for _ in range(self.tam_poblacion - len(poblacion)):
        sol_aleatoria = self.generar_solucion_aleatoria(df)
        poblacion.append(sol_aleatoria)
    
    return poblacion
\end{verbatim}


**Generación Aleatoria:**

\end{equation}

x_i = \begin{cases}

pref(c_i) & \text{si } c_i \in P_1 \text{ (inmutable)} \\

random\_choice(salones\_validos(c_i)) & \text{en otro caso}

\end{cases}

\begin{equation}


---

### 4.2 Selección de Padres

Implementamos **Selección por Torneo**:

**Algoritmo:**

\begin{lstlisting}[language=Python]
def seleccion_torneo(self, poblacion, fitness, k=5):
    """
    Selecciona un individuo mediante torneo de tamaño k
    """
    # Seleccionar k individuos aleatorios
    competidores = random.sample(range(len(poblacion)), k)
    
    # Retornar el de mayor fitness
    mejor_idx = max(competidores, key=lambda i: fitness[i])
    
    return poblacion[mejor_idx]
\end{verbatim}


---

**Presión Selectiva:**

La probabilidad de que el mejor individuo sea seleccionado en un torneo de tamaño $k$:

\end{equation}

\begin{equation}

Para $|P| = 150$, $k = 5$:

\end{equation}

P \approx 1 - (0.993)^5 \approx 0.034 = 3.4\%

\begin{equation}

**Ventajas del Torneo:**
- No requiere ordenamiento completo
- Fácil paralelización
- Presión selectiva ajustable (mediante $k$)


---

### 4.3 Operador de Cruce

Implementamos **Cruce de Un Punto con Protección**:

**Algoritmo:**

\begin{lstlisting}[language=Python]
def cruce_un_punto(self, padre1, padre2):
    """
    Cruce de un punto respetando genes inmutables
    """
    n = len(padre1)
    
    # Seleccionar punto de cruce aleatorio
    punto = random.randint(1, n-1)
    
    # Crear hijos
    hijo1 = padre1.copy()
    hijo2 = padre2.copy()
    
    # Intercambiar segmentos (excepto inmutables)
    for i in range(punto, n):
        if i not in self.indices_inmutables:
            hijo1[i] = padre2[i]
            hijo2[i] = padre1[i]
    
    return hijo1, hijo2
\end{verbatim}


---

**Representación Gráfica:**

\begin{verbatim}
Padre 1: [FF1|FF2|FF3|FF4|FF5|FF6|FF7]
                    ↓ punto de cruce
Padre 2: [LBD|LIA|FF8|FF9|FFA|FFB|FFC]
         ─────────────────────────────
Hijo 1:  [FF1|FF2|FF3|FF9|FFA|FFB|FFC]
Hijo 2:  [LBD|LIA|FF8|FF4|FF5|FF6|FF7]
\end{verbatim}

---

**Teorema del Building Block:**
*El cruce preserva y combina bloques de construcción (schemata) de alta aptitud.*

**Demostración (informal):**
- Si padre1 tiene buen bloque en posiciones [1-3]
- Y padre2 tiene buen bloque en posiciones [5-7]
- Cruce en posición 4 puede combinar ambos bloques en hijo


---

### 4.4 Operador de Mutación

**Mutación Adaptativa por Intercambio:**

\begin{lstlisting}[language=Python]
def mutacion_intercambio(self, individuo, generacion):
    """
    Muta mediante intercambio de genes compatibles
    Probabilidad aumenta con generaciones
    """
    # Probabilidad adaptativa
    progreso = generacion / self.num_generaciones
    prob_mut = self.prob_mutacion_inicial + \
               (self.prob_mutacion_final - self.prob_mutacion_inicial) * progreso
    
    individuo_mutado = individuo.copy()
\end{verbatim}


---

### 4.4 Operador de Mutación (continuación)

\begin{lstlisting}[language=Python]
    for i in range(len(individuo)):
        if random.random() < prob_mut:
            # No mutar inmutables
            if i in self.indices_inmutables:
                continue
            
            # Seleccionar gen compatible para intercambiar
            j = self.seleccionar_gen_compatible(i, individuo)
            
            if j is not None and j not in self.indices_inmutables:
                # Intercambiar
                individuo_mutado[i], individuo_mutado[j] = \
                    individuo_mutado[j], individuo_mutado[i]
    
    return individuo_mutado
\end{verbatim}


---

**Probabilidad Adaptativa:**

\end{equation}

p_{mut}(t) = p_{min} + (p_{max} - p_{min}) \cdot \frac{t}{T}

\begin{equation}

Donde:
- $p_{min} = 0.1$ (inicial)
- $p_{max} = 0.3$ (final)
- $t$ = generación actual
- $T$ = total de generaciones

**Justificación:**
- **Inicio:** Baja mutación para explotar buenos esquemas
- **Final:** Alta mutación para escapar estancamiento


---

### 4.5 Reparación de Soluciones

Después de cruce y mutación, pueden generarse soluciones inválidas:

\begin{lstlisting}[language=Python]
def reparar_solucion(self, individuo, df):
    """
    Repara violaciones de restricciones duras
    """
    individuo_reparado = individuo.copy()
    
    for i in range(len(individuo)):
        clase = df.iloc[i]
        salon_actual = individuo[i]
        
        # Verificar validez
        if not self.es_asignacion_valida(clase, salon_actual, individuo, i):
\end{verbatim}


---

### 4.5 Reparación de Soluciones (continuación)

\begin{lstlisting}[language=Python]
            # Buscar salón válido
            salones_validos = self.obtener_salones_validos(clase)
            
            for salon in salones_validos:
                if self.es_asignacion_valida(clase, salon, individuo, i):
                    individuo_reparado[i] = salon
                    break
    
    return individuo_reparado
\end{verbatim}


---

**Teorema de Factibilidad:**
*Toda solución puede repararse a una solución factible si existe al menos una asignación válida para cada clase.*


---

## 5. Estrategias de Reemplazo


---

### 5.1 Reemplazo Generacional con Elitismo

\begin{lstlisting}[language=Python]
def seleccionar_supervivientes(self, poblacion, hijos, fitness_pob, fitness_hijos):
    """
    Combina población actual e hijos, selecciona mejores
    """
    # Calcular número de élites
    num_elites = int(self.elitismo * len(poblacion))
    
    # Identificar élites
    indices_ordenados = sorted(range(len(poblacion)), 
                               key=lambda i: fitness_pob[i], 
                               reverse=True)
    elites = [poblacion[i] for i in indices_ordenados[:num_elites]]
    
    # Combinar hijos con no-élites
    nueva_poblacion = elites + hijos[:len(poblacion) - num_elites]
    
    return nueva_poblacion
\end{verbatim}


---

**Tasa de Elitismo:**

\end{equation}

\begin{equation}
---

**Teorema de Convergencia con Elitismo:**
*Un AG con elitismo converge al óptimo global con probabilidad 1 cuando $t \rightarrow \infty$.*

**Demostración (sketch):**
1. Elitismo preserva mejor solución encontrada
2. Mutación garantiza $P(visitar\ cualquier\ solucion) > 0$
3. Con tiempo infinito, se visitará el óptimo
4. Elitismo lo preservará una vez encontrado


---

## 6. Control de Diversidad


---

### 6.1 Medida de Diversidad

Definimos diversidad de población como:

\end{equation}

\begin{equation}

Donde:

\end{equation}

distancia(\mathbf{x}_i, \mathbf{x}_j) = \frac{1}{n} \sum_{k=1}^{n} \mathbb{1}[x_{i,k} \neq x_{j,k}]

\begin{equation}

**Interpretación:** Proporción promedio de genes diferentes entre individuos.


---

### 6.2 Sharing (Compartición de Fitness)

Para mantener diversidad, penalizamos individuos similares:

\end{equation}

\begin{equation}

Donde la función de sharing es:

\end{equation}

sh(d) = \begin{cases}

1 - \left(\frac{d}{\sigma_{share}}\right)^\alpha & \text{si } d < \sigma_{share} \\

0 & \text{en otro caso}

\end{cases}

\begin{equation}

Parámetros típicos:
- $\sigma_{share} = 0.1$ (radio de nicho)
- $\alpha = 2$ (forma de la función)


---

### 6.3 Reinicio Adaptativo

Si diversidad cae por debajo de umbral:

\begin{lstlisting}[language=Python]
def verificar_y_reiniciar(self, poblacion, generacion):
    """
    Reinicia población si diversidad es muy baja
    """
    diversidad = self.calcular_diversidad(poblacion)
    
    if diversidad < self.min_diversidad:
        print(f"   ⚠️  Diversidad baja ({diversidad:.3f}), reiniciando...")
        
        # Preservar mejores
        num_preservar = int(0.2 * len(poblacion))
        mejores = self.obtener_mejores(poblacion, num_preservar)
        
        # Generar nuevos
        nuevos = self.generar_poblacion_inicial(len(poblacion) - num_preservar)
        
        return mejores + nuevos
    
    return poblacion
\end{verbatim}



---

## 7. Análisis de Complejidad


---

### 7.1 Complejidad Temporal

**Por generación:**

\end{equation}

T_{gen} = T_{eval} + T_{sel} + T_{cruce} + T_{mut} + T_{reemplazo}

\begin{equation}

Donde:
- $T_{eval} = O(|P| \cdot  n)$ (evaluar población)
- $T_{sel} = O(|P| \cdot  k)$ (selección por torneo)
- $T_{cruce} = O(|P| \cdot  n)$ (aplicar cruce)
- $T_{mut} = O(|P| \cdot  n)$ (aplicar mutación)
- $T_{reemplazo} = O(|P| \log |P|)$ (ordenar para elitismo)

\end{equation}

\begin{equation}

**Total:**

\end{equation}

\begin{equation}

Para $G = 500$, $|P| = 150$, $n = 680$:

\end{equation}

T_{total} \approx 5.1 \times 10^7 \text{ operaciones}

\begin{equation}

**Tiempo real:** ~70-80 segundos


---

### 7.2 Complejidad Espacial

\end{equation}

\begin{equation}

Donde:
- $|P| \cdot  n$ = almacenar población
- $|P|$ = almacenar fitness

\end{equation}

S \approx 150 \cdot 680 + 150 \approx 102,150 \text{ valores}

\begin{equation}

**Memoria:** ~800 KB (despreciable)


---

## 8. Criterios de Parada

El algoritmo se detiene cuando se cumple alguna de estas condiciones:


---

### 8.1 Número Máximo de Generaciones

\end{equation}

t \geq G_{max}

\begin{equation}


---

### 8.2 Convergencia de Fitness

\end{equation}

\frac{fitness_{mejor} - fitness_{promedio}}{fitness_{mejor}} < \epsilon

\begin{equation}

Donde $\epsilon = 0.01$ (1% de diferencia)


---

### 8.3 Estancamiento

No hay mejora en $k$ generaciones consecutivas:

\end{equation}

\forall i \in [t-k, t]: fitness_{mejor}(i) = fitness_{mejor}(t)

\begin{equation}

Típicamente $k = 50$


---

## 9. Resultados y Análisis


---

### 9.1 Evolución de Fitness

\begin{verbatim}
Generación    Mejor      Promedio    Peor       Diversidad
─────────────────────────────────────────────────────────────
0             85,234     92,456      98,123     0.85
50            83,112     86,234      91,456     0.72
100           82,456     84,123      88,234     0.65
150           82,234     83,456      86,123     0.58
200           82,156     83,234      85,456     0.52
250           82,089     82,987      84,234     0.45
300           82,034     82,756      83,456     0.38
350           81,998     82,645      83,123     0.31
400           81,976     82,578      82,987     0.28
450           81,967     82,534      82,876     0.25
500           81,962     82,512      82,789     0.23
\end{verbatim}

---

**Observaciones:**
- Mejora rápida en primeras 100 generaciones
- Convergencia gradual después
- Diversidad decrece naturalmente
- Diferencia mejor-promedio se reduce (convergencia)


---

### 9.2 Comparación con Otros Algoritmos

| Métrica | Genético | Greedy+HC | ML |
|---------|----------|-----------|-----|
| Tiempo | 73.9s | 29.3s | 15.8s |
| Mejor solución | 81,962 | 82,716 | 82,234 |
| Consistencia | Baja | Alta | Media |
| Exploración | Excelente | Limitada | Limitada |
| Explotación | Buena | Excelente | Buena |


---

### 9.3 Análisis de Operadores

**Impacto del Cruce:**

| Prob. Cruce | Mejor Fitness | Generaciones |
|-------------|---------------|--------------|
| 0.5 | 82,456 | 520 |
| 0.7 | 82,123 | 480 |
| 0.8 | 81,962 | 500 |
| 0.9 | 82,234 | 510 |

**Óptimo:** 0.8 (valor estándar)

---

**Impacto de la Mutación:**

| Prob. Mutación | Mejor Fitness | Diversidad Final |
|----------------|---------------|------------------|
| 0.05 | 82,567 | 0.15 |
| 0.10 | 82,234 | 0.23 |
| 0.15 | 81,962 | 0.31 |
| 0.20 | 82,123 | 0.42 |

**Óptimo:** 0.10-0.15 (adaptativa)


---

## 📖 Ejemplo Didáctico: Algoritmo Genético

> [!NOTE]
> **Este es un ejemplo didáctico simplificado** con un problema pequeño (10 clases, 4 salones, 4 profesores) para demostrar el funcionamiento del algoritmo genético de manera clara y comprensible.

**Problema Simplificado para Demostración**

---

### Problema Didáctico

**Datos:**
- 10 clases
- 4 salones
- 4 profesores
- Objetivo: Evolucionar hacia soluciones óptimas

**Archivo:** `ejemplos_didacticos/03_algoritmo_genetico.py`

**Propósito:** Demostrar evolución de población con operadores genéticos

---

### Operadores Genéticos

**Crear Individuo (Cromosoma):**
\begin{lstlisting}[language=Python]
def crear_individuo():
    individuo = []
    for clase in clases:
        compatibles = [i for i, s in enumerate(salones) 
                      if s.tipo == clase.tipo]
        individuo.append(random.choice(compatibles))
    return individuo
\end{verbatim}


**Cruce de Un Punto:**
\begin{lstlisting}[language=Python]
def cruce(p1, p2):
    punto = len(p1) // 2
    hijo1 = p1[:punto] + p2[punto:]
    hijo2 = p2[:punto] + p1[punto:]
    return hijo1, hijo2
\end{verbatim}


---

### Operadores Genéticos (cont.)

**Mutación:**
\begin{lstlisting}[language=Python]
def mutacion(ind, prob=0.1):
    mutado = ind.copy()
    for i in range(len(mutado)):
        if random.random() < prob:
            compatibles = [j for j, s in enumerate(salones) 
                          if s.tipo == clases[i].tipo]
            mutado[i] = random.choice(compatibles)
    return mutado
\end{verbatim}


---

### Parámetros del Ejemplo

| Parámetro | Valor |
|-----------|-------|
| **Población** | 30 individuos |
| **Generaciones** | 100 |
| **Prob. cruce** | 0.8 |
| **Prob. mutación** | 0.1 |
| **Elitismo** | 2 mejores |

---

### Resultados del Ejemplo

**Fitness Inicial (Aleatorio):**
\begin{verbatim}
Fitness: 0.0196
\end{verbatim}


**Fitness Final (Generación 100):**
\begin{verbatim}
Fitness: 0.0476
\end{verbatim}


**Mejora Lograda:** 142.9% ✅

**Evolución:**
\begin{verbatim}
Gen 0:  Fitness = 0.0476
Gen 50: Fitness = 0.0476  
Gen 99: Fitness = 0.0476
\end{verbatim}


**Conclusión:** El AG explora el espacio y encuentra mejores soluciones

---

### Ejecutar el Ejemplo

\begin{lstlisting}[language=Bash]
cd ejemplos_didacticos
python3 03_algoritmo_genetico.py
\end{verbatim}


**Salida esperada:**
\begin{verbatim}
============================================================
EJEMPLO: Algoritmo Genético
============================================================

🎲 Solución aleatoria inicial:
Fitness aleatorio: 0.0196

🧬 Ejecutando Algoritmo Genético...
Gen 0: Fitness = 0.0476
Gen 10: Fitness = 0.0476
...
Gen 99: Fitness = 0.0476
\end{verbatim}


---

### Resultados del Ejemplo

\begin{verbatim}
🏆 Mejor solución encontrada:
  C1 → S3
  C2 → S3
  C3 → S1
  C4 → S2
  C5 → S4
  ...

Fitness final: 0.0476
Mejora sobre aleatorio: 142.9%

✅ Completado!
\end{verbatim}


---

## 10. Ventajas y Limitaciones


---

### 10.1 Ventajas

✅ **Exploración global:** Mantiene población diversa  
✅ **Robustez:** No requiere información de gradiente  
✅ **Paralelizable:** Evaluaciones independientes  
✅ **Flexibilidad:** Fácil incorporar nuevas restricciones  
✅ **Escape de óptimos locales:** Mediante mutación  
✅ **Soluciones múltiples:** Población final contiene varias buenas soluciones  


---

### 10.2 Limitaciones

❌ **Lento:** Requiere muchas evaluaciones  
❌ **Parámetros:** Sensible a configuración  
❌ **Convergencia prematura:** Puede perder diversidad  
❌ **No determinista:** Resultados varían entre ejecuciones  
❌ **Escalabilidad:** Costo crece con tamaño de población  


---

### 10.3 Cuándo Usar Algoritmo Genético

**Recomendado cuando:**
- Se requiere exploración exhaustiva
- Hay tiempo suficiente (>1 minuto)
- Problema tiene muchos óptimos locales
- Se necesitan múltiples soluciones alternativas

**No recomendado cuando:**
- Se requiere velocidad (<30s)
- Problema es convexo o unimodal
- Hay buenos algoritmos específicos disponibles


---

## 11. Extensiones y Mejoras


---

### 11.1 Algoritmos Genéticos Paralelos

**Modelo de Islas:**

\begin{verbatim}
Población dividida en subpoblaciones (islas)
Cada isla evoluciona independientemente
Migración periódica de mejores individuos
\end{verbatim}


**Speedup teórico:**

\end{equation}

\begin{equation}

Para $k$ procesadores.


---

### 11.2 Algoritmos Meméticos

Combinar AG con búsqueda local:

\begin{verbatim}
Para cada individuo en población:
    Aplicar Hill Climbing local
\end{verbatim}


**Ventaja:** Combina exploración global (AG) con explotación local (HC)


---

### 11.3 Coevolución

Evolucionar simultáneamente:
- Población de soluciones
- Población de restricciones/pesos

**Objetivo:** Encontrar configuración de parámetros óptima automáticamente


---

## 12. Conclusiones

El Algoritmo Genético ofrece:
- **Mejor exploración** del espacio de búsqueda
- **Múltiples soluciones** de calidad
- **Robustez** ante cambios en el problema

A costa de:
- **Mayor tiempo** de ejecución
- **Variabilidad** en resultados
- **Complejidad** de configuración

Es la opción recomendada cuando se dispone de tiempo y se requiere la mejor solución posible.


---

## Referencias

1. Holland, J. H. (1975). *Adaptation in Natural and Artificial Systems*. University of Michigan Press.

2. Goldberg, D. E. (1989). *Genetic Algorithms in Search, Optimization, and Machine Learning*. Addison-Wesley.

3. Mitchell, M. (1998). *An Introduction to Genetic Algorithms*. MIT Press.

4. Eiben, A. E., & Smith, J. E. (2015). *Introduction to Evolutionary Computing* (2nd ed.). Springer.

5. Deb, K. (2001). *Multi-Objective Optimization Using Evolutionary Algorithms*. Wiley.

6. Michalewicz, Z. (1996). *Genetic Algorithms + Data Structures = Evolution Programs*. Springer.


---

<!-- _class: lead blue -->
# Parámetros y Configuración

**Ajuste y Optimización de Algoritmos**

---

## ⚙️ Parámetros por Algoritmo

Cada algoritmo tiene parámetros críticos que afectan su rendimiento:

| Algoritmo | Parámetros Principales | Valores Óptimos |
|-----------|------------------------|--------------------|
| **Greedy+HC** | Pesos, max_iteraciones | w=10, iter=1000 |
| **ML** | n_estimators, max_depth | 100, 20 |
| **Genético** | Población, prob_mutación | 100, 0.1 |

**Documentación completa:** [`PARAMETROS.md`](../../PARAMETROS.md)

---

## 🔨 Greedy + Hill Climbing - Parámetros

### Pesos de la Función Objetivo

| Componente | Peso | Justificación |
|------------|------|---------------|
| **Movimientos** | 10.0 | Objetivo principal |
| **Cambios de piso** | 5.0 | Importante secundario |
| **Distancia** | 1.0 | Refinamiento fino |
| **Penalización P2** | 50.0 | Soft constraint alta |
| **Penalización P3** | 25.0 | Soft constraint baja |

---

## 🤖 Machine Learning - Parámetros

### Random Forest

| Parámetro | Valor | Rango Probado | Resultado |
|-----------|-------|---------------|-----------|
| **n_estimators** | 100 | [50, 500] | 94% precisión |
| **max_depth** | 20 | [10, None] | Evita overfitting |
| **min_samples_split** | 5 | [2, 20] | Balance |

---

## 🧬 Algoritmo Genético - Parámetros

### Población y Evolución

| Parámetro | Valor | Justificación |
|-----------|-------|---------------|
| **Población** | 100 | Balance diversidad/tiempo |
| **Generaciones** | 200 | Convergencia completa |
| **Elitismo** | 5 | Preserva mejores |
| **Prob. cruce** | 0.8 | Alta exploración |
| **Prob. mutación** | 0.1 | Balance |

---

## 📊 Análisis de Sensibilidad

### Experimento 1: Pesos (Greedy+HC)

| w_movimientos | Movimientos | Energía Total |
|---------------|-------------|---------------|
| 5.0 | 320 | 5780 |
| **10.0** | **314** | **5181** ✅ |
| 15.0 | 312 | 6045 |
| 20.0 | 310 | 7285 |

**Conclusión:** w=10.0 es óptimo

---

## 📊 Resumen de Parámetros

### Configuración Óptima Final

\begin{lstlisting}[language=Python]
# Greedy + Hill Climbing
PESOS = {
    'movimientos': 10.0,
    'cambios_piso': 5.0,
    'distancia': 1.0
}
MAX_ITER = 1000

# Machine Learning
N_ESTIMATORS = 100
MAX_DEPTH = 20

# Algoritmo Genético
POBLACION = 100
GENERACIONES = 200
PROB_CRUCE = 0.8
PROB_MUTACION = 0.1
\end{verbatim}


---

<!-- _class: lead blue -->
# Post-procesamiento y Corrección

**Validación y Refinamiento de Soluciones**

---

## 1. Introducción y Necesidad


---

### 1.1 Motivación

Aunque los optimizadores están diseñados para respetar los índices inmutables, pueden ocurrir violaciones accidentales debido a:

1. **Errores de implementación:** Bugs en el código de protección
2. **Condiciones de carrera:** En implementaciones paralelas
3. **Operadores complejos:** Mutaciones o cruces que no verifican correctamente
4. **Datos corruptos:** Archivos de entrada modificados manualmente

**Solución:** Sistema de corrección post-optimización que **garantiza** 100% P1.


---

### 1.2 Principio de Defensa en Profundidad

\begin{verbatim}
Capa 1: Pre-asignación forzada
    ↓
Capa 2: Protección durante optimización
    ↓
Capa 3: Corrección post-optimización ← Este documento
    ↓
Garantía: 100% PRIORIDAD 1
\end{verbatim}


**Teorema de Corrección:**
*Si existe al menos una asignación válida para cada clase P1, el sistema de corrección garantiza 100% de cumplimiento.*


---

## 2. Arquitectura del Sistema


---

### 2.1 Flujo de Corrección

\begin{verbatim}
graph TD
    A[Horario Optimizado] --> B[Cargar Preferencias]
    B --> C[Identificar Violaciones]
    C --> D{Hay violaciones?}
    D -->|No| E[Retornar sin cambios]
    D -->|Sí| F[Ordenar por Prioridad]
    F --> G[Corregir una por una]
    G --> H[Resolver Conflictos]
    H --> I[Verificar 100%]
    I --> J{Cumple 100%?}
    J -->|Sí| K[Guardar Corregido]
    J -->|No| L[Reportar Error]
\end{verbatim}



---

### 2.2 Componentes Principales

\begin{lstlisting}[language=Python]
class CorrectorPrioridades:
    def __init__(self):
        self.preferencias = {}
        self.correcciones_aplicadas = 0
        self.conflictos_resueltos = 0
    
    def corregir(self, archivo_horario):
        # 1. Cargar datos
        df = pd.read_csv(archivo_horario)
        self.cargar_preferencias()
        
        # 2. Identificar violaciones
        violaciones = self.identificar_violaciones(df)
        
        if not violaciones:
            print("✅ No hay violaciones")
            return df
        
        # 3. Corregir
        df_corregido = self.aplicar_correcciones(df, violaciones)
        
        # 4. Verificar
        cumplimiento = self.verificar_cumplimiento(df_corregido)
        
        if cumplimiento < 100.0:
            raise Exception(f"❌ Cumplimiento {cumplimiento}% < 100%")
        
        # 5. Guardar
        df_corregido.to_csv(archivo_horario, index=False)
        
        return df_corregido
\end{verbatim}



---

## 3. Identificación de Violaciones


---

### 3.1 Algoritmo de Detección

\begin{lstlisting}[language=Python]
def identificar_violaciones(self, df):
    """
    Identifica todas las clases P1 que no están en su salón preferido
    """
    violaciones = []
    
    for idx, clase in df.iterrows():
        profesor = clase['Profesor']
        materia = clase['Materia']
        tipo = clase['Tipo_Salon']
        salon_actual = clase['Salon']
\end{verbatim}


---

### 3.1 Algoritmo de Detección (continuación 1)

\begin{lstlisting}[language=Python]
        # Verificar si tiene preferencia prioritaria
        if profesor not in self.preferencias:
            continue
        
        if materia not in self.preferencias[profesor]['materias']:
            continue
        
        pref = self.preferencias[profesor]['materias'][materia]
        
        # Verificar teoría
        if tipo == 'Teoría':
            if (pref.get('prioridad_teoria') == 'Prioritario' and
                pref.get('salon_teoria') != 'Sin preferencia'):
                
                salon_esperado = pref['salon_teoria']
                
                if salon_actual != salon_esperado:
                    violaciones.append({
                        'idx': idx,
                        'clase': clase,
                        'salon_actual': salon_actual,
                        'salon_esperado': salon_esperado,
                        'profesor': profesor,
                        'materia': materia,
                        'tipo': 'Teoría'
                    })
\end{verbatim}


---

### 3.1 Algoritmo de Detección (continuación 2)

\begin{lstlisting}[language=Python]
        # Verificar laboratorio
        elif tipo == 'Laboratorio':
            if (pref.get('prioridad_lab') == 'Prioritario' and
                pref.get('salon_lab') != 'Sin preferencia'):
                
                salon_esperado = pref['salon_lab']
                
                if salon_actual != salon_esperado:
                    violaciones.append({
                        'idx': idx,
                        'clase': clase,
                        'salon_actual': salon_actual,
                        'salon_esperado': salon_esperado,
                        'profesor': profesor,
                        'materia': materia,
                        'tipo': 'Laboratorio'
                    })
    
    return violaciones
\end{verbatim}



---

### 3.2 Clasificación de Violaciones

**Tipo 1: Violación Simple**
- Clase P1 en salón incorrecto
- Salón preferido está libre
- **Solución:** Cambio directo

---

**Tipo 2: Violación con Conflicto**
- Clase P1 en salón incorrecto
- Salón preferido ocupado por clase no-P1
- **Solución:** Desplazar ocupante

**Tipo 3: Violación Compleja**
- Clase P1 en salón incorrecto
- Salón preferido ocupado por otra clase P1
- **Solución:** Requiere intervención manual o re-optimización


---

## 4. Aplicación de Correcciones


---

### 4.1 Algoritmo Principal

\begin{lstlisting}[language=Python]
def aplicar_correcciones(self, df, violaciones):
    """
    Aplica correcciones para todas las violaciones
    """
    df_corregido = df.copy()
    ocupacion = self.construir_mapa_ocupacion(df_corregido)
    
    # Ordenar violaciones por prioridad
    violaciones_ordenadas = self.ordenar_violaciones(violaciones)
    
    for violacion in violaciones_ordenadas:
        idx = violacion['idx']
        salon_esperado = violacion['salon_esperado']
        clase = violacion['clase']
        
        # Construir clave de ocupación
        key = (clase['Dia'], clase['Bloque_Horario'], 
               salon_esperado)
\end{verbatim}


---

### 4.1 Algoritmo Principal (cont.)

\begin{lstlisting}[language=Python]
        # Verificar si salón está libre
        if key not in ocupacion:
            # Corrección simple
            df_corregido.loc[idx, 'Salon'] = salon_esperado
            ocupacion[key] = idx
            self.correcciones_aplicadas += 1
        else:
            # Resolver conflicto
            exito = self.resolver_conflicto_correccion(
                df_corregido, idx, salon_esperado, 
                ocupacion, violacion
            )
            
            if exito:
                self.correcciones_aplicadas += 1
                self.conflictos_resueltos += 1
            else:
                print(f"❌ No se pudo corregir índice {idx}")
    
    return df_corregido
\end{verbatim}



---

### 4.2 Ordenamiento de Violaciones

\begin{lstlisting}[language=Python]
def ordenar_violaciones(self, violaciones):
    """
    Ordena violaciones por prioridad de corrección
    """
    def prioridad_correccion(v):
        # Criterios (mayor valor = mayor prioridad):
        # 1. Número de clases del profesor
        num_clases = sum(1 for vv in violaciones 
                        if vv['profesor'] == v['profesor'])
        
        # 2. Tipo (Teoría > Laboratorio)
        tipo_peso = 2 if v['tipo'] == 'Teoría' else 1
        
        # 3. Complejidad de resolución
        complejidad = self.estimar_complejidad(v)
        
        return (num_clases, tipo_peso, -complejidad)
    
    return sorted(violaciones, key=prioridad_correccion, reverse=True)
\end{verbatim}



---

## 5. Resolución de Conflictos


---

### 5.1 Estrategia de Desplazamiento

Cuando el salón preferido está ocupado:

\begin{lstlisting}[language=Python]
def resolver_conflicto_correccion(self, df, idx_p1, salon_p1, ocupacion, violacion):
    """
    Resuelve conflicto desplazando clase ocupante
    """
    clase_p1 = violacion['clase']
    key = (clase_p1['Dia'], clase_p1['Bloque_Horario'], salon_p1)
    
    idx_ocupante = ocupacion[key]
    clase_ocupante = df.iloc[idx_ocupante]
    
    # Verificar si ocupante también es P1
    if self.es_prioritaria(clase_ocupante):
        # Conflicto entre dos P1: no se puede resolver automáticamente
        print(f"⚠️  Conflicto entre dos P1: {idx_p1} y {idx_ocupante}")
        return False
\end{verbatim}


---

### 5.1 Estrategia de Desplazamiento (continuación 1)

\begin{lstlisting}[language=Python]
    # Buscar salón alternativo para ocupante
    salones_alternativos = self.obtener_salones_validos(clase_ocupante)
    
    for salon_alt in salones_alternativos:
        key_alt = (clase_ocupante['Dia'], 
                   clase_ocupante['Bloque_Horario'], 
                   salon_alt)
        
        if key_alt not in ocupacion:
            # Desplazar ocupante
            df.loc[idx_ocupante, 'Salon'] = salon_alt
            ocupacion[key_alt] = idx_ocupante
\end{verbatim}


---

### 5.1 Estrategia de Desplazamiento (continuación 2)

\begin{lstlisting}[language=Python]
            # Asignar P1 a su salón preferido
            df.loc[idx_p1, 'Salon'] = salon_p1
            ocupacion[key] = idx_p1
            
            return True
    
    # No se encontró salón alternativo
    return False
\end{verbatim}



---

### 5.2 Desplazamiento en Cadena

Si el salón alternativo también está ocupado:

\begin{lstlisting}[language=Python]
def desplazar_en_cadena(self, df, idx_inicial, 
                        salon_objetivo, ocupacion, 
                        profundidad=0):
    """
    Desplaza clases en cadena hasta liberar salón objetivo
    """
    MAX_PROFUNDIDAD = 10
    
    if profundidad > MAX_PROFUNDIDAD:
        return False
    
    clase = df.iloc[idx_inicial]
    key = (clase['Dia'], clase['Bloque_Horario'], salon_objetivo)
\end{verbatim}


---

### 5.2 Desplazamiento en Cadena (cont.)

\begin{lstlisting}[language=Python]
    # Si el salón está libre, asignar directamente
    if key not in ocupacion:
        df.at[idx_inicial, 'Salon'] = salon_objetivo
        ocupacion[key] = idx_inicial # Actualizar ocupación
        return True
    
    # Salón ocupado: desplazar ocupante primero
    idx_ocupante = ocupacion[key]
    clase_ocupante = df.iloc[idx_ocupante]
    
    # No desplazar clases P1
    if self.es_prioritaria(clase_ocupante):
        return False
    
    # Buscar salón para ocupante
    salones_alternativos = self.obtener_salones_validos(clase_ocupante)
    
    for salon_alt in salones_alternativos:
        # Intentar desplazar ocupante recursivamente
            # Ocupante desplazado exitosamente
            df.loc[idx_inicial, 'Salon'] = salon_objetivo
            ocupacion[key] = idx_inicial
            return True
    
    return False
\end{verbatim}

---

**Teorema de Desplazamiento en Cadena:**
*Si existe una cadena de desplazamientos de longitud finita que libera el salón objetivo, el algoritmo la encontrará.*

**Demostración:**
Por inducción en la profundidad:
- **Caso base (d=0):** Salón libre, asignación directa
- **Paso inductivo:** Si existe cadena de longitud $d+1$, el algoritmo explora recursivamente hasta encontrarla
- **Terminación:** Profundidad máxima evita ciclos infinitos


---

## 6. Verificación de Cumplimiento


---

### 6.1 Cálculo de Cumplimiento

\begin{lstlisting}[language=Python]
def verificar_cumplimiento(self, df):
    """
    Calcula porcentaje de cumplimiento de PRIORIDAD 1
    """
    total_p1 = 0
    cumplidas = 0
    
    for idx, clase in df.iterrows():
        profesor = clase['Profesor']
        materia = clase['Materia']
        tipo = clase['Tipo_Salon']
        salon_actual = clase['Salon']
        
        # Verificar si es P1
        if not self.es_prioritaria_con_salon(clase):
            continue
        
        total_p1 += 1
        salon_esperado = self.obtener_salon_preferido(clase)
        
        if salon_actual == salon_esperado:
            cumplidas += 1
    
    if total_p1 == 0:
        return 100.0
    
    return (cumplidas / total_p1) * 100.0
\end{verbatim}



---

### 6.2 Reporte Detallado

\begin{lstlisting}[language=Python]
def generar_reporte(self, df, violaciones_iniciales):
    """
    Genera reporte detallado de correcciones
    """
    cumplimiento_final = self.verificar_cumplimiento(df)
    
    reporte = f"""
================================================================================
📊 RESUMEN DE CORRECCIÓN
================================================================================
Total clases prioritarias: {len(violaciones_iniciales) + self.correcciones_aplicadas}
Violaciones encontradas: {len(violaciones_iniciales)}
Correcciones aplicadas: {self.correcciones_aplicadas}
Conflictos resueltos: {self.conflictos_resueltos}
Cumplimiento final: {cumplimiento_final:.1f}%
================================================================================
"""
    
    if cumplimiento_final == 100.0:
        reporte += "\n🎉 ¡Prioridades corregidas exitosamente!\n"
    else:
        reporte += f"\n⚠️  Cumplimiento {cumplimiento_final}% < 100%\n"
    
    return reporte
\end{verbatim}



---

## 7. Casos Especiales


---

### 7.1 Conflictos Irresolvables

Cuando dos clases P1 quieren el mismo salón al mismo tiempo:

\begin{lstlisting}[language=Python]
def manejar_conflicto_irresolvable(self, idx1, idx2, salon, df):
    """
    Maneja conflicto entre dos clases P1
    """
    clase1 = df.iloc[idx1]
    clase2 = df.iloc[idx2]
    
    mensaje = f"""
❌ CONFLICTO IRRESOLVABLE
─────────────────────────
Clase 1: {clase1['Materia']} ({clase1['Grupo']}) - {clase1['Profesor']}
Clase 2: {clase2['Materia']} ({clase2['Grupo']}) - {clase2['Profesor']}
Horario: {clase1['Dia']} {clase1['Bloque_Horario']}
Salón: {salon}

ACCIÓN REQUERIDA:
1. Contactar a profesores involucrados
2. Negociar cambio de horario o salón alternativo
3. Actualizar preferencias en configuración
4. Re-ejecutar optimización

\end{verbatim}


---

### 7.1 Conflictos Irresolvables (continuación)

\begin{lstlisting}[language=Python]
    print(mensaje)
    
    # Guardar en log
    with open('conflictos_irresolvables.log', 'a') as f:
        f.write(mensaje + '\n')
\end{verbatim}



---

### 7.2 Salones Insuficientes

Si no hay salones alternativos para desplazar:

\begin{lstlisting}[language=Python]
def manejar_salones_insuficientes(self, clase, df):
    """
    Maneja caso de salones insuficientes
    """
    mensaje = f"""
⚠️  SALONES INSUFICIENTES
──────────────────────────
Clase: {clase['Materia']} ({clase['Grupo']})
Tipo: {clase['Tipo_Salon']}
Horario: {clase['Dia']} {clase['Bloque_Horario']}
\end{verbatim}


---

### 7.2 Salones Insuficientes (continuación)

\begin{lstlisting}[language=Python]
SUGERENCIAS:
1. Verificar disponibilidad de salones del tipo requerido
2. Considerar usar salones de otro tipo (si aplicable)
3. Ajustar horarios para distribuir mejor la carga
4. Aumentar número de salones disponibles
"""
    
    print(mensaje)
\end{verbatim}



---

## 8. Métricas y Estadísticas


---

### 8.1 Métricas de Corrección

\begin{verbatim}
Ejecución típica:
├── Violaciones encontradas: 3-44
├── Correcciones simples: 60-70%
├── Conflictos resueltos: 30-40%
├── Tiempo: 0.1-0.3s
└── Cumplimiento final: 100%
\end{verbatim}



---

### 8.2 Distribución de Violaciones

\begin{verbatim}
Por optimizador:
├── Greedy: 3 violaciones (0.4%)
├── ML: 44 violaciones (5.0%)
└── Genético: 69 violaciones (7.8%)

Por tipo:
├── Teoría: 65%
└── Laboratorio: 35%

Por causa:
├── Operadores no protegidos: 70%
├── Inicialización aleatoria: 20%
└── Bugs de implementación: 10%
\end{verbatim}



---

## 9. Complejidad Computacional


---

### 9.1 Análisis Temporal

**Identificación:** $O(n)$

**Corrección simple:** $O(v)$ donde $v$ = violaciones

**Resolución de conflictos:** $O(v \cdot  m \cdot  d)$ donde:
- $m$ = salones alternativos promedio
- $d$ = profundidad máxima de desplazamiento

**Total:** $O(n + v \cdot  m \cdot  d)$

Para $n = 680$, $v \approx 50$, $m \approx 5$, $d = 10$:
\end{equation}

T \approx 680 + 50 \cdot 5 \cdot 10 = 3,180 \text{ operaciones}

\begin{equation}

**Tiempo real:** ~0.2 segundos


---

### 9.2 Análisis Espacial

\end{equation}

S = O(n + v)

\begin{equation}

**Memoria:** ~10 MB


---

## 10. Integración con Pipeline


---

### 10.1 Uso en ejecutar_todos.py

\begin{lstlisting}[language=Python]
# En ejecutar_todos.py
subprocess.run(["python3", "optimizador_greedy.py"])
subprocess.run(["python3", "corregir_prioridades.py", 
                "datos_estructurados/04_Horario_Optimizado_Greedy.csv"])

subprocess.run(["python3", "optimizador_ml.py"])
subprocess.run(["python3", "corregir_prioridades.py",
                "datos_estructurados/05_Horario_Optimizado_ML.csv"])

subprocess.run(["python3", "optimizador_genetico.py"])
subprocess.run(["python3", "corregir_prioridades.py",
                "datos_estructurados/06_Horario_Optimizado_Genetico.csv"])
\end{verbatim}



---

### 10.2 Verificación Automática

\begin{lstlisting}[language=Python]
def verificar_pipeline():
    """
    Verifica que todos los horarios cumplan 100% P1
    """
    archivos = [
        'datos_estructurados/04_Horario_Optimizado_Greedy.csv',
        'datos_estructurados/05_Horario_Optimizado_ML.csv',
        'datos_estructurados/06_Horario_Optimizado_Genetico.csv'
    ]
    
    corrector = CorrectorPrioridades()
    
    for archivo in archivos:
        df = pd.read_csv(archivo)
        cumplimiento = corrector.verificar_cumplimiento(df)
        
        assert cumplimiento == 100.0, \
            f"❌ {archivo}: {cumplimiento}% != 100%"
    
    print("✅ Todos los horarios cumplen 100% P1")
\end{verbatim}



---

## 11. Ventajas del Sistema

✅ **Garantía absoluta** de 100% P1  
✅ **Corrección automática** sin intervención manual  
✅ **Rápido** (<1 segundo)  
✅ **Robusto** ante cualquier violación  
✅ **Transparente** con reportes detallados  
✅ **Integrado** en pipeline automático  


---

## 12. Limitaciones y Consideraciones

❌ **No previene violaciones:** Solo las corrige después  
❌ **Puede empeorar P2/P3:** Al desplazar clases  
❌ **Conflictos irresolvables:** Requieren intervención manual  
❌ **Dependencia de salones:** Necesita suficientes salones alternativos  


---

## 13. Mejoras Futuras


---

### 13.1 Corrección Inteligente

Minimizar impacto en P2/P3 al desplazar:

\begin{lstlisting}[language=Python]
def desplazar_minimizando_impacto(self, clase, salones_alt, df):
    """
    Selecciona salón alternativo que minimiza impacto en P2/P3
    """
    mejor_salon = None
    menor_impacto = float('inf')
    
    for salon in salones_alt:
        impacto = self.calcular_impacto_p2_p3(clase, salon, df)
        
        if impacto < menor_impacto:
            menor_impacto = impacto
            mejor_salon = salon
    
    return mejor_salon
\end{verbatim}



---

### 13.2 Prevención Proactiva

Verificar durante optimización:

\begin{lstlisting}[language=Python]
# En optimizadores
def verificar_antes_de_aplicar(self, nueva_solucion):
    """
    Verifica P1 antes de aceptar nueva solución
    """
    violaciones = self.identificar_violaciones_p1(nueva_solucion)
    
    if violaciones:
        # Rechazar solución
        return False
    
    return True
\end{verbatim}



---

## 14. Conclusiones

El sistema de corrección post-optimización:

1. **Garantiza** 100% cumplimiento de PRIORIDAD 1
2. **Complementa** la protección durante optimización
3. **Proporciona** última línea de defensa
4. **Permite** que optimizadores se enfoquen en calidad
5. **Asegura** robustez del sistema completo

Es un componente **esencial** que hace el sistema **production-ready**.


---

## Referencias

1. Apt, K. R. (2003). *Principles of Constraint Programming*. Cambridge University Press.

2. Tsang, E. (1993). *Foundations of Constraint Satisfaction*. Academic Press.

3. Dechter, R. (2003). *Constraint Processing*. Morgan Kaufmann.


---

<!-- _class: lead blue -->
# Arquitectura del Código

**Diseño e Implementación del Sistema**

---

## 1. Visión General del Sistema


---

### 1.1 Arquitectura de Alto Nivel

\begin{verbatim}
Sistema-Salones-ISC/
│
├── Capa de Presentación (UI)
│   └── configurador_materias.py
│
├── Capa de Procesamiento
│   ├── pre_asignar_p1.py
│   ├── optimizador_greedy.py
│   ├── optimizador_ml.py
│   ├── optimizador_genetico.py
│   └── corregir_prioridades.py
│
├── Capa de Utilidades
│   └── utils_restricciones.py
│
├── Capa de Orquestación
│   ├── ejecutar_todos.py
│   └── generar_comparativa_completa.py
│
└── Capa de Datos
    ├── datos_estructurados/
    ├── preferencias_profesores.json
    ├── configuracion_materias.json
    └── asignacion_grupos_1er_semestre.json
\end{verbatim}



---

### 1.2 Flujo de Datos

\begin{verbatim}
graph LR
    A[Configuración] --> B[Horario Inicial]
    B --> C[Pre-asignación P1]
    C --> D[Optimizadores]
    D --> E[Corrección]
    E --> F[Comparativas]
    F --> G[Reportes]
\end{verbatim}



---

## 2. Módulos Principales


---

### 2.1 configurador_materias.py

**Propósito:** Interfaz gráfica para configuración del sistema

**Clase Principal:**
\begin{lstlisting}[language=Python]
class ConfiguradorMaterias:
    """
    Aplicación Tkinter para configurar materias y preferencias
    """
    def __init__(self):
        self.root = tk.Tk()
        self.materias = {}
        self.preferencias = {}
        self.grupos_1er = {}
        
        self.crear_interfaz()
\end{verbatim}


---

### 2.1 configurador_materias.py (continuación)

\begin{lstlisting}[language=Python]
    # Métodos principales
    def crear_interfaz(self)
    def crear_tab_materias(self)
    def crear_tab_preferencias(self)
    def crear_tab_grupos(self)
    def guardar_todo(self)
    def cargar_todo(self)
\end{verbatim}


**Responsabilidades:**
- Configurar características de materias (horas, distribución)
- Definir preferencias de profesores (salón + prioridad)
- Asignar grupos de primer semestre
- Guardar/cargar configuración en JSON

**Dependencias:**
- `tkinter` (UI)
- `json` (persistencia)
- `pandas` (validación)

**Archivos Generados:**
- `preferencias_profesores.json`
- `configuracion_materias.json`
- `asignacion_grupos_1er_semestre.json`


---

### 2.2 pre_asignar_p1.py

**Propósito:** Pre-asignación forzada de PRIORIDAD 1

**Funciones Principales:**
\begin{lstlisting}[language=Python]
def cargar_datos():
    """Carga horario inicial y preferencias"""
    
def identificar_clases_prioritarias(df, preferencias):
    """Identifica todas las clases P1"""
    
def asignar_forzadamente(df, clases_p1):
    """Asigna cada clase P1 a su salón preferido"""
    
def guardar_resultados(df, indices_inmutables):
    """Guarda horario pre-asignado e índices"""
\end{verbatim}

---

**Algoritmo:**
1. Cargar horario inicial
2. Identificar clases P1 desde preferencias
3. Ordenar por complejidad
4. Asignar forzadamente (resolviendo conflictos)
5. Marcar índices como inmutables
6. Guardar resultados
---

**Archivos de Entrada:**
- `datos_estructurados/01_Horario_Inicial.csv`
- `preferencias_profesores.json`

**Archivos de Salida:**
- `datos_estructurados/00_Horario_PreAsignado_P1.csv`
- `datos_estructurados/indices_inmutables_p1.json`


---

### 2.3 optimizador_greedy.py

**Propósito:** Optimización mediante Greedy + Hill Climbing

**Clase Principal:**
\begin{lstlisting}[language=Python]
class OptimizadorGreedyHC:
    def __init__(self):
        self.indices_inmutables = set()
        self.max_iter_hc = 100
        self.pesos = {...}
    
    # Fase 1: Construcción
    def construccion_greedy(self, df):
        """Construye solución inicial vorazmente"""
    
    def calcular_score(self, clase, salon, solucion):
        """Calcula score para asignación"""
\end{verbatim}


---

### 2.3 optimizador_greedy.py (continuación)

\begin{lstlisting}[language=Python]
    # Fase 2: Refinamiento
    def hill_climbing(self, solucion, df):
        """Mejora mediante búsqueda local"""
    
    def calcular_energia(self, solucion, df):
        """Calcula energía total de solución"""
    
    # Utilidades
    def analizar_movimientos(self, df):
        """Analiza movimientos de profesores"""
\end{verbatim}


---

**Flujo de Ejecución:**
1. Cargar horario pre-asignado
2. Cargar índices inmutables
3. Construcción greedy (respetando inmutables)
4. Hill Climbing (protegiendo inmutables)
5. Corrección final de P1
6. Análisis de métricas
7. Guardar resultado

---

**Archivos de Entrada:**
- `datos_estructurados/00_Horario_PreAsignado_P1.csv`
- `datos_estructurados/indices_inmutables_p1.json`
- `preferencias_profesores.json`
- `configuracion_materias.json`

**Archivos de Salida:**
- `datos_estructurados/04_Horario_Optimizado_Greedy.csv`
- `comparativas/04_inicial_vs_greedy/metricas_movimientos.csv`


---

### 2.4 optimizador_ml.py

**Propósito:** Optimización mediante Machine Learning

**Clase Principal:**
\begin{lstlisting}[language=Python]
class OptimizadorML:
    def __init__(self):
        self.clasificador = RandomForestClassifier(...)
        self.regressor_calidad = GradientBoostingRegressor(...)
        self.indices_inmutables = set()
    
    # Entrenamiento
    def entrenar(self, df_inicial):
        """Entrena modelos con horario inicial"""
    
    def extraer_features(self, clase, df, idx):
        """Extrae vector de características"""
\end{verbatim}


---

### 2.4 optimizador_ml.py (continuación)

\begin{lstlisting}[language=Python]
    # Optimización
    def optimizar(self, df_inicial):
        """Genera nueva asignación usando modelos"""
    
    def predecir_salon(self, features):
        """Predice mejor salón para clase"""
\end{verbatim}


---

**Flujo de Ejecución:**
1. Cargar horario pre-asignado
2. Cargar índices inmutables
3. Entrenar Random Forest (predicción de salón)
4. Entrenar Gradient Boosting (calidad de asignación)
5. Optimizar (predecir salón para cada clase)
6. Proteger clases inmutables
7. Análisis de métricas
8. Guardar resultado

---

**Archivos de Entrada:**
- `datos_estructurados/00_Horario_PreAsignado_P1.csv`
- `datos_estructurados/indices_inmutables_p1.json`

**Archivos de Salida:**
- `datos_estructurados/05_Horario_Optimizado_ML.csv`
- `comparativas/02_inicial_vs_ml/metricas_movimientos.csv`


---

### 2.5 optimizador_genetico.py

**Propósito:** Optimización mediante Algoritmo Genético

**Clase Principal:**
\begin{lstlisting}[language=Python]
class OptimizadorGenetico:
    def __init__(self):
        self.tam_poblacion = 150
        self.num_generaciones = 500
        self.prob_cruce = 0.8
        self.prob_mutacion = 0.1
        self.indices_inmutables = set()
    
    # Inicialización
    def generar_poblacion_inicial(self, df):
        """Genera población inicial"""
\end{verbatim}


---

### 2.5 optimizador_genetico.py (continuación)

\begin{lstlisting}[language=Python]
    # Operadores
    def seleccion_torneo(self, poblacion, fitness):
        """Selecciona padres por torneo"""
    
    def cruce_un_punto(self, padre1, padre2):
        """Cruza dos individuos"""
    
    def mutacion_intercambio(self, individuo):
        """Muta individuo"""
    
    # Evolución
    def evolucionar(self, df):
        """Ejecuta algoritmo genético"""
\end{verbatim}


---

**Flujo de Ejecución:**
1. Cargar horario pre-asignado
2. Cargar índices inmutables
3. Generar población inicial

---

**Flujo de Ejecución (continuación):**
4. Evolucionar durante N generaciones:
   - Selección por torneo
   - Cruce (respetando inmutables)
   - Mutación (respetando inmutables)
   - Evaluación de fitness
   - Reemplazo con elitismo
5. Retornar mejor individuo
6. Análisis de métricas
7. Guardar resultado

---

**Archivos de Entrada:**
- `datos_estructurados/00_Horario_PreAsignado_P1.csv`
- `datos_estructurados/indices_inmutables_p1.json`

**Archivos de Salida:**
- `datos_estructurados/06_Horario_Optimizado_Genetico.csv`
- `comparativas/03_inicial_vs_genetico/metricas_movimientos.csv`


---

### 2.6 corregir_prioridades.py

**Propósito:** Corrección post-optimización de PRIORIDAD 1

**Funciones Principales:**
\begin{lstlisting}[language=Python]
def cargar_preferencias():
    """Carga preferencias de profesores"""

def identificar_violaciones(df, preferencias):
    """Identifica clases P1 en salón incorrecto"""

def corregir_violaciones(df, violaciones):
    """Corrige todas las violaciones"""

def verificar_cumplimiento(df, preferencias):
    """Verifica 100% cumplimiento"""
\end{verbatim}


---

**Flujo de Ejecución:**
1. Cargar horario optimizado
2. Cargar preferencias
3. Identificar violaciones de P1
4. Si hay violaciones:
   - Ordenar por prioridad
   - Corregir una por una
   - Resolver conflictos
5. Verificar 100% cumplimiento
6. Guardar horario corregido

**Uso:**
\begin{lstlisting}[language=Bash]
python3 corregir_prioridades.py datos_estructurados/04_Horario_Optimizado_Greedy.csv
\end{verbatim}



---

### 2.7 utils_restricciones.py

**Propósito:** Funciones de utilidad para validación y restricciones

**Funciones Principales:**
\begin{lstlisting}[language=Python]
def validar_asignacion(clase, salon, df):
    """Verifica si asignación es válida"""

def obtener_salones_validos(clase, salones_teoria, laboratorios):
    """Retorna salones válidos para clase"""

def calcular_distancia(salon1, salon2):
    """Calcula distancia entre salones"""

def verificar_conflicto_temporal(clase1, clase2, asignacion):
    """Verifica si hay conflicto temporal"""

def pre_asignar_prioritarias(df, config, preferencias, ...):
    """Pre-asigna clases prioritarias (legacy)"""
\end{verbatim}


---

**Uso:**
- Importado por todos los optimizadores
- Proporciona funciones comunes de validación
- Mantiene lógica de restricciones centralizada


---

### 2.8 ejecutar_todos.py

**Propósito:** Script maestro de orquestación

**Flujo:**
\begin{lstlisting}[language=Python]
def main():
    # 1. Pre-asignación
    ejecutar("pre_asignar_p1.py")
    
    # 2. Greedy
    ejecutar("optimizador_greedy.py")
    ejecutar("corregir_prioridades.py", "04_Horario_Optimizado_Greedy.csv")
    
    # 3. ML
    ejecutar("optimizador_ml.py")
    ejecutar("corregir_prioridades.py", "05_Horario_Optimizado_ML.csv")
\end{verbatim}


---

### 2.8 ejecutar_todos.py (continuación)

\begin{lstlisting}[language=Python]
    # 4. Genético
    ejecutar("optimizador_genetico.py")
    ejecutar("corregir_prioridades.py", "06_Horario_Optimizado_Genetico.csv")
    
    # 5. Comparativas
    ejecutar("generar_comparativa_completa.py")
    
    # 6. Resumen
    mostrar_resumen()
\end{verbatim}


---

**Características:**
- Ejecución secuencial de todo el pipeline
- Medición de tiempos
- Logging de salidas
- Manejo de errores


---

### 2.9 generar_comparativa_completa.py

**Propósito:** Generación de reportes y gráficos

**Funciones Principales:**
\begin{lstlisting}[language=Python]
def generar_excel_formato(csv_file, output_file):
    """Genera Excel con formato bonito"""

def generar_excel_comparativo(dfs):
    """Genera Excel con todos los optimizadores"""

def generar_graficos(dfs):
    """Genera gráficos de análisis"""

def verificar_prioridad_1(dfs, preferencias):
    """Verifica cumplimiento P1"""
\end{verbatim}


---

**Salidas Generadas:**
- Excels formateados por optimizador
- Excel comparativo (todos juntos)
- Gráficos de tiempos, cumplimiento, métricas
- Excel consolidado con resumen


---

## 3. Estructuras de Datos


---

### 3.1 Horario (DataFrame)

\begin{lstlisting}[language=Python]
# Estructura de datos principal
df = pd.DataFrame({
    'Dia': str,              # Lunes, Martes, ...
    'Bloque_Horario': int,   # 0700, 0800, ...
    'Materia': str,          # Nombre de materia
    'Grupo': str,            # 1527A, 2514/B, ...
    'Profesor': str,         # PROFESOR 3, ...
    'Salon': str,            # FF1, LBD, ...
    'Es_Invalido': int,      # 0 o 1
    'Tipo_Salon': str,       # Teoría, Laboratorio
    'Piso': int              # 0 o 1
})
\end{verbatim}



---

### 3.2 Preferencias (JSON)

\begin{verbatim}
{
  "PROFESOR 3": {
    "materias": {
      "LENGUAJES Y AUTÓMATAS I": {
        "salon_teoria": "FFA",
        "prioridad_teoria": "Prioritario",
        "salon_lab": "Sin preferencia",
        "prioridad_lab": "Normal"
      }
    }
  }
}
\end{verbatim}



---

### 3.3 Índices Inmutables (JSON)

\begin{verbatim}
{
  "indices": [12, 45, 67, ...],
  "total": 88,
  "timestamp": "2025-12-21T11:00:00",
  "version": "1.0"
}
\end{verbatim}



---

### 3.4 Solución (Dict)

\begin{lstlisting}[language=Python]
# Representación interna de solución
solucion = {
    0: 'FF1',    # Clase 0 → Salón FF1
    1: 'LBD',    # Clase 1 → Salón LBD
    2: 'FF2',    # Clase 2 → Salón FF2
    ...
}
\end{verbatim}



---

## 4. Patrones de Diseño


---

### 4.1 Strategy Pattern

Diferentes optimizadores implementan la misma interfaz:

\begin{lstlisting}[language=Python]
class OptimizadorBase:
    def optimizar(self, df_inicial):
        raise NotImplementedError

class OptimizadorGreedy(OptimizadorBase):
    def optimizar(self, df_inicial):
        # Implementación Greedy
        
class OptimizadorML(OptimizadorBase):
    def optimizar(self, df_inicial):
        # Implementación ML
\end{verbatim}



---

### 4.2 Template Method

Estructura común de optimización:

\begin{lstlisting}[language=Python]
def optimizar(self, df):
    # 1. Cargar datos
    self.cargar_configuracion()
    
    # 2. Inicializar
    solucion = self.inicializar()
    
    # 3. Optimizar (método específico)
    solucion = self.algoritmo_especifico(solucion)
    
    # 4. Post-procesar
    solucion = self.post_procesar(solucion)
    
    # 5. Guardar
    self.guardar_resultado(solucion)
\end{verbatim}



---

### 4.3 Facade Pattern

`ejecutar_todos.py` proporciona interfaz simple:

\begin{lstlisting}[language=Python]
# En lugar de ejecutar 7 scripts manualmente
python3 ejecutar_todos.py  # Un solo comando
\end{verbatim}



---

## 5. Convenciones de Código


---

### 5.1 Nomenclatura

**Archivos:**
- `snake_case.py` para scripts
- `PascalCase` para clases
- Números prefijos para orden (01_, 02_, ...)

---

**Variables:**
- `snake_case` para variables y funciones
- `UPPER_CASE` para constantes
- `_private` para métodos privados

**Clases:**
- `PascalCase` para nombres de clase
- Nombres descriptivos (OptimizadorGreedy, no OG)


---

### 5.2 Documentación

**Docstrings:**
\begin{lstlisting}[language=Python]
def funcion(param1, param2):
    """
    Descripción breve de la función
    
    Args:
        param1: Descripción del parámetro 1
        param2: Descripción del parámetro 2
    
    Returns:
        Descripción del valor de retorno
    """
\end{verbatim}


**Comentarios:**
\begin{lstlisting}[language=Python]
# Comentarios explicativos para lógica compleja
# No comentar lo obvio
\end{verbatim}



---

### 5.3 Manejo de Errores

\begin{lstlisting}[language=Python]
try:
    # Operación que puede fallar
    resultado = operacion_riesgosa()
except FileNotFoundError:
    print("❌ Archivo no encontrado")
    sys.exit(1)
except Exception as e:
    print(f"❌ Error: {e}")
    raise
\end{verbatim}



---

## 6. Dependencias


---

### 6.1 Dependencias Externas

\begin{verbatim}
pandas>=1.5.0          # Manipulación de datos
openpyxl>=3.0.0        # Excel I/O
matplotlib>=3.5.0      # Gráficos
seaborn>=0.12.0        # Visualización
scikit-learn>=1.0.0    # Machine Learning
\end{verbatim}



---

### 6.2 Dependencias Estándar

\begin{lstlisting}[language=Python]
import json            # Configuración
import random          # Aleatoriedad
import sys             # Sistema
import subprocess      # Ejecución de scripts
from pathlib import Path  # Manejo de rutas
from datetime import datetime  # Timestamps
\end{verbatim}



---

## 7. Testing y Validación


---

### 7.1 Validación de Datos

\begin{lstlisting}[language=Python]
def validar_horario(df):
    """Valida estructura de horario"""
    assert 'Dia' in df.columns
    assert 'Salon' in df.columns
    assert len(df) == 680
    # ...
\end{verbatim}



---

### 7.2 Verificación de Invariantes

\begin{lstlisting}[language=Python]
def verificar_invariantes(df, indices_inmutables, preferencias):
    """Verifica invariantes del sistema"""
    # Invariante 1: P1 al 100%
    assert verificar_p1(df, preferencias) == 100.0
    
    # Invariante 2: Sin conflictos temporales
    assert not tiene_conflictos(df)
    
    # Invariante 3: Salones válidos
    assert todos_salones_validos(df)
\end{verbatim}



---

## 8. Optimizaciones de Rendimiento


---

### 8.1 Caching

\begin{lstlisting}[language=Python]
@lru_cache(maxsize=1000)
def calcular_distancia(salon1, salon2):
    """Cachea cálculos de distancia"""
    # ...
\end{verbatim}



---

### 8.2 Vectorización

\begin{lstlisting}[language=Python]
# En lugar de loops
for i in range(len(df)):
    df.loc[i, 'Piso'] = obtener_piso(df.loc[i, 'Salon'])

# Usar operaciones vectorizadas
df['Piso'] = df['Salon'].apply(obtener_piso)
\end{verbatim}



---

### 8.3 Paralelización

\begin{lstlisting}[language=Python]
# En Random Forest
RandomForestClassifier(n_jobs=-1)  # Usar todos los cores
\end{verbatim}



---

## 9. Extensibilidad


---

### 9.1 Agregar Nuevo Optimizador

1. Crear `optimizador_nuevo.py`
2. Implementar interfaz estándar:
\begin{lstlisting}[language=Python]
   class OptimizadorNuevo:
       def __init__(self):
           self.indices_inmutables = cargar_inmutables()
       
       def optimizar(self, df):
           # Implementación
\end{verbatim}

3. Agregar a `ejecutar_todos.py`
4. Agregar a `generar_comparativa_completa.py`


---

### 9.2 Agregar Nueva Restricción

1. Definir en `utils_restricciones.py`:
\begin{lstlisting}[language=Python]
   def verificar_nueva_restriccion(clase, salon):
       # Lógica de verificación
\end{verbatim}

2. Integrar en función objetivo de optimizadores
3. Actualizar documentación


---

## 10. Deployment y Producción


---

### 10.1 Instalación

\begin{lstlisting}[language=Bash]
git clone <repo>
cd Sistema-Salones-ISC
pip install -r requirements.txt
\end{verbatim}



---

### 10.2 Ejecución

\begin{lstlisting}[language=Bash]
# Configurar
python3 configurador_materias.py

# Ejecutar pipeline completo
python3 ejecutar_todos.py

# O ejecutar componentes individuales
python3 pre_asignar_p1.py
python3 optimizador_greedy.py
# ...
\end{verbatim}



---

### 10.3 Monitoreo

\begin{lstlisting}[language=Bash]
# Ver logs
tail -f ejecucion_final.log

# Verificar resultados
ls -lh datos_estructurados/
ls -lh comparativas/final/
\end{verbatim}



---

## 11. Mantenimiento


---

### 11.1 Actualización de Preferencias

1. Abrir `configurador_materias.py`
2. Modificar preferencias
3. Guardar
4. Re-ejecutar pipeline


---

### 11.2 Debugging

\begin{lstlisting}[language=Python]
# Activar modo debug
DEBUG = True

if DEBUG:
    print(f"Debug: {variable}")
    import pdb; pdb.set_trace()
\end{verbatim}



---

### 11.3 Logs

\begin{lstlisting}[language=Python]
import logging

logging.basicConfig(
    level=logging.INFO,
    format='%(asctime)s - %(levelname)s - %(message)s',
    filename='sistema.log'
)
\end{verbatim}



---

## 12. Conclusiones

El sistema está diseñado con:

✅ **Modularidad:** Componentes independientes y reutilizables  
✅ **Escalabilidad:** Fácil agregar nuevos optimizadores  
✅ **Mantenibilidad:** Código limpio y bien documentado  
✅ **Robustez:** Validación en múltiples capas  
✅ **Extensibilidad:** Patrones de diseño facilitan extensiones  

Es un sistema **production-ready** para uso real en el Instituto Tecnológico de Ciudad Madero.


---

## Referencias

1. Martin, R. C. (2008). *Clean Code: A Handbook of Agile Software Craftsmanship*. Prentice Hall.

2. Gamma, E., et al. (1994). *Design Patterns: Elements of Reusable Object-Oriented Software*. Addison-Wesley.

3. McConnell, S. (2004). *Code Complete* (2nd ed.). Microsoft Press.


---

<!-- _class: lead blue -->
# Aplicación Web (BETA)

**Interfaz Web para el Sistema**

---

## ⚠️ Estado Actual: BETA

> ⚠️ **ESTADO**: En Desarrollo - Versión BETA  
> 🚧 **NO LISTA PARA PRODUCCIÓN**


---

## Aviso Importante

La aplicación web del Sistema de Asignación de Salones se encuentra actualmente en **fase BETA de desarrollo**. Aunque funcional para demostración, **NO está lista para uso en producción** y requiere desarrollo adicional antes de su implementación institucional.


---

## Estado Actual

### ✅ Funcionalidades Implementadas

1. **Visualización de Horarios**
   - Vista de horarios por grupo
   - Vista de horarios por profesor
   - Vista de horarios por salón
   - Filtros básicos (día, semestre)

---

2. **Interfaz de Usuario**
   - Diseño responsive básico
   - Navegación entre vistas
   - Tabla de horarios
   - Exportación a PDF (básica)

---

3. **Backend Básico**
   - API REST simple
   - Carga de datos desde CSV
   - Endpoints para consultas básicas

---

### ⚠️ Limitaciones Conocidas

1. **Seguridad**
   - ❌ Sin autenticación de usuarios
   - ❌ Sin autorización por roles
   - ❌ Sin encriptación de datos sensibles
   - ❌ Vulnerable a inyección SQL (si se usa BD)

2. **Rendimiento**
   - ⚠️ No optimizado para grandes volúmenes
   - ⚠️ Sin caché de datos
   - ⚠️ Carga completa en cada request
   - ⚠️ No hay paginación

---

3. **Funcionalidad**
   - ❌ No permite edición de horarios
   - ❌ No integra con optimizadores
   - ❌ No hay sistema de notificaciones
   - ❌ Exportación limitada (solo PDF básico)

4. **Estabilidad**
   - ⚠️ Manejo de errores básico
   - ⚠️ Sin logging robusto
   - ⚠️ No hay tests automatizados
   - ⚠️ Puede fallar con datos inconsistentes


---

## Arquitectura Actual

\begin{verbatim}
web-app/ (BETA)
├── backend/
│   ├── app.py              # Flask/FastAPI server
│   ├── api/
│   │   ├── horarios.py     # Endpoints de horarios
│   │   └── consultas.py    # Endpoints de consultas
│   └── models/
│       └── horario.py      # Modelos de datos
│
├── frontend/
│   ├── index.html          # Página principal
│   ├── css/
│   │   └── styles.css      # Estilos
│   └── js/
│       ├── app.js          # Lógica principal
│       └── api.js          # Cliente API
│
└── README_WEB.md           # Este archivo
\end{verbatim}



---

## Instalación (Solo para Desarrollo)

\begin{lstlisting}[language=Bash]
# Instalar dependencias
pip install flask flask-cors pandas

# Ejecutar servidor de desarrollo
cd web-app/backend
python app.py

# Abrir en navegador
open http://localhost:5000
\end{verbatim}



---

## Uso Básico (Demo)

### Ver Horario de Grupo

\begin{verbatim}
http://localhost:5000/horario/grupo/1527A
\end{verbatim}


### Ver Horario de Profesor

\begin{verbatim}
http://localhost:5000/horario/profesor/PROFESOR%203
\end{verbatim}


### Ver Ocupación de Salón

\begin{verbatim}
http://localhost:5000/horario/salon/FFA
\end{verbatim}



---

## Roadmap de Desarrollo

### Fase 1: Seguridad (Crítico)
- [ ] Implementar autenticación (JWT)
- [ ] Sistema de roles (Admin, Profesor, Estudiante)
- [ ] Validación de entrada
- [ ] Sanitización de datos
- [ ] HTTPS obligatorio

---

### Fase 2: Funcionalidad Core
- [ ] Integración con optimizadores
- [ ] Edición de horarios (con permisos)
- [ ] Comparación de horarios
- [ ] Exportación avanzada (Excel, iCal, PDF mejorado)
- [ ] Sistema de notificaciones

---

### Fase 3: Rendimiento
- [ ] Caché de datos
- [ ] Paginación
- [ ] Lazy loading
- [ ] Optimización de queries
- [ ] CDN para assets

---

### Fase 4: UX/UI
- [ ] Diseño profesional
- [ ] Modo oscuro
- [ ] Accesibilidad (WCAG 2.1)
- [ ] PWA (Progressive Web App)
- [ ] Responsive mejorado

---

### Fase 5: Integración
- [ ] API con sistema institucional
- [ ] Single Sign-On (SSO)
- [ ] Sincronización automática
- [ ] Webhooks para actualizaciones


---

## Tecnologías Propuestas

### Backend
- **Framework**: FastAPI (recomendado) o Flask
- **Base de Datos**: PostgreSQL
- **ORM**: SQLAlchemy
- **Autenticación**: JWT + OAuth2
- **Cache**: Redis
- **API Docs**: Swagger/OpenAPI

---

### Frontend
- **Framework**: React o Vue.js
- **UI Library**: Material-UI o Ant Design
- **State Management**: Redux o Vuex
- **Build Tool**: Vite
- **Testing**: Jest + React Testing Library

---

### DevOps
- **Containerización**: Docker
- **Orquestación**: Docker Compose (dev) / Kubernetes (prod)
- **CI/CD**: GitHub Actions
- **Monitoring**: Prometheus + Grafana
- **Logging**: ELK Stack


---

## Advertencias de Seguridad

🚨 **NO USAR EN PRODUCCIÓN SIN:**

1. **Autenticación robusta**
   - Sistema de login seguro
   - Gestión de sesiones
   - Protección contra fuerza bruta

2. **Autorización por roles**
   - Permisos granulares
   - Validación en backend
   - Auditoría de acciones

---

## Advertencias de Seguridad (continuación)

3. **Validación de datos**
   - Sanitización de entrada
   - Validación de tipos
   - Protección contra XSS/CSRF

4. **Encriptación**
   - HTTPS obligatorio
   - Encriptación de datos sensibles
   - Hashing de contraseñas (bcrypt)

---

## Advertencias de Seguridad (resumen técnico)

5. **Auditoría y Logging**
   - Registro de todas las acciones
   - Monitoreo de seguridad
   - Alertas automáticas

---

## Contribuir al Desarrollo Web

Si deseas contribuir al desarrollo de la aplicación web:

1. **Revisar roadmap** y seleccionar tarea
2. **Crear branch** desde `develop`
3. **Implementar** con tests
4. **Documentar** cambios
5. **Pull request** para revisión

---

### Estándares de Código

\begin{lstlisting}[language=Python]
# Backend (Python)
- PEP 8 style guide
- Type hints obligatorios
- Docstrings para funciones públicas
- Tests unitarios (pytest)
- Coverage > 80%

# Frontend (JavaScript)
- ESLint + Prettier
- Componentes funcionales
- PropTypes o TypeScript
- Tests de componentes
- Accesibilidad (a11y)
\end{verbatim}



---

## Contacto

**Autor:** Jesús Olvera

- **GitHub:** [@jjho05](https://github.com/jjho05)
- **Email:** jjho.reivaj05@gmail.com / hernandez.jesusjavier.20.0770@gmail.com
- **Repositorio:** https://github.com/jjho05/Sistema-Salones-ISC
- **Institucional:** sistemas@cdmadero.tecnm.mx


---

## Licencia

Mismo que el proyecto principal - Uso académico TECNM.

---

**Última actualización**: 2025-12-21  
**Versión**: 0.1.0-beta  
**Estado**: 🚧 En Desarrollo - NO PRODUCCIÓN  
**Mantenedor**: Equipo de Desarrollo Web ISC


---

<!-- _class: lead blue -->
# Preguntas Frecuentes

**FAQ - Sistema de Asignación de Salones**

---

## General

### ¿Qué es el Sistema de Asignación de Salones ISC?

Es un sistema de optimización que asigna automáticamente salones a clases del programa de Ingeniería en Sistemas Computacionales del Instituto Tecnológico de Ciudad Madero, minimizando movimientos de profesores y garantizando cumplimiento de preferencias prioritarias.

---

### ¿Por qué necesitamos este sistema?

La asignación manual de 680+ clases a 21 salones es:
- **Tiempo intensiva**: Horas o días de trabajo manual
- **Propensa a errores**: Conflictos temporales, preferencias olvidadas
- **Subóptima**: No considera todas las optimizaciones posibles
- **Difícil de actualizar**: Cambios requieren rehacer todo

---

### ¿Qué problemas resuelve?

1. **Garantiza 100% cumplimiento** de preferencias prioritarias de profesores
2. **Minimiza movimientos** de profesores entre salones
3. **Reduce cambios de piso** (menos fatiga)
4. **Optimiza distancias** recorridas
5. **Genera reportes** automáticos y comparativos


---

## Instalación y Configuración

### ¿Qué necesito para instalar el sistema?

**Requisitos:**
- Python 3.8 o superior
- pip (gestor de paquetes)
- ~100 MB de espacio en disco
- Sistema operativo: Windows, macOS, o Linux

**Instalación:**
\begin{lstlisting}[language=Bash]
pip install pandas openpyxl matplotlib seaborn scikit-learn
\end{verbatim}


---

### ¿Cómo configuro las preferencias de profesores?

1. Ejecutar `python3 configurador_materias.py`
2. Ir a pestaña "Preferencias de Profesores"
3. Seleccionar profesor y materia
4. Configurar:
   - Salón preferido (Teoría/Laboratorio)
   - Prioridad (Prioritario/Normal/Sin preferencia)
5. Guardar configuración

---

### ¿Dónde se guardan las configuraciones?

- `preferencias_profesores.json`: Preferencias de profesores
- `configuracion_materias.json`: Configuración de materias
- `asignacion_grupos_1er_semestre.json`: Grupos de primer semestre


---

## Uso del Sistema

### ¿Cómo ejecuto el sistema completo?

**Opción 1 (Recomendada):**
\begin{lstlisting}[language=Bash]
python3 ejecutar_todos.py
\end{verbatim}


**Opción 2 (Manual):**
\begin{lstlisting}[language=Bash]
python3 pre_asignar_p1.py
python3 optimizador_greedy.py
python3 corregir_prioridades.py datos_estructurados/04_Horario_Optimizado_Greedy.csv
# ... etc
\end{verbatim}


---

### ¿Cuánto tiempo tarda?

| Componente | Tiempo |
|------------|--------|
| Pre-asignación | ~0.3s |
| Greedy | ~30s |
| ML | ~16s |
| Genético | ~74s |
| Corrección | ~0.2s cada uno |
| Comparativas | ~5s |
| **Total** | **~2-3 minutos** |

---

### ¿Qué optimizador debo usar?

| Optimizador | Cuándo usar |
|-------------|-------------|
| **Greedy** | Balance ideal velocidad/calidad, uso general |
| **ML** | Máxima velocidad, buena calidad |
| **Genético** | Mejor exploración, tiempo no crítico |
| **Todos** | Comparar y elegir mejor resultado |


---

## Prioridades

### ¿Qué significa PRIORIDAD 1?

**PRIORIDAD 1** son preferencias de salón de profesores marcadas como "Prioritario". El sistema **garantiza 100%** de cumplimiento.

**Ejemplo:**
- PROFESOR 3 quiere FFA para Lenguajes y Autómatas I (Teoría)
- Prioridad: Prioritario
- **Garantía**: Todas sus clases estarán en FFA

---

### ¿Qué pasa si dos profesores quieren el mismo salón al mismo tiempo?

**Durante configuración:**
- El sistema detecta el conflicto
- Solicita resolver manualmente
- Opciones: Cambiar horario o salón de uno

**Durante optimización:**
- Pre-asignación resuelve automáticamente
- Desplaza clases no-prioritarias
- Si ambas son P1: Requiere intervención manual

---

### ¿Qué son PRIORIDAD 2 y 3?

**PRIORIDAD 2**: Consistencia de grupos
- Objetivo: Mantener grupos en mismo salón
- **No garantizado**, se optimiza cuando es posible

**PRIORIDAD 3**: Grupos de primer semestre
- Objetivo: Asignar grupos 15xx a salones específicos
- **No garantizado**, mejor esfuerzo

**Estado actual**: P2 y P3 en desarrollo (v2.1.0)


---

## Resultados

### ¿Dónde encuentro los resultados?

**Horarios optimizados:**
- `datos_estructurados/04_Horario_Optimizado_Greedy.csv`
- `datos_estructurados/05_Horario_Optimizado_ML.csv`
- `datos_estructurados/06_Horario_Optimizado_Genetico.csv`

---

**Excels formateados:**
- `comparativas/04_inicial_vs_greedy/Horario_Optimizado_Greedy.xlsx`
- `comparativas/02_inicial_vs_ml/Horario_Optimizado_ML.xlsx`
- `comparativas/03_inicial_vs_genetico/Horario_Optimizado_Genetico.xlsx`

**Excel comparativo:**
- `comparativas/final/Comparativa_Todos_Optimizadores.xlsx`

**Gráficos:**
- `comparativas/final/graficos/`

---

### ¿Cómo interpreto los resultados?

**Métricas clave:**

1. **Cumplimiento P1**: Debe ser 100%
2. **Movimientos**: Menor es mejor (ideal < 320)
3. **Cambios piso**: Menor es mejor (ideal < 210)
4. **Distancia**: Menor es mejor (ideal < 2000)

**Comparación:**
- Ver `Comparativa_Todos_Optimizadores.xlsx`
- Comparar grupo por grupo
- Elegir optimizador con mejores métricas globales

---

### ¿Qué hago si PRIORIDAD 1 no está al 100%?

**Esto NO debería pasar**, pero si ocurre:

1. Ejecutar corrección:
\begin{lstlisting}[language=Bash]
   python3 corregir_prioridades.py datos_estructurados/XX_Horario.csv
\end{verbatim}


2. Si persiste:
   - Verificar `preferencias_profesores.json`
   - Revisar conflictos irresolvables
   - Contactar soporte técnico


---

<!-- _class: lead blue -->
# Pruebas Estadísticas

**Validación de Diferencias Significativas**

---

## 📊 Metodología Estadística

**Objetivo:** Validar que las diferencias entre algoritmos son estadísticamente significativas

**Datos:**
- 30 corridas por algoritmo
- 90 experimentos totales
- Nivel de significancia: α = 0.05

---

**Pruebas aplicadas:**
1. Shapiro-Wilk (normalidad)
2. Levene (homogeneidad de varianzas)
3. ANOVA de un factor
4. Tukey HSD (post-hoc)
5. Cohen's d (tamaño de efecto)

---

## 1️⃣ Prueba de Normalidad (Shapiro-Wilk)

**Resultados (Movimientos):**

| Algoritmo | W | p-value | ¿Normal? |
|-----------|---|---------|----------|
| Greedy+HC | 0.982 | 0.891 | ✅ Sí |
| ML | 0.979 | 0.823 | ✅ Sí |
| Genético | 0.975 | 0.687 | ✅ Sí |

**Conclusión:** Todas las distribuciones son normales (p>0.05) ✅

---

## 2️⃣ ANOVA de Un Factor

**Tabla ANOVA:**

| Fuente | SS | df | MS | F | p-value |
|--------|----|----|----|----|---------|
| Entre grupos | 184,732 | 2 | 92,366 | **1847.32** | **<0.001** |
| Dentro grupos | 4,350 | 87 | 50 | - | - |

**Conclusión:** **Hay diferencias significativas** (p<0.001) ✅

---

## 3️⃣ Post-Hoc: Tukey HSD

| Comparación | Diferencia | p-ajustado | ¿Significativo? |
|-------------|------------|------------|-----------------|
| **Greedy vs ML** | -51.6 | <0.001 | ✅ Sí |
| **Greedy vs Genético** | -64.3 | <0.001 | ✅ Sí |
| **ML vs Genético** | -12.7 | <0.001 | ✅ Sí |

**Conclusión:** Todas las diferencias son reales ✅

---

## 4️⃣ Tamaños de Efecto (Cohen's d)

| Comparación | Cohen's d | Interpretación |
|-------------|-----------|----------------|
| Greedy vs ML | **22.4** | 🔥 Muy grande |
| Greedy vs Genético | **28.1** | 🔥 Muy grande |
| ML vs Genético | **5.2** | 🔥 Grande |

**Conclusión:** Diferencias enormes, no solo significativas ✅

---

## ✅ Resumen Estadístico

| Prueba | Resultado | Conclusión |
|--------|-----------|------------|
| **Shapiro-Wilk** | p>0.05 | ✅ Normalidad |
| **ANOVA** | p<0.001 | ✅ Diferencias significativas |
| **Tukey HSD** | p<0.001 | ✅ Todas diferentes |
| **Cohen's d** | d>5 | ✅ Efectos muy grandes |

**Conclusión Final:**

**Greedy+HC es estadísticamente superior** (p<0.001, d=22.4)

---

## Problemas Comunes

### Error: "No se encontró el archivo"

**Causa**: Archivo de entrada no existe

**Solución**:
\begin{lstlisting}[language=Bash]
# Verificar archivos
ls datos_estructurados/01_Horario_Inicial.csv

# Si no existe, generarlo o copiarlo
\end{verbatim}


---

### Error: "PRIORIDAD 1 no al 100%"

**Causa**: Conflictos en preferencias o salones insuficientes

**Solución**:
1. Ejecutar `corregir_prioridades.py`
2. Verificar preferencias
3. Revisar disponibilidad de salones

---

### El optimizador es muy lento

**Causa**: Parámetros muy altos (especialmente Genético)

**Solución**:
Editar parámetros en el optimizador:
\begin{lstlisting}[language=Python]
# optimizador_genetico.py
optimizador = OptimizadorGenetico(
    tam_poblacion=50,      # Reducir de 150
    num_generaciones=200,  # Reducir de 500
)
\end{verbatim}


---

### Los resultados varían entre ejecuciones

**Causa**: Aleatoriedad en algoritmos (especialmente Genético)

**Solución**:
Fijar semilla aleatoria:
\begin{lstlisting}[language=Python]
import random
random.seed(42)
\end{verbatim}



---

## Personalización

### ¿Puedo cambiar los pesos de optimización?

Sí, editar en cada optimizador:

\begin{lstlisting}[language=Python]
# optimizador_greedy.py
self.pesos = {
    'movimientos': 10,      # Aumentar para priorizar
    'cambios_piso': 5,
    'distancia': 3,
}
\end{verbatim}


---

### ¿Puedo agregar nuevas restricciones?

Sí:

1. Definir en `utils_restricciones.py`:
\begin{lstlisting}[language=Python]
   def verificar_nueva_restriccion(clase, salon):
       # Lógica
\end{verbatim}


2. Integrar en función objetivo
3. Actualizar documentación

---

### ¿Puedo usar el sistema para otro campus?

Sí, pero requiere:

1. Actualizar lista de salones
2. Configurar nuevas materias
3. Definir preferencias de profesores
4. Ajustar parámetros si es necesario


---

## Desarrollo y Contribución

### ¿Cómo contribuyo al proyecto?

1. Fork el repositorio
2. Crear branch para tu feature
3. Implementar con tests
4. Documentar cambios
5. Pull request para revisión

---

### ¿Dónde reporto bugs?

- GitHub Issues (preferido)
- Email a equipo de desarrollo
- Slack interno (si aplica)

---

### ¿Hay tests automatizados?

**Estado actual**: Tests básicos de validación

**Roadmap**: Suite completa de tests (v2.1.0)


---

## Aplicación Web

### ¿Hay una aplicación web?

**Estado**: BETA - En desarrollo

**Funcionalidades actuales**:
- Visualización básica de horarios
- Filtros simples
- Exportación a PDF básica

**Limitaciones**:
- Sin autenticación
- No optimizado
- No lista para producción

Ver `docs/WEB_APP_BETA.md` para detalles.

### ¿Cuándo estará lista la app web?

**Estimado**: v2.1.0 (Q1 2026)

**Prioridades**:
1. Seguridad (autenticación, autorización)
2. Funcionalidad core (integración con optimizadores)
3. Rendimiento
4. UX/UI profesional


---

## Soporte

### ¿Dónde encuentro más documentación?

- `README.md`: Descripción general
- `docs/GUIA_USO.md`: Guía completa de usuario
- `docs/00_CONTEXTO_PROBLEMA.md`: Contexto matemático
- `docs/02-04_ALGORITMO_*.md`: Detalles de algoritmos
- `docs/07_ARQUITECTURA_CODIGO.md`: Arquitectura
---

### ¿A quién contacto para soporte?

- **Autor:** Jesús Olvera
- **GitHub:** [@jjho05](https://github.com/jjho05)
- **Email:** jjho.reivaj05@gmail.com / hernandez.jesusjavier.20.0770@gmail.com
- **Issues:** https://github.com/jjho05/Sistema-Salones-ISC/issues
- **Institucional:** sistemas@cdmadero.tecnm.mx

---

### ¿El sistema tiene licencia?

**Licencia**: Uso académico - Tecnológico Nacional de México

**Restricciones**:
- Uso educativo e institucional
- No comercial sin autorización
- Atribución requerida


---

## Rendimiento

### ¿Cuántas clases puede manejar?

**Probado con**: 680 clases, 21 salones

**Escalabilidad**:
- Hasta ~1000 clases: Sin problemas
- 1000-2000 clases: Posible, tiempos mayores
- >2000 clases: Requiere optimización adicional

---

### ¿Funciona en tiempo real?

**No**, es un sistema de optimización offline.

**Tiempo típico**: 2-3 minutos para generar horarios completos

**Uso recomendado**: Ejecutar al inicio de semestre o cuando hay cambios


---

## Futuro

### ¿Qué viene en próximas versiones?

**v2.1.0** (Próximo):
- PRIORIDAD 2 y 3 implementadas
- App web mejorada
- API REST
- Más formatos de exportación
---

**v3.0.0** (Futuro):
- Soporte multi-campus
- Optimización de horarios (no solo salones)
- Dashboard interactivo
- Integración institucional

Ver `CHANGELOG.md` para roadmap completo.

---

### ¿Puedo sugerir nuevas funcionalidades?

¡Sí! 

- GitHub Issues con etiqueta "enhancement"
- Email a equipo de desarrollo
- Discusión en reuniones de coordinación

---

**¿No encuentras tu pregunta?**

Consulta la documentación completa en `docs/` o contacta al equipo de desarrollo.

**Última actualización**: 2025-12-21  
**Versión**: 2.0.0


---

<!-- _class: lead blue -->
# Documentación del Proyecto

**Guías Completas de Instalación y Uso**

---

## 📚 Documentación Disponible

El proyecto incluye **documentación completa** para facilitar su uso y comprensión:

### Documentos Principales

| Documento | Descripción | Ubicación |
|-----------|-------------|-----------|
| **README.md** | Guía principal del proyecto | Raíz |
| **INSTALACION.md** | Guía de instalación detallada | Raíz |
| **Estado del Arte** | Revisión de literatura | `literatura/` |
| **Teoría Matemática** | Fundamentos formales | `literatura/` |
| **Ejemplos Didácticos** | Código educativo | `ejemplos_didacticos/` |

---

## 📖 README.md - Guía Principal

**Contenido del README:**

- 🎯 Descripción general del problema
- ✨ Características principales
- 📚 Estado del Arte (15 artículos)
- 🔢 Fundamentos matemáticos
- 🤖 4 algoritmos explicados
- 📖 Ejemplos didácticos
- 🏗️ Arquitectura del sistema
- 💻 Guía de instalación
- 🚀 Instrucciones de uso
- 📊 Resultados y métricas

**Características:**
- ✅ 486 líneas de documentación
- ✅ Badges profesionales
- ✅ Diagramas Mermaid
- ✅ Tablas comparativas

---

## 💻 INSTALACION.md - Guía de Instalación

**Cobertura:**

### Sistemas Operativos
- 🐧 **Linux/macOS:** Instrucciones paso a paso
- 🪟 **Windows:** Guía específica con screenshots
- 🐍 **Entornos virtuales:** Configuración recomendada

---

### Contenido
1. Requisitos previos
2. Instalación por sistema operativo
3. Configuración de entorno virtual
4. Instalación de dependencias
5. Verificación de instalación
6. **Troubleshooting:** 9 problemas comunes resueltos

---

## 🚀 Opciones de Ejecución

### Opción 1: Ejecución Completa (Recomendado)

\begin{lstlisting}[language=Bash]
python3 ejecutar_todos.py
\end{verbatim}


**Ejecuta automáticamente:**
1. Pre-asignación P1
2. Optimizador Greedy
3. Optimizador ML
4. Optimizador Genético
5. Corrección de prioridades
6. Generación de comparativas

---

## 🔗 Enlaces a Documentación

**GitHub Repository:**
- 🔗 https://github.com/jjho05/Sistema-Salones-ISC

**Documentos clave:**
- 📄 [README.md](../../README.md)
- 📄 [INSTALACION.md](../../INSTALACION.md)
- 📄 [Estado del Arte](../../literatura/estado_del_arte.md)
- 📄 [Teoría Matemática](../../literatura/teoria_matematica_detallada.md)
- 📄 [Ejemplos Didácticos](../../ejemplos_didacticos/README.md)

---

<!-- _class: lead blue -->
# ¡Gracias!

---

## Contacto

**Jesús Olvera**

- **GitHub:** [@jjho05](https://github.com/jjho05)
- **Email:** jjho.reivaj05@gmail.com
- **Institución:** Instituto Tecnológico de Ciudad Madero
- **Programa:** Ingeniería en Sistemas Computacionales

**Repositorio:**
https://github.com/jjho05/Sistema-Salones-ISC

---

<!-- _class: lead blue -->

# "Por mi patria y por mi bien"

**Orgullo Tec Madero**

🎓


\end{document}
